\documentclass[11pt]{article}
\usepackage[margin=1in]{geometry}
\usepackage{amsmath,amssymb,amsthm}
\usepackage[hidelinks]{hyperref}
\newcommand{\hashpart}[1]{\texttt{#1}\allowbreak}
\newtheorem{theorem}{Theorem}[section]
\newtheorem{lemma}[theorem]{Lemma}
\newtheorem{proposition}[theorem]{Proposition}
\newtheorem{corollary}[theorem]{Corollary}
\newtheorem*{auxlemma}{Lemma}
\title{Affine companion groups in characteristic two: scalar realizations and density}
\author{Baohua Li}
\date{October 12, 2026}
\begin{document}
\maketitle
\begin{abstract}
We study scalar realizations of affine companion groups in characteristic two.
For $q=2^n$, $d\mid q-1$, and $k=(q-1)/d$, we give an exact
M\"obius formula and an explicit error for the number of $a$ satisfying
$a,1+a\in C_d$ and $\langle a,1+a\rangle=C_d$. At the quartic norm boundary
$t=2^s$, $s\ge2$, and $d=(t+1)(t^2+1)$, the full-family joint-generation
count is
\[
 |\mathcal R_d|=\theta_2(d)(t+1)^2+O(t^{3/2}2^{\omega(d)}),
 \qquad \theta_2(d)=\prod_{p\mid d}(1-p^{-2}).
\]
Each norm fiber has $O(t)$ points. For every nonexceptional fiber we give
an explicit $O(t)$ main term with a norm-dependent coefficient and a
square-root-scale error; the exceptional fibers are bounded separately.
Exact local averaging over all $t+1$ fibers yields the quadratic
full-family main term. The number $D_d$ of slopes satisfying
$\operatorname{ord}(a)=\operatorname{ord}(1+a)=d$ obeys
$D_d\sim(\varphi(d)/d)^2(t+1)^2$, with an explicit error, and is positive
for every $s\ge2$. This exact-order pair is a standard
$(t-1)$-primitive-pair object; we do not claim its first formulation.
Our contribution is the norm-fiber density analysis, its companion-group
interpretation, and the all-dimensional existence proof.
We also determine affine generation criteria and prove the optimal uniform
index exponent $1/4$; further sections treat spread statistics and a
nonlinear example generating $A_{16}$. Lean checks only the stated triality
and field-Singer results; the density and existence arguments use classical
curve bounds and exact arithmetic.
\end{abstract}

\section{Introduction}\label{sec:intro}
Which primitive affine groups can be generated by a complete mapping
and its prescribed companion, and how many slopes realize each group?
For a complete mapping $f$ of an elementary abelian $2$-group, write
$F(f)=(\mathrm{id}\oplus f)^{-1}$. Theorem~\ref{thm:affine-group}
determines the group generated by each affine pair:
for $f(x)=Mx\oplus c$, it is
$\mathbb F_2[M]c\rtimes\langle M,(I+M)^{-1}\rangle$.
For scalar maps, this reduces realization to the joint generation
condition $\langle a,1+a\rangle=C_d$ in a specified cyclic subgroup.

\paragraph{The quartic boundary: existence and density.}
Our principal arithmetic family has $t=2^s$, $s\ge2$,
$L=\mathbb F_{t^4}$, $K=\mathbb F_t$, and
$d=(t+1)(t^2+1)$. Put $C_d=\ker N_{L/K}$, and let
$\mathcal R_d$ consist of the slopes $a,1+a\in C_d$ that
jointly generate $C_d$. Let $\mathcal D_d$ be the subset of these slopes for which both orders equal $d$, and write $D_d=|\mathcal D_d|$. Since each full-order element individually generates $C_d$, we have $\mathcal D_d\subseteq\mathcal R_d\subseteq\{a\in C_d:1+a\in C_d\}$. Thus $\mathcal R_d$ counts ordinary normalized triangles whose two ratios jointly generate $C_d$, without requiring either ratio individually to have order $d$. Theorems~\ref{thm:inner-joint-density}
and~\ref{thm:inner-double-count} give
\[
 |\mathcal R_d|
   =\theta_2(d)(t+1)^2+O(t^{3/2}2^{\omega(d)}),
 \qquad
 D_d\sim\theta(d)^2(t+1)^2,
\]
where $\theta_2(d)=\prod_{p\mid d}(1-p^{-2})$,
$\theta(d)=\varphi(d)/d$, and both errors are explicit.
Moreover, $D_d>0$ for every $s\ge2$. Thus every member of
this inner quartic family realizes its primitive affine group
with both slopes of full subgroup order.

The count and the all-dimensional existence statement use
complementary arguments. The density proof descends to genus-one
norm fibers, classifies the constant norm-character exceptions,
and proves an $8\sqrt t$ mixed-character bound
(Proposition~\ref{prop:inner-mixed-character}).
A fixed fiber has size $O(t)$ and coefficient
$\lambda_{t+1}(b)\theta_2(t^2+1)$, which depends on its norm
parameter $b$. Averaging these coefficients gives the quadratic
full-family term $\theta_2(d)(t+1)^2$.
For existence at every parameter, Proposition~\ref{prop:inner-diagonal-double}
uses the trace-one construction
\[
 a=(1+\gamma)^{-1},\qquad
 \gamma\in C_{t^2+1},\qquad
 \gamma+\gamma^{-1}\in C_{t+1},\qquad
 1+a=a^{t^2}.
\]
The two required orders are tied by Frobenius conjugacy.
A separate four-point Lang cover and a two-character sieve prove
positivity for every $s\ge33$; exact field certificates cover
$2\le s\le32$. The intermediate range $33\le s\le113$ is
covered by exact divisor bounds, and a uniform estimate covers
every $s\ge114$. This is a complete parameter partition.
The trace-one subfamily itself is empty at $s=3,5$; the
certificates at these parameters belong to the full family.

\paragraph{Prescribed-order predecessors and quantitative comparison.}
The exact-order pair above is a $(t-1)$-primitive pair in the
standard terminology: $d=(t^4-1)/(t-1)$, and exact order $d$
already implies norm one because $N_{L/K}(a)=a^d$.
Thus the norm condition is automatic for $D_d$.
Prescribed-order pairs, their freeness indicators, and their
multiplicative-character sieves have established predecessors
\cite{CKR22,AN22,SSL24}. The distinction developed here is the
effective estimate after descent to the specified norm fibers.
For example, \cite[Corollary~17]{CKR22}, with ambient field
size $t^4$, polynomials $X,X+1$, subgroup indices $t-1,t-1$,
and freeness parameters $d,d$, requires
\[
 t^2>(t-1)^2W(d)^2,\qquad W(d)=2^{\omega(d)}.
\]
This fails for every $t\ge4$: the coprime odd factors $t+1$
and $t^2+1$ are both greater than one, so $W(d)\ge4$.
More directly, \cite[Proposition~3.1 and Theorem~1]{SSL24}
count the same adjacent exact-order pairs:
$D_d=N_{t-1}(t^4,1,2)$ in their notation. Their sufficient
condition at these parameters is
\[
 t^2>\frac{((t-1)W(d))^3}{(t-1)W(d)-1},
\]
which also fails for every $t\ge4$.
Our estimates instead exploit the lower-dimensional norm geometry.
Failure of these sufficient criteria is not a claim that
every possible refinement of the earlier methods has been excluded.

The Frobenius relation in the trace-one construction is classical.
The reciprocal-trace reduction in \cite[\S4]{Coh} already connects
norm-one elements with trace-one elements by a fractional linear
change of variable. With its parameter $q=t^2$,
\cite[Theorem~1.5]{Coh} gives
$\operatorname{ord}(\gamma)=t^2+1$ and
$\operatorname{ord}(\gamma+\gamma^{-1})=t^2-1$.
The latter order differs from the required $t+1$; its norm
to $K$ has order $t-1$, rather than one.
Proposition~\ref{prop:inner-diagonal-double} treats this
additional reciprocal-trace norm fiber, its four-point conductor,
and its prescribed-order count. We do not claim the fractional
linear construction or Frobenius coupling as a new mechanism.
Hachenberger's quartic projective treatment is a closer geometric
predecessor \cite[Remark~4.2 and Proposition~7.1]{Ha21}.
In characteristic two, $\gcd(t-1,d)=1$ gives
$L^\times=K^\times\times C_d$, so each projective point has a
unique representative in $C_d$; a primitive projective point
then corresponds exactly to a representative of order $d$.
Moreover, every primitive ovoid in his partition meets the
absolute trace-zero plane in an oval of $t+1$ points.
For the construction in Proposition~\ref{prop:inner-diagonal-double},
put $x=(1+\gamma)^{-1}$ and $\delta=\gamma+\gamma^{-1}$.
Then $x^{t^2}=1+x$ and $x^{t^2+1}=\delta^{-1}$. Since
$t^2+1\equiv2\pmod{t+1}$ and $t+1$ is odd, the
$C_{t+1}$-component of $x$ has order $t+1$; hence $[x]$ lies in a
primitive ovoid, and $[x]$ and $[1+x]$ lie in its trace-zero
oval. Thus this projective model and ambient oval incidence are
prior. The additional condition here is that the paired points
both have exact order $d$ and are exchanged by $t^2$-Frobenius.
The cited incidence theorem does not count this arithmetic
subfamily. We claim neither the ambient ovoid--plane incidence
nor the reciprocal-trace change of variables as new.

The line and translate properties of $r$-primitive elements
were studied in \cite{CK21}; Remark~4.1 there already allows
sufficiently slow growth of $r$. Varying the index alone
therefore does not establish novelty.
Choudhary--Sharma \cite[Definition~3.1 and Theorem~1]{CSh25}
provide a direct predecessor for the exact count $D_d$. In their
notation, set $q=t$, $n=4$, $m=2$, $\gamma=1$,
$r_1=r_2=t-1$, and $a_1=a_2=1$. Then
$r_i'=t-1$, $R_i=d$, $R_i'=1$, and $Q_{R_i}=d$, so their
$N_{1,1}(d)$ is exactly $D_d$. We therefore do not claim priority
for formulating or counting this adjacent exact-order pair.
Their Theorem~1 criterion becomes
$1>2(t-1)^2W(d)^2$, which fails for every $t\ge4$.
For Theorem~3, write $U=\sum_i u_i$ for the total number of
sieved primes and $\Delta=(U-1)/\delta+2$. When $U=0$,
$\Delta=1$ and the right-hand side is at least
$2(t-1)^2\ge18$; when $U\ge1$, $\Delta\ge2$ and it is
at least $4(t-1)^2\ge36$. Theorem~4 also fails for every legal
split: its positive denominator is at most $1$, its numerator is
at least $\sum_i c_i\ge2$, and the full right-hand side is at
least $4(t-1)^2\ge36$. Theorem~2 assumes $q\equiv3\pmod4$
and odd $n$, and the paper's application theorem treats $r=3$,
$n\ge7$; neither applies here. These are failures of the stated
sufficient criteria, not of every possible refinement.
The quantitative contribution is the explicit $O(t)$ estimate on
individual norm fibers, with their norm-dependent main coefficients, and
its exact averaging to the $O(t^2)$ full-family count. The adjacent
exact-order statistic is treated as an application, not claimed as a new
formulation. The ordinary cyclotomic intersection count is studied,
in particular, in \cite{KN26}; it is not claimed as new.
The Lang construction, curve character bound, and primitive
indicators are classical \cite{Katz89,FFSTV13,CKR22}.
Our proofs verify geometric monodromy, ramification and
the resulting errors for the specific realization problem.

\paragraph{Group structure and the other index ranges.}
The affine group formula gives transitivity, primitivity and
double-transitivity criteria (Corollary~\ref{cor:affine-actions}),
inverse realization and a divisor-sum sieve
(Corollaries~\ref{cor:scalar-realization} and~\ref{cor:scalar-sieve}).
A Jacobi-sum estimate realizes every fixed odd index greater
than one in all sufficiently large admissible dimensions
(Theorem~\ref{thm:fixed-index-realization}), with the exact
index-three threshold in Corollary~\ref{cor:index-three-realization}.
Theorem~\ref{thm:scalar-density} and
Corollary~\ref{cor:growing-index-realization} give the uniform
range $k\le q^{1/4}/\sqrt{\log_2q}$.
Norm descent (Theorem~\ref{thm:norm-subfield-descent}) supplies
nonrealizable primitive groups at $k=q^{1/4}+1$ for every
$q=2^{4s}$ (Corollary~\ref{cor:quarter-power-barrier}).
Thus the uniform exponent $1/4$ is optimal, and the adjacent
quartic families $k=q^{1/4}-1$ and $k=q^{1/4}+1$ have opposite
realization behavior. The universal guarantee still has a
logarithmic gap; we do not classify all indices at the boundary.
The single full-order count and all-dimensional realization
also appear in Theorems~\ref{thm:inner-primitive-count}
and~\ref{thm:inner-all}. Primitive inverse classification and
generating triangles in generalized Paley graphs are given in
Propositions~\ref{prop:primitive-scalar-classification}
and~\ref{prop:Paley-generating-triangles}.

\paragraph{Complete mappings, triality and spread geometry.}
Complete mappings connect finite groups and Latin-square
constructions \cite{HP55}; over a finite field they also have
permutation-polynomial representatives. In characteristic two
complete mappings and orthomorphisms coincide.
Cycle types and groups generated by their full collections
are studied in \cite{BW21,BW22}. Here the individual prescribed
pair is the group-generating object.
The fixed maps of $F$ correspond to oriented line spreads.
An orientation-sensitive Burnside formula
(Proposition~\ref{prop:oriented-Burnside}) gives, in dimension
four, exactly three classes and their stabilizers
(Theorem~\ref{thm:fixed-n4-conjugacy}).
The local invariant-orientation criterion in
Proposition~\ref{prop:orientation-local} is a deduction from
classical group actions. All six-dimensional standard-spread
counts concern one specified underlying spread type.

For the standard Desarguesian spread on $\mathbb F_4^m$, semilinear
conjugation permutes the line orientations and reverses all their signs.
Since the number of lines is odd, the semilinear coset fixes no orientation.
This yields a simplified class formula and an explicit asymptotic estimate
(Theorem~\ref{thm:standard-asymptotic}). The six-dimensional standard family
has $80$ classes (Corollary~\ref{cor:standard-n6}), with all representatives
and stabilizer orders supplied in the finite data.
The asymptotic argument specializes the classical motion method for
distinguishing colorings \cite{RS98}; existence of regular subsets is studied
more generally in \cite{CNS84,Se97}. These methods are not claimed as new.

The nonlinear group question is already different from the affine formula:
Proposition~\ref{prop:nonlinear-A16} gives a specified pair generating $A_{16}$
whose three companions all have cycle type $1+15$. Its proof uses an explicit
three-cycle certificate and elementary permutation identities.
Proposition~\ref{prop:translated-fixed} explains why simply translating a
fixed map cannot provide transitive examples beyond dimension two.

Cycle-type existence has also been established in broader settings: M\"uyesser
\cite{Muy26} proves the Friedlander--Gordon--Tannenbaum conjecture for sufficiently
large abelian groups. The linked three-map cycle data and the subgroup generated
by an individual map and its companion are the specific objects considered here.

The function identity $F^3=\mathrm{id}$ requires no linearity assumption
(Lemma~\ref{lem:F3}). In the transversal model of the additive Cayley table,
$F$ is induced by the cyclic permutation of the row, column, and symbol
coordinates (Proposition~\ref{prop:F-parastrophy}). This places the identity
within the classical conjugate, or parastrophic, symmetry of Latin squares
\cite{MW22}. Its fixed maps have a geometric description as oriented
line spreads (Proposition~\ref{prop:F-spreads}), which includes nonlinear maps.
For the entire affine class on $\mathbb F_2^n$, Proposition~\ref{prop:F-affine}
gives exact orbit counts. The matrix fixed-space count used there is classical
\cite[Theorem~2.1]{FS15}; a short subspace-lattice proof is included to make
the application to $F$ self-contained. These statements place the scalar affine
classification of Proposition~\ref{prop:F-orbits} in a wider setting.

For the coordinate application, we read the $64$ words of $\mathbb{F}_2^6$
through a fixed ordered basis and complement the resulting coordinates.
This transformation $T$ is an affine Singer map (layer), meaning that its
linear part is a Singer cycle. The
associated triple $(T,D,K)$ and the group $\langle T,D\rangle$ then reduce to
elementary finite-field statements, with a uniform analogue for $2k$-bit systems
(an existence theorem for all $k$, \S\ref{sec:2k} and Appendix~\ref{app:pte}; the concrete coordinate presentation is verified for $k=1,2,3,4$). Related families of
permutations over $\mathbb{F}_2^n$ with prescribed cycle structure and algebraic
degree were studied by Gravel, Panario and Thomson \cite{GPT}. The underlying
even-characteristic existence statement is classical (see \cite{Mor} and
\cite[Theorem~1, p.~1]{Coh90}; an equivalent square-order formulation is \cite[Theorem~1.5, p.~170]{Coh});
Cohen and Presern \cite{CP05} subsequently gave a self-contained account of
the prescribed-trace theorem with reduced computation.
Appendix~\ref{app:pte} reconstructs the even-characteristic trace-one case and
provides an explicit finite certificate; neither the existence theorem nor its
character-sum method is claimed as new.

The paper uses this coordinate-permutation symmetry as a common framework for
the geometric fixed-map description, the affine orbit counts, and the Singer
realization. The geometric classification concerns all complete mappings,
whereas the orbit formula concerns the entire affine class, including matrices
that are not multiplication in any single chosen field structure.
Sections~\ref{sec:T}--\ref{sec:tdk} give the six-dimensional coordinate data,
and Section~\ref{sec:group} determines the general affine companion group; Section~\ref{sec:2k} supplies Singer triples
in every even dimension using the classical existence input.
Section~\ref{sec:rem} develops the orbit classifications, including the standard-spread asymptotics.
Section~\ref{sec:nonlinear} studies nonlinear companion groups. Exact finite data and
Lean sources with pinned dependencies are supplied with the scope stated in
Appendix~\ref{app:code}. No priority is asserted for the coordinate-permutation
symmetry, the classical matrix count, the $56$ spreads of $\mathrm{PG}(3,2)$,
or the prescribed-trace theorem.
The classical names of the words play no role in any proof.

\section{Notation}\label{sec:not}
For the six-bit coordinate example, $V=\mathbb{F}_2^6$; other dimensions
are stated explicitly. Coordinates $x=(x_1,\dots,x_6)$ have $x_1$ the most significant
bit (value $32$) and $x_6$ the least (value $1$); a word is identified with its
integer value in $\{0,\dots,63\}$. Basis $G=(\texttt{63},\texttt{33},\texttt{18},
\texttt{49},\texttt{1},\texttt{21})$ (a fixed ordered basis --- the six words are linearly independent over $\mathbb{F}_2$, by Gaussian elimination; non-uniqueness is discussed in \S\ref{sec:2k}); write $g_j$ for
the $j$-th entry, $j=0,\dots,5$.
$\mathrm{coords}_G(x)=(c_1,\dots,c_6)$ is the unique vector with
$x=\bigoplus_{i=1}^6 c_i\,g_{i-1}$, read as the six-bit word $c_1\cdots c_6$
($c_1$ most significant); $\oplus$ is bitwise XOR, i.e.~addition in $\mathbb{F}_2^n$.
When a field structure is chosen in characteristic two, $\oplus$ is also its
field addition. A monic polynomial is called \emph{primitive} if it is
irreducible and its roots generate the multiplicative group of its root field.
Superscripts on functions and on $F$ denote composition iterates; orders of
field elements denote multiplicative orders.

\section[Affine Singer map]{$T$ is an affine Singer map}\label{sec:T}
\subsection*{3.1 Definition and basic checks}
$T(x)=63\oplus \mathrm{coords}_G(x)$. Verified identities:
$T(0)=63$, $T(41)=0$, $T(38)=38$; for example
$\mathrm{coords}_G(38)=(0,1,1,0,0,1)$, read as $011001=25$, so $T(38)=63\oplus25=38$,
and $\mathrm{coords}_G(41)=(1,1,1,1,1,1)=63$, so $T(41)=63\oplus63=0$.
The integer labels $38,41,63$ belong to this chosen coordinate presentation;
they are not invariants under changes of basis (see \S\ref{sec:2k}).
\begin{lemma}\label{lem:aff}
$T$ is affine: $T(x)=L_Tx+b$ with $b=T(0)=63$ and $L_T(x)=T(x)\oplus 63$.
\end{lemma}
\begin{proof}
Since $G$ is a basis, $\mathrm{coords}_G$ is an $\mathbb F_2$-linear
isomorphism. Thus $L_T=\mathrm{coords}_G$, and adding the constant $63$
gives the stated affine map.
\end{proof}
\begin{lemma}\label{lem:fix}
$\mathrm{Fix}(T)=\{38\}$; on the remaining $63$ points $T$ is a single cycle
(type $1+63$, order $63$), equivalently $L_T$ has order $63$ (Proposition~\ref{prop:singer}).
\end{lemma}
\begin{proposition}\label{prop:singer}
Identify $(V,+)$ with $\mathbb{F}_{64}$; then $L_T$ is
multiplication by a primitive element of $\mathbb{F}_{64}$, a Singer cycle \cite{LN}. The
characteristic polynomial of $L_T$ is
\begin{equation}\label{eq:pT}
p_T(t)=t^6+t^5+t^3+t^2+1,
\end{equation}
computed directly as $\det(tI+L_T)$ (code appendix), primitive of degree $6$.
\end{proposition}
\begin{proof}
The matrix in Appendix~\ref{app:code} gives the displayed determinant.
Modulo $p_T$, repeated squaring and multiplication give
\[
t^{64}=t,\qquad \gcd(p_T,t^8+t)=\gcd(p_T,t^4+t)=1.
\]
The degree-$6$ irreducibility criterion, using the prime divisors $2,3$ of
$6$, therefore shows that $p_T$ is irreducible. For completeness,
$p_T(0)=p_T(1)=1$ also gives $\gcd(p_T,t^2+t)=1$.
In $\mathbb F_2[t]/(p_T)$ the residue $\bar t$ satisfies
\[
\bar t^{63}=1,\qquad
\bar t^{21}=\bar t^4+\bar t^2+\bar t+1\ne1,\qquad
\bar t^9=\bar t^2+\bar t\ne1.
\]
Since $63=3^2\cdot7$, its order is $63$. The irreducible characteristic
polynomial identifies $L_T$ with multiplication by $\bar t$ on this field.
Finally, $T(38)=38$ implies $T(x)\oplus38=L_T(x\oplus38)$, proving
the fixed-point and cycle claims of Lemma~\ref{lem:fix}.
\end{proof}


\subsection*{3.2 Powers and $63=7\cdot 9=3\cdot 21$}
Since $L_T$ has order $63$, $T^d$ for $d\mid 63$ has type $1+d\times(63/d)$ (one fixed point plus $d$ cycles of length $63/d$):
$T^1=1+63$, $T^3=1+3\cdot21$, $T^7=1+7\cdot9$, $T^9=1+9\cdot7$,
$T^{21}=1+21\cdot3$ (verified).

\section[Complete mappings and the (T,D,K) triple]{Complete mappings and the $(T,D,K)$ triple}\label{sec:tdk}
A permutation $f$ of an additive group of exponent two is a \emph{complete
mapping} if $\mathrm{id}\oplus f$ is again a permutation. Define its
\emph{triality operator} $F(f)=(\mathrm{id}\oplus f)^{-1}$, where
$(\mathrm{id}\oplus f)(x)=x\oplus f(x)$. By Proposition~\ref{prop:singer},
identify the linear part of $T$ with multiplication by a primitive
$\alpha\ne1$. Then $\mathrm{id}\oplus T$ has nonzero slope
$1\oplus\alpha$ and is therefore a permutation. Its inverse $D=F(T)$ has slope
$(1\oplus\alpha)^{-1}$, and
$1\oplus(1\oplus\alpha)^{-1}=\alpha/(1\oplus\alpha)\ne0$, so
$\mathrm{id}\oplus D$ is also a permutation. Thus
$D=F(T)=(\mathrm{id}\oplus T)^{-1}$ and $K=F(D)=(\mathrm{id}\oplus D)^{-1}$ are
defined; since $F^3(T)=T$ below, $\{T,D,K\}$ is the $F$-orbit of $T$.
Verified facts:
\begin{itemize}
\item cycle types $T=D=1+63$, $K=1+7\cdot 9$ (seven $9$-cycles plus one fixed point);
\item fixed points $\mathrm{Fix}(T)=\{38\}$, $\mathrm{Fix}(D)=\{41\}$,
$\mathrm{Fix}(K)=\{63\}$;
\item $F(K)=T$, i.e.\ $(\mathrm{id}\oplus K)^{-1}=T$; hence $F^3=\mathrm{id}$ with
$T\to D\to K\to T$;
\item offset-fixed three-cycle: $D(0)=38=\mathrm{Fix}(T)$,
$K(0)=41=\mathrm{Fix}(D)$, $T(0)=63=\mathrm{Fix}(K)$;
\item linear parts: for any $n\times n$ matrix $A$ in characteristic two with
$p_A(t)=\det(tI+A)$ and $p_A(1)\ne0$, taking determinants gives
\[
p_{(I+A)^{-1}}(t)
=\frac{\det\bigl(t(I+A)+I\bigr)}{\det(I+A)}
=\frac{t^n p_A(1+t^{-1})}{p_A(1)}.
\]
Here $p_T(1)=1$ and $p_T(t+1)=p_T(t)$ (directly,
$(t+1)^6+(t+1)^5+(t+1)^3+(t+1)^2+1$ reduces to
$t^6+t^5+t^3+t^2+1$). Hence $L_D=(I+L_T)^{-1}$ has
$p_D(t)=t^6p_T(1+t^{-1})=t^6p_T(t^{-1})=t^6+t^4+t^3+t+1$
(primitive). The reciprocal form uses this shift invariance.
$L_K$ has $\Phi_9(t)=t^6+t^3+1$
(cyclotomic, non-primitive, order $9$).
\end{itemize}
\begin{lemma}\label{lem:fp}
Every complete mapping $f$ of an additive group of exponent two has exactly one
fixed point, namely $F(f)(0)$.
\end{lemma}
\begin{proof}
The equation $f(x)=x$ is equivalent to $x\oplus f(x)=0$.
Since $\mathrm{id}\oplus f$ is bijective, its unique solution is $F(f)(0)$.
\end{proof}

\begin{lemma}\label{lem:F3}
On complete mappings of any additive group of exponent two, $F$ is defined
at every iterate and satisfies
\[
F^2(f)=\mathrm{id}\oplus f^{-1},\qquad F^3(f)=f.
\]
\end{lemma}
\begin{proof}
Put $g=(\mathrm{id}\oplus f)^{-1}$. The identity
$x=g(x)\oplus f(g(x))$ gives
\[
\mathrm{id}\oplus g=f\circ g.
\]
Thus $g$ is complete. Taking inverses yields
\[
F^2(f)=(f\circ g)^{-1}
=g^{-1}\circ f^{-1}
=(\mathrm{id}\oplus f)\circ f^{-1}
=\mathrm{id}\oplus f^{-1}.
\]
Consequently $\mathrm{id}\oplus F^2(f)=f^{-1}$ and $F^3(f)=f$.
Only addition of exponent two and bijectivity were used.
\end{proof}

\begin{proposition}\label{prop:F-parastrophy}
For a complete mapping $f$ of an additive group $V$ of exponent two, put
\[
\Gamma_f=\{(x,f(x),x\oplus f(x)):x\in V\},\qquad J(f)=f^{-1}.
\]
Permuting the three coordinates preserves this family of sets. The permutations
$\rho(x,y,z)=(z,x,y)$ and $\sigma(x,y,z)=(y,x,z)$ induce $F$ and $J$,
respectively. Thus $J$ preserves complete mappings and
\[
F^3=J^2=\mathrm{id},\qquad JFJ=F^{-1}=F^2.
\]
For finite $V$, $\Gamma_f$ is a transversal of its additive Cayley table.
\end{proposition}
\begin{proof}
Every triple of $\Gamma_f$ has sum zero, and each of its three coordinate
projections is a bijection to $V$. Both properties are invariant under
coordinate permutations. Conversely these properties recover a unique complete
mapping from the first two coordinates. With $z=x\oplus f(x)$, one has
$F(f)(z)=x$ and $z\oplus x=f(x)$, so
\[
\rho(\Gamma_f)=\Gamma_{F(f)}.
\]
With $y=f(x)$, swapping the first two coordinates gives
$(y,f^{-1}(y),y\oplus f^{-1}(y))$, hence
$\sigma(\Gamma_f)=\Gamma_{J(f)}$.
The coordinate permutations satisfy
$\rho^3=\sigma^2=\mathrm{id}$ and $\sigma\rho\sigma=\rho^{-1}$;
injectivity of $f\mapsto\Gamma_f$ transfers these relations to $F,J$.
For finite $V$, the three bijective projections say exactly that the set
contains one cell in each row, each column, and each symbol class of the
additive Cayley table. This is the definition of a transversal.
\end{proof}
This is the cyclic part of the classical coordinate-permutation symmetry of
Latin squares \cite{MW22}, acting here on transversals of a table invariant
under all six coordinate permutations. The proof identifies this relationship
explicitly; the order-three symmetry itself is not claimed as new.

For a scalar affine map $f(x)=\alpha x\oplus c$, put
$\beta=(1\oplus\alpha)^{-1}$ and $\kappa=1\oplus\alpha^{-1}$.
Then $F(f)(x)=\beta x\oplus\beta c$ and
$F^2(f)(x)=\kappa x\oplus\alpha^{-1}c$, so the slopes follow
$\alpha\mapsto\beta\mapsto\kappa\mapsto\alpha$, with
$\alpha\beta\kappa=1$. The exponent-two hypothesis matters: over
$\mathbb F_5$, replacing $\oplus$ by ordinary field addition $+$,
all three iterates of $f(x)=2x+1$ are defined, but
$F^3(f)(x)=2x+2\ne f(x)$.

For the Singer example, in field coordinates
$T(x)=\alpha x\oplus c$, so
$F(T)(0)=D(0)=(1\oplus\alpha)^{-1}c$ solves $x=\alpha x\oplus c$, i.e.\
$D(0)=\mathrm{Fix}(T)$, and the slope cycle
$\alpha\mapsto(1\oplus\alpha)^{-1}\mapsto1\oplus\alpha^{-1}\mapsto\alpha$ gives the
two remaining identities. For the six-bit layer the values are $F(T)(0)=38$,
$F(D)(0)=41$, $F(K)(0)=63$.
Conceptually, $T$ is the launching map --- an affine Singer layer read through the
fixed generator basis (a change of basis) --- while the triality operator $F$ keeps
the third-order cycle running: $T\mapsto D\mapsto K\mapsto T$, carrying the fixed
point $38\mapsto41\mapsto63\mapsto38$.

\section{Groups generated by affine companions}\label{sec:group}
For a vector $u$ write $\tau_u(x)=x\oplus u$.
If $U\le V$ is invariant under $M\le\mathrm{GL}(V)$, we realize
$U\rtimes M$ as $\{x\mapsto Lx\oplus u:L\in M,\ u\in U\}$ in the given coordinates.
Composition has the usual rule $(u,L)(v,N)=(u\oplus Lv,LN)$.

\begin{theorem}[Affine companion group]\label{thm:affine-group}
Let $V=\mathbb F_2^n$, $n\ge1$, and let $f(x)=Ax\oplus c$ be complete,
so $A,I+A\in\mathrm{GL}(n,2)$. Put
\[
\begin{aligned}
B&=(I+A)^{-1},\qquad R=\mathbb F_2[A],\qquad M_A=\langle A,B\rangle,\\
U&=Rc=\operatorname{span}_{\mathbb F_2}\{c,Ac,\ldots,A^{n-1}c\}.
\end{aligned}
\]
For every $c\in V$, including $c=0$,
\[
\langle f,F(f)\rangle=U\rtimes M_A,
\qquad |\langle f,F(f)\rangle|=2^{\dim U}|M_A|.
\]
In particular its pure translations are exactly $\{\tau_u:u\in U\}$.
\end{theorem}
\begin{proof}
The minimal polynomial of $A$ is coprime to $t$ and $t+1$.
Bezout identities therefore show that $A^{-1}$ and $B$ belong to $R$;
their inverses do as well. Thus $M_A\le R^\times$ is abelian and preserves $U$.
Cayley--Hamilton gives the displayed finite spanning set for $U$.
Since $g=F(f)$ has the explicit form $g(x)=Bx\oplus Bc$, both generators
lie in $U\rtimes M_A$, proving one inclusion.

Using $AB=I+B$ and $[f,g]=f\circ g\circ f^{-1}\circ g^{-1}$,
direct composition gives
\[
g\circ f(x)=ABx,\qquad
f\circ g(x)=ABx\oplus Bc,\qquad [f,g]=\tau_{Bc}.
\]
Conjugation by $f$ sends $\tau_u$ to $\tau_{Au}$.
Hence the generated group contains all translations whose vectors lie in
$\operatorname{span}\{A^jBc:j\ge0\}=RBc=Rc=U$, where $B$ is a unit of $R$.
It contains $\tau_c\circ f=A$ and $\tau_{Bc}\circ g=B$ as pure linear maps,
so it contains $M_A$ as well. This proves the reverse inclusion.
Distinct pairs $(u,L)$ give distinct affine maps, yielding the order and
translation statements.
\end{proof}

\begin{corollary}[Action criteria]\label{cor:affine-actions}
In Theorem~\ref{thm:affine-group}, write $H=\langle f,F(f)\rangle$.
Then $H$ is transitive on $V$ if and only if $U=V$, equivalently $c$ is a
cyclic vector for $A$. It is primitive if and only if $c\ne0$ and the
characteristic polynomial of $A$ is irreducible.
If $A$ is identified with multiplication by $\alpha\in\mathbb F_{2^n}$
in this irreducible case, then for $c\ne0$
\[
|H|=2^n\operatorname{lcm}\bigl(\operatorname{ord}(\alpha),
             \operatorname{ord}(1\oplus\alpha)\bigr).
\]
It is doubly transitive if and only if $c\ne0$, $A$ is irreducible, and
this least common multiple is $2^n-1$. In that case
$H=\mathrm{AGL}(1,2^n)$ in these field coordinates.
\end{corollary}
\begin{proof}
The orbit of $0$ is exactly $U$, giving transitivity.
When $U=V$, the group contains the full regular translation group and its
stabilizer of $0$ is $M_A$. A block containing $0$ is an additive subspace:
if $u$ belongs to the block, its translate by $u$ meets the block at $u$,
so the block property makes that translate equal to the block.
The block is also $M_A$-invariant. Conversely the cosets of any
$M_A$-invariant subspace form a block system (cf.\ \cite{DM}).
An $A$-invariant subspace is invariant under the polynomial maps
$B,A^{-1},B^{-1}$, so the $M_A$-invariant subspaces are exactly the
$A$-invariant ones. The action is therefore primitive precisely when
$A$ acts irreducibly; this is equivalent to an irreducible characteristic
polynomial. Such an $A$ makes every nonzero vector cyclic, which explains
the condition $c\ne0$.

An irreducible $A$ identifies $V$ with $\mathbb F_{2^n}$ and $A$ with
multiplication by $\alpha$; the two multipliers generate
$\langle\alpha,(1\oplus\alpha)^{-1}\rangle$ in the cyclic multiplicative group.
Its order is the least common multiple of the two element orders.
Double transitivity means that $M_A$ is transitive on $V\setminus\{0\}$,
which holds exactly when this subgroup is the entire multiplicative group.
Double transitivity implies primitivity, so the conditions are also necessary
outside the irreducible case.
\end{proof}

\begin{corollary}[Realizing prescribed scalar affine subgroups]\label{cor:scalar-realization}
Fix $n\ge2$, put $q=2^n$, and regard $\mathbb F_q$ as an
$n$-dimensional vector space over $\mathbb F_2$. For a divisor
$d>1$ of $q-1$, let $C_d$ be the unique subgroup of order $d$
of $\mathbb F_q^\times$, and define
\[
\mathcal R_d=\{a\in C_d:\ 1+a\in C_d,\
                  \langle a,1+a\rangle=C_d\}.
\]
The particular affine subgroup
$H_d=\mathbb F_q^+\rtimes C_d\le\mathrm{AGL}(1,q)$
is generated by $f$ and $F(f)$ for a scalar affine complete mapping
$f(x)=ax+c$ if and only if
\[
\operatorname{ord}_d(2)=n\quad\text{and}\quad\mathcal R_d\ne\varnothing.
\]
Here $\operatorname{ord}_d(2)$ denotes the multiplicative order of $2$
modulo the odd integer $d$. When the conditions hold, the exact number
of scalar affine maps $f$ realizing this specified subgroup $H_d$
(in these fixed field coordinates) is
\[
(q-1)|\mathcal R_d|.
\]
\end{corollary}
\begin{proof}
If $\langle f,F(f)\rangle=H_d$, Theorem~\ref{thm:affine-group}
shows that its linear part is $\langle a,(1+a)^{-1}\rangle=C_d$,
so $a\in\mathcal R_d$. Its pure translations fill the field,
which forces $c\ne0$ and $\mathbb F_2[a]=\mathbb F_q$.
Since $C_d\le\mathbb F_q^\times$ and $\langle a,1+a\rangle=C_d$,
the least subfield containing $C_d$ is $\mathbb F_2[a]$.
The subfield generated by a multiplicative subgroup of order $d$
has degree $\operatorname{ord}_d(2)$ over $\mathbb F_2$;
thus this order is $n$.

Conversely, suppose $\operatorname{ord}_d(2)=n$ and
$a\in\mathcal R_d$. Because $1+a\in\mathbb F_2[a]$ and
$C_d=\langle a,1+a\rangle$, the field $\mathbb F_2[a]$
contains $C_d$, hence equals $\mathbb F_q$. In particular
$a\ne0,1$, multiplication by $a$ is irreducible over
$\mathbb F_2$, and for any $c\ne0$ one has
$\mathbb F_2[a]c=\mathbb F_q$. Theorem~\ref{thm:affine-group}
now gives $\langle f,F(f)\rangle=H_d$. There are exactly
$|\mathcal R_d|$ admissible slopes and $q-1$ admissible
translations; different pairs $(a,c)$ define different maps.
\end{proof}

\noindent\emph{Scope.} This is an exact realizability test for
specified scalar affine subgroups, not a classification of all
subgroups of $\mathrm{AGL}(n,2)$ up to conjugacy. The remaining
arithmetic question is to determine the nonemptiness and size of
$\mathcal R_d$ uniformly in $n,d$; the criterion by itself is a
consequence of Theorem~\ref{thm:affine-group}.

\begin{corollary}[An arithmetic sieve for realizable slopes]
\label{cor:scalar-sieve}
Retain the notation of Corollary~\ref{cor:scalar-realization},
and for every divisor $e$ of $q-1$ set
\[
T_e=|\{a\in C_e:1+a\in C_e\}|.
\]
Then the number of slopes realizing $C_d$ satisfies
\begin{equation}\label{eq:realization-mobius}
|\mathcal R_d|=\sum_{e\mid d}\mu(d/e)T_e,
\end{equation}
where $\mu$ is the number-theoretic Moebius function. In fact,
\begin{equation}\label{eq:intersection-gcd}
T_e=\deg\gcd_{\mathbb F_2[t]}\bigl(t^e+1,(t+1)^e+1\bigr),
\end{equation}
so these numbers depend on $e$, not on the ambient field containing $C_e$.
Moreover, for $d>3$ one has
\[
|\mathcal R_d|\equiv0\pmod3;
\]
for $d=3$ (necessarily $n$ even), $|\mathcal R_3|=2$.
In particular, the full-translation subgroup
$\mathbb F_{2^n}^+\rtimes C_3$ is realizable by a scalar affine
companion pair in fixed field coordinates exactly when $n=2$.
\end{corollary}
\begin{proof}
Every pair $(a,1+a)$ counted by $T_e$ generates a unique
subgroup $C_s\le C_e$. Therefore
$T_e=\sum_{s\mid e}|\mathcal R_s|$, extending the definition
by $\mathcal R_1=\varnothing$. Moebius inversion yields
\eqref{eq:realization-mobius}.
Since $e\mid q-1$, both polynomials in
\eqref{eq:intersection-gcd} split into simple linear factors over
$\mathbb F_q$ (their derivatives are nonzero, as $e$ is odd).
Their common roots are precisely the elements $a$ counted by $T_e$,
proving the polynomial-gcd formula.

The fractional-linear transformation
$\varphi(a)=(1+a)^{-1}$ preserves $\mathcal R_d$:
indeed $\varphi(a)$ and
$1+\varphi(a)=a/(1+a)$ belong to $C_d$ and generate
exactly the subgroup generated by $a$ and $1+a$.
Direct simplification gives $\varphi^3(a)=a$.
A fixed point obeys $a^2+a+1=0$, hence has multiplicative
order three and can lie in $\mathcal R_d$ only when $d=3$.
Thus for $d>3$ all $\varphi$-orbits in $\mathcal R_d$
have length three. For $d=3$, the two nonidentity elements
of $\mathbb F_4^\times$ comprise $\mathcal R_3$.
Finally $\operatorname{ord}_3(2)=2$,
so Corollary~\ref{cor:scalar-realization} gives the last claim.
\end{proof}

As a symbolic nonrealizability example in dimension six,
\[
\gcd(t^9+1,(t+1)^9+1)=t^2+t+1,
\]
while the analogous gcd for exponent $21$ is
\[
(t^2+t+1)(t^3+t+1)(t^3+t^2+1).
\]
Together with the divisor-sum formula, these identities give
$|\mathcal R_9|=2-2=0$ and
$|\mathcal R_{21}|=8-6-2=0$.
Although $\operatorname{ord}_9(2)=\operatorname{ord}_{21}(2)=6$,
neither $\mathbb F_{64}^+\rtimes C_9$ nor
$\mathbb F_{64}^+\rtimes C_{21}$ can be realized by a scalar affine
companion pair. These are exact polynomial identities, not claims
of a complete nonrealizability classification.

The quantities $T_e$ are standard cyclotomic intersection counts.
For example, if $k=(q-1)/e$ and $\chi$ is a multiplicative
character of order $k$ on $\mathbb F_q$, extending all powers
(including the trivial character) by zero at zero, then
\[
T_e=\frac1{k^2}\sum_{i,j=0}^{k-1}
       \sum_{x\in\mathbb F_q}\chi^i(x)\chi^j(1+x).
\]
This identity is simply the multiplicative-character indicator
of $C_e$. Accounting separately for the trivial-character terms gives
the effective classical estimate in
Proposition~\ref{prop:scalar-positive-bound}; its combination with a
maximal-subgroup sieve will prove that $\mathcal R_d$ is nonempty
in the families below.

\begin{theorem}[An infinite family of primitive nonrealizable groups]
\label{thm:fermat-obstruction}
Let $r\ge2$, $q=2^{2r}$, and $d=2^r+1$. The scalar affine group
\[
H_d=\mathbb F_q^+\rtimes C_d\le\mathrm{AGL}(1,q)
\]
is primitive of degree $q$, but no affine complete mapping $f$
on $\mathbb F_2^{2r}$ satisfies
$\langle f,F(f)\rangle=H_d$.
The same nonrealizability holds for every
$\mathrm{GL}(2r,2)$-conjugate of this subgroup.
More generally, if an odd integer $d>3$ divides $2^r+1$,
then $\mathcal R_d=\varnothing$ in every characteristic-two
field containing $C_d$.
\end{theorem}
\begin{proof}
Since $2^r\equiv-1\pmod d$, the order of $2$ modulo $d$
divides $2r$. It cannot be at most $r$, as
$0<2^j-1<d$ for $1\le j\le r$.
Every proper divisor of $2r$ is at most $r$, so
$\operatorname{ord}_d(2)=2r$.
Consequently the multiplicative action of $C_d$ on
$\mathbb F_q$ is irreducible over $\mathbb F_2$, and
its affine extension by all translations is primitive.

More generally, suppose $d>3$ divides $2^r+1$ and
$a,1+a\in C_d$. Because $a^d=1$ and
$(1+a)^d=1$, we have $a^{2^r}=a^{-1}$ and
\[
1=(1+a)^{2^r+1}
  =(1+a^{-1})(1+a)=a+a^{-1}.
\]
Multiplying by $a$ yields $a^2+a+1=0$.
Thus $a$ and $1+a=a^2$ both have order $3$,
and $\langle a,1+a\rangle=C_3$.
Since $d>3$, no $a$ belongs to $\mathcal R_d$.
Corollary~\ref{cor:scalar-realization} proves that no scalar
affine companion pair generates $H_d$. Moreover, every affine
permutation belonging to the specified group $H_d$ has precisely
the scalar form $x\mapsto ax+c$, so this excludes arbitrary
affine complete mappings whose generated group equals $H_d$.

Finally, for an invertible $\mathbb F_2$-linear change of
coordinates $S$, one has
$F(SfS^{-1})=SF(f)S^{-1}$. Hence an affine
companion realization of a conjugate subgroup would conjugate
back to one of $H_d$, which is impossible.
\end{proof}

The theorem gives an obstruction beyond irreducibility and
primitivity; it does not claim that all nonrealizable primitive
affine groups have a complement of this form.

\begin{proposition}[A positive lower bound for realizing slopes]
\label{prop:scalar-positive-bound}
Let $q=2^n$, let $d>1$ divide $q-1$, and set $k=(q-1)/d$.
Write $\omega(d)$ for the number of distinct prime divisors of $d$, and put
\[
 b_k=\frac{(k-1)(k-2)}{k^2}.
\]
For every divisor $e$ of $q-1$, with $\ell=(q-1)/e$, one has
\begin{equation}\label{eq:scalar-intersection-bound}
 \left|T_e-\frac{q+1-3\ell}{\ell^2}\right|
 \le b_\ell\sqrt q.
\end{equation}
Consequently,
\begin{equation}\label{eq:scalar-positive-bound}
 |\mathcal R_d|\ge
 \frac{q}{k^2}\left(1-\sum_{p\mid d}\frac1{p^2}\right)
 -\frac{3k-1}{k^2}-(\omega(d)+b_k)\sqrt q,
\end{equation}
where the sum runs over the distinct prime divisors of $d$. In particular,
\begin{equation}\label{eq:scalar-uniform-bound}
 |\mathcal R_d|>
 \frac{3q}{4k^2}-\frac{3k-1}{k^2}
 -(\omega(d)+b_k)\sqrt q.
\end{equation}
\end{proposition}
\begin{proof}
The input is the classical Jacobi-sum calculation
\cite[Theorems~1--2]{vW02}; we include its character-sum proof.
Choose a character $\chi$ of order $\ell$, and extend all its powers,
including the trivial character, by zero at zero. The subgroup indicator gives
\[
 T_e=\frac1{\ell^2}\sum_{i,j=0}^{\ell-1}J(\chi^i,\chi^j),
 \qquad J(\eta,\theta)=\sum_{x\in\mathbb F_q}\eta(x)\theta(1+x).
\]
The term $i=j=0$ is $q-2$. Each of the $2(\ell-1)$ terms with
exactly one nontrivial character is $-1$. The $\ell-1$ terms with
both characters nontrivial and product trivial are also $-1$:
the substitution $u=x/(1+x)$ identifies their sum with the sum of
a nontrivial character over $\mathbb F_q^\times\setminus\{1\}$.

For completeness, fix a nontrivial additive character $\psi$ and set
$G(\eta)=\sum_{x\ne0}\eta(x)\psi(x)$. For nontrivial $\eta$,
substituting $x=ty$ into its squared modulus gives
\[
 |G(\eta)|^2
 =\sum_{t\ne0}\eta(t)\sum_{y\ne0}\psi((t+1)y)
 =(q-1)-\sum_{t\ne0,1}\eta(t)=q.
\]
If $\eta$, $\theta$, and $\eta\theta$ are all nontrivial,
grouping $G(\eta)G(\theta)$ by $z=x+y$ gives
\[
 G(\eta)G(\theta)=J(\eta,\theta)G(\eta\theta).
\]
Here the $z=0$ contribution vanishes by multiplicative-character
orthogonality; for $z\ne0$ put $x=zu$, $y=z(1+u)$.
Thus $|J(\eta,\theta)|=\sqrt q$.
There are $(\ell-1)(\ell-2)$ remaining terms.
Their triangle bound, together with the terms evaluated above, proves
\eqref{eq:scalar-intersection-bound}, also when $\ell=1$.

A pair counted by $T_d$ that does not generate $C_d$ lies in at least
one maximal subgroup $C_{d/p}$, for a prime $p\mid d$. Therefore
\[
 |\mathcal R_d|\ge T_d-\sum_{p\mid d}T_{d/p}.
\]
Applying \eqref{eq:scalar-intersection-bound} to the first term gives
$T_d\ge(q+1-3k)/k^2-b_k\sqrt q$. For the subgroup of index $kp$,
the same bound gives
\[
 T_{d/p}\le\frac{q+1-3kp}{k^2p^2}+b_{kp}\sqrt q
 \le\frac{q}{k^2p^2}+\sqrt q.
\]
This proves \eqref{eq:scalar-positive-bound}.
All the primes dividing $d$ are odd. The elementary telescoping bound
\[
 \sum_{p\mid d}\frac1{p^2}
 \le\sum_{j=1}^\infty\frac1{(2j+1)^2}
 <\frac14\sum_{j=1}^\infty\frac1{j(j+1)}=\frac14
\]
proves \eqref{eq:scalar-uniform-bound}.
\end{proof}

\begin{theorem}[Realization at any fixed odd index]
\label{thm:fixed-index-realization}
Let $k>1$ be odd. For every sufficiently large $n$ satisfying
$k\mid2^n-1$, put $q=2^n$ and $d=(q-1)/k$. Then
\[
 H_d=\mathbb F_q^+\rtimes C_d
\]
is a primitive subgroup of index $k$ in $\mathrm{AGL}(1,q)$ and
is generated by $f$ and $F(f)$ for a scalar affine complete mapping.
An explicit sufficient condition is $d>\sqrt q-1$ and
\begin{equation}\label{eq:fixed-index-sufficient}
 L_{q,k,d}:=\frac{3q}{4k^2}-\frac{3k-1}{k^2}
 -(\omega(d)+b_k)\sqrt q>0.
\end{equation}
Under these conditions the number of realizing maps in the specified
field coordinates is $(q-1)|\mathcal R_d|>(q-1)L_{q,k,d}$.
In particular, each fixed odd index $k>1$ gives an infinite family of
realizable proper primitive affine groups.
\end{theorem}
\begin{proof}
The degree $s=\operatorname{ord}_d(2)$ of the field generated by $C_d$
divides $n$. If $s<n$, then $s\le n/2$ and
$d\mid2^s-1\le\sqrt q-1$, contrary to the assumed inequality.
Hence $s=n$. Proposition~\ref{prop:scalar-positive-bound} and
\eqref{eq:fixed-index-sufficient} imply $\mathcal R_d\ne\varnothing$.
Corollaries~\ref{cor:scalar-realization} and~\ref{cor:affine-actions}
give the realization, primitivity, and exact count. The index is $k$
because the translation subgroup is full and the complement has order
$(q-1)/k$.

To establish the infinite-family assertion without an arithmetic
assumption on $d$, note that $3^{\omega(d)}\le d<2^n$.
Since $3>2^{3/2}$, this yields $\omega(d)<2n/3$.
Also $0\le b_k\le1$. For fixed $k$, the positive term
$3\cdot2^n/(4k^2)$ in \eqref{eq:fixed-index-sufficient} eventually
exceeds the constant term and $(2n/3+1)2^{n/2}$.
Likewise $(2^n-1)/k>2^{n/2}-1$ eventually.
There are infinitely many admissible $n$, namely multiples of
$\operatorname{ord}_k(2)$. The sufficient conditions therefore hold
for every sufficiently large admissible dimension.
\end{proof}

\begin{corollary}[The exact threshold at index three]
\label{cor:index-three-realization}
For even $n\ge4$, the specified subgroup
\[
 \mathbb F_{2^n}^+\rtimes C_{(2^n-1)/3}
\]
is primitive and is generated by a scalar affine complete mapping
and its prescribed companion if and only if $n\ge8$.
Thus index-three proper primitive companion groups exist in every
even dimension from eight onwards.
\end{corollary}
\begin{proof}
For every even $n\ge4$, writing $t=2^{n/2}>2$ gives
$(t^2-1)/3>t-1$. The field-degree argument in
Theorem~\ref{thm:fixed-index-realization} proves primitivity.
At $n=4$ the complement has order $5=2^2+1$, excluded by
Theorem~\ref{thm:fermat-obstruction}. At $n=6$ it has order $21$,
and the polynomial identities following Corollary~\ref{cor:scalar-sieve}
give $\mathcal R_{21}=\varnothing$.

For $k=3$, $b_3=2/9$. Using $\omega(d)<2n/3$, the uniform lower
bound in \eqref{eq:scalar-uniform-bound} is greater than
\[
 \frac{2^n}{12}-\frac89-
 \left(\frac{2n}{3}+\frac29\right)2^{n/2}.
\]
This expression equals $424/3>0$ at $n=14$. After division by $2^{n/2}$,
increasing $n$ by two doubles its positive exponential term,
whereas its linear term increases by $4/3$ and its negative
exponential term decreases in magnitude. Its increase is greater than
$2^{n/2}/12-4/3>0$ for even $n\ge14$, proving positivity there.
For $n=10$, $d=341=11\cdot31$ and $\omega(d)=2$;
for $n=12$, $d=1365=3\cdot5\cdot7\cdot13$ and $\omega(d)=4$.
The bound $L_{q,3,d}$ equals $40/3$ and $632/9$, respectively,
so these dimensions are covered as well.

Finally take $n=8$ and
\[
 \mathbb F_{256}=\mathbb F_2[t]/(m(t)),\qquad
 m(t)=t^8+t^4+t^3+t+1,
 \qquad a=t^5+t^3+t^2+t.
\]
The irreducibility criterion for $m$ is certified by
$t^{256}\equiv t\pmod m$ and
$\gcd(m,t^{16}+t)=1$; these are polynomial identities over $\mathbb F_2$.
Direct reduction modulo $m$ gives
\[
 a^{85}=(1+a)^{85}=1,\qquad
 a^5=t^7+t^5+t^3+t+1\ne1,\qquad
 a^{17}=t^7+t^6+t^5+t^3+t^2+1\ne1.
\]
Since $85=5\cdot17$, these identities show that $a$ has exact order
$85$ and $1+a\in C_{85}$. Thus $a\in\mathcal R_{85}$;
$f(x)=ax+c$ realizes the group for every $c\ne0$ by
Corollary~\ref{cor:scalar-realization}.
This closes all the even dimensions claimed.
\end{proof}

The square-root estimate above is classical; the results specialize it
to the subgroup generated by the linked slopes $a$ and $1+a$.
The sufficient bound is not a general necessary condition for
realizability, and no priority claim is made for the character-sum
method or these particular applications.


\subsection{Quantitative realization and the quarter-power obstruction}

The divisor sieve yields a main term as well as positivity. The next
two results allow the subgroup index to grow with the field size.

\begin{theorem}[Quantitative slope counts]\label{thm:scalar-density}
Let $q=2^n$, $d>1$ divide $q-1$, and $k=(q-1)/d$. Set
\[
 \theta_j(d)=\prod_{p\mid d}(1-p^{-j})\quad(j=1,2),\qquad
 S_{k,d}=\sum_{h\mid\operatorname{rad}(d)}b_{kh},
 \qquad b_\ell=\frac{(\ell-1)(\ell-2)}{\ell^2}.
\]
Here $\operatorname{rad}(d)$ is the product of its distinct prime divisors.
Then
\begin{equation}\label{eq:scalar-density}
 |\mathcal R_d|=\frac{q+1}{k^2}\theta_2(d)
                -\frac3k\theta_1(d)+E_{q,k,d},
 \qquad |E_{q,k,d}|\le S_{k,d}\sqrt q\le2^{\omega(d)}\sqrt q.
\end{equation}
The sharper error coefficient is
\begin{equation}\label{eq:scalar-error-coefficient}
 S_{k,d}=2^{\omega(d)}-\frac3k\prod_{p\mid d}(1+p^{-1})
                    +\frac2{k^2}\prod_{p\mid d}(1+p^{-2}).
\end{equation}
For fixed $k$, as $n$ tends to infinity through admissible dimensions,
\begin{equation}\label{eq:scalar-density-asymptotic}
 |\mathcal R_d|=\frac q{k^2}\theta_2(d)+O_k(q^{1/2+o(1)}),
 \qquad
 \frac{|\mathcal R_d|}{T_d}
       =\theta_2(d)+O_k(q^{-1/2+o(1)}).
\end{equation}
In particular $\theta_2(d)>3/4$: asymptotically at least three quarters
of the slopes satisfying the two subgroup-membership conditions generate
the specified complement jointly.
\end{theorem}
\begin{proof}
Write $T_{d/h}=(q+1-3kh)/(kh)^2+\varepsilon_h$ in
\eqref{eq:realization-mobius}. Proposition~\ref{prop:scalar-positive-bound}
gives $|\varepsilon_h|\le b_{kh}\sqrt q$. Only squarefree $h$ contribute.
The identities
\[
 \sum_{h\mid d}\frac{\mu(h)}{h^j}=\theta_j(d),\qquad
 \sum_{h\mid\operatorname{rad}(d)}h^{-j}
       =\prod_{p\mid d}(1+p^{-j})
\]
give \eqref{eq:scalar-density} and, since $b_\ell=1-3/\ell+2/\ell^2$,
\eqref{eq:scalar-error-coefficient}. All these coefficients are
nonnegative and at most one for integer $\ell\ge1$.

For each $\epsilon>0$, choose a fixed cutoff beyond which $2\le p^\epsilon$.
The bounded contribution of smaller primes can be absorbed into a constant,
so $2^{\omega(d)}\le C_\epsilon d^\epsilon$. Consequently
$2^{\omega(d)}=q^{o(1)}$. The first asymptotic follows from
\eqref{eq:scalar-density}; the omitted terms are bounded for fixed $k$.
Also $T_d=q/k^2+O_k(\sqrt q)$ by
\eqref{eq:scalar-intersection-bound}, which gives the ratio assertion.
Finally $\prod(1-x_i)\ge1-\sum x_i$ for $0\le x_i\le1$, and the
odd-prime bound in Proposition~\ref{prop:scalar-positive-bound} is
strictly less than $1/4$. Thus $\theta_2(d)>3/4$.
\end{proof}

\begin{corollary}[A uniform range of growing indices]
\label{cor:growing-index-realization}
Let $n\ge16$, $q=2^n$, and $k>1$ divide $q-1$. If
\begin{equation}\label{eq:growing-index-range}
 k\le\frac{q^{1/4}}{\sqrt n},
\end{equation}
then $H_{(q-1)/k}$ is a realizable proper primitive affine companion group.
\end{corollary}
\begin{proof}
Put $d=(q-1)/k$. Since $k<\sqrt q+1$, we have $d>\sqrt q-1$,
so the field-degree argument in Theorem~\ref{thm:fixed-index-realization}
gives $\operatorname{ord}_d(2)=n$. The index is odd and at least three;
hence $(3k-1)/k^2<1$. Using $\omega(d)<2n/3$ and $b_k\le1$ gives
\[
 L_{q,k,d}>\left(\frac n{12}-1\right)\sqrt q-1>0
 \qquad(n\ge16).
\]
Theorem~\ref{thm:fixed-index-realization} proves the assertion.
\end{proof}

The exponent $1/4$ has a structural obstruction. It is not merely
an artifact of bounding each Jacobi sum separately.

\begin{theorem}[Norm descent into a proper subfield]
\label{thm:norm-subfield-descent}
Let $1\le s<r$, $s\mid r$, $u\mid2^s-1$, and work in
$L=\mathbb F_{2^{2r}}$. Put
\[
 d=u(2^r+1),\qquad
 g=\begin{cases}
     u,&r/s\text{ even},\\
     u(2^s+1),&r/s\text{ odd}.
   \end{cases}
\]
The two membership conditions descend exactly:
\begin{equation}\label{eq:norm-subfield-descent}
 \{a\in C_d:1+a\in C_d\}
 =\{a\in C_g\subset\mathbb F_{2^{2s}}^\times:1+a\in C_g\}.
\end{equation}
Thus $T_d=T_g$, $\mathcal R_d=\varnothing$, and
$H_d=L^+\rtimes C_d$ is primitive but cannot be generated by an
affine complete mapping and its prescribed companion.
The nonrealizability also holds for every linear conjugate of $H_d$.
\end{theorem}
\begin{proof}
Set $Q=2^r$, $K=\mathbb F_Q$, and $E=\mathbb F_{2^s}\subset K$.
The norm $N_{L/K}(x)=x^{Q+1}$ has kernel $C_{Q+1}$ and maps $C_d$
onto $C_u\subset E^\times$: its image has order
$d/\gcd(d,Q+1)=u$, since $\gcd(u,Q+1)=1$.
If $a,1+a\in C_d$, put $U=N(a)$ and $W=N(1+a)$. Both lie in $E$, and
\[
 a+a^Q=1+U+W\in E,\qquad
 a^2+(1+U+W)a+U=0.
\]
Thus $a$ has degree at most two over $E$ and belongs to
$\mathbb F_{2^{2s}}\subset L$. The same holds for $1+a$.

Write $t=2^s$ and $m=r/s$. Cyclic subgroup intersection gives
$C_d\cap\mathbb F_{t^2}^\times=C_{\gcd(d,t^2-1)}$.
If $m$ is even, $t^m+1\equiv2\pmod{t^2-1}$, so this gcd is $u$.
If $m$ is odd, $t^m+1\equiv t+1\pmod{t^2-1}$; since
$u\mid t-1$ and $\gcd(t-1,t+1)=1$, the gcd is $u(t+1)$.
This proves one inclusion in \eqref{eq:norm-subfield-descent};
the other follows from $C_g\subset C_d$.
In both cases $g<d$, so every pair in the left-hand set generates
a proper subgroup of $C_d$. Therefore $\mathcal R_d=\varnothing$.

The group $C_d$ contains $C_{2^r+1}$, whose action on $L$ is irreducible
by the field-degree proof in Theorem~\ref{thm:fermat-obstruction}.
Hence $H_d$ is primitive. Corollary~\ref{cor:scalar-realization}
excludes scalar realizations, and every affine permutation in the
specified group is scalar. Linear-conjugacy equivariance of $F$
excludes realizations of its linear conjugates as well.
\end{proof}

\begin{corollary}[The quarter-power barrier]
\label{cor:quarter-power-barrier}
For every $s\ge1$, put $t=2^s$, $q=t^4$, and
\[
 k=t+1=q^{1/4}+1,\qquad d=(t-1)(t^2+1)=\frac{q-1}{k}.
\]
Then $H_d$ is a proper primitive nonrealizable affine companion group,
and its membership pairs are exactly
\begin{equation}\label{eq:quarter-subfield-pairs}
 \{a\in C_d:1+a\in C_d\}=\mathbb F_t\setminus\{0,1\}.
 \qquad T_d=t-2,\quad |\mathcal R_d|=0.
\end{equation}
Consequently no exponent $\beta>1/4$ can give universal eventual
realizability for all admissible indices $1<k\le q^\beta$.
Even the range $1<k\le Cq^{1/4}$ fails for every fixed $C>1$.
\end{corollary}
\begin{proof}
Apply Theorem~\ref{thm:norm-subfield-descent} with $r=2s$ and
$u=t-1$. Here $g=t-1$ and $C_g=\mathbb F_t^\times$, yielding
\eqref{eq:quarter-subfield-pairs}. The indices are proper and
$k/q^{1/4}=1+1/t\to1$. Hence these infinitely many exceptions
eventually lie in either range in the statement.
\end{proof}

Together, Corollaries~\ref{cor:growing-index-realization}
and~\ref{cor:quarter-power-barrier} determine the exponent $1/4$
for a universal growing-index guarantee. They leave a factor of
$\sqrt n=\sqrt{\log_2q}$ between the sufficient range and the
obstructed family. In particular, the exact boundary $k\le q^{1/4}$
has not been decided by these arguments.

\subsection{Primitive classification and generating triangles}

\begin{proposition}[Primitive inverse classification and linear conjugacy]
\label{prop:primitive-scalar-classification}
Up to $\mathrm{GL}(n,2)$ conjugacy, the primitive groups generated by
an affine complete mapping and its prescribed companion are exactly
\[
 H_d=\mathbb F_{2^n}^+\rtimes C_d,
 \qquad d\mid2^n-1,\quad \operatorname{ord}_d(2)=n,
 \quad \mathcal R_d\ne\varnothing.
\]
Distinct admissible $d$ give distinct group-conjugacy classes.
For each such $d$, the normalizer in $\mathrm{GL}(n,2)$ has order
$n(2^n-1)$ and is the semilinear scalar group $\Gamma\mathrm L(1,2^n)$.
The affine complete mappings realizing a linear conjugate of $H_d$
form exactly $|\mathcal R_d|/n$ linear-conjugacy classes, each of size
$|\mathrm{GL}(n,2)|$.
\end{proposition}
\begin{proof}
Corollary~\ref{cor:affine-actions} makes the linear part $A$ of every
primitive realization irreducible and its translation nonzero.
The field $\mathbb F_2[A]$ then identifies $V$ with $\mathbb F_{2^n}$;
the complement is generated by $a$ and $1+a$ and has some order $d$.
The characterization follows from Corollary~\ref{cor:scalar-realization}.
Conversely that corollary realizes every listed $H_d$.
Different orders $2^n d$ distinguish the group classes.

A linear transformation normalizes $H_d$ exactly when it normalizes
its stabilizer $C_d$ of zero. Since $\mathbb F_2[C_d]=\mathbb F_{2^n}$,
conjugation by such a transformation induces an automorphism of this
field, hence a power of Frobenius. After removing that power, the
transformation centralizes all scalar multiplications and is itself
scalar. Conversely every map $x\mapsto b x^{2^i}$ normalizes $C_d$.
The normalizer therefore has the asserted order.

For maps realizing the specified $H_d$, normalizer conjugation sends
$(a,c)$ to $(a^{2^i},bc^{2^i})$. Scalar transformations act transitively
on $c\ne0$. Every $a\in\mathcal R_d$ has degree $n$ over $\mathbb F_2$:
a proper subfield containing $a$ would contain $1+a$ and all of $C_d$.
Thus its Frobenius orbit has length $n$, giving $|\mathcal R_d|/n$
normalizer orbits. Conjugacy between two maps with group $H_d$ must
normalize $H_d$, so these are also the full linear-conjugacy classes
when group conjugates are included. Finally, a linear transformation
centralizing $f(x)=ax+c$ centralizes the irreducible multiplier $a$,
so is scalar; fixing $c\ne0$ forces it to be the identity.
All full linear-conjugacy classes consequently have size
$|\mathrm{GL}(n,2)|$.
\end{proof}

For an explicit polynomial encoding, set
\[
 G_e(t)=\gcd(t^e+1,(t+1)^e+1),\qquad
 P_d(t)=\prod_{e\mid d}G_e(t)^{\mu(d/e)}.
\]
Although written as a quotient, $P_d$ is a monic squarefree polynomial
over $\mathbb F_2$ with root set $\mathcal R_d$. Indeed the exponent
of a root whose pair generates $C_h$ is
$\sum_{h\mid e\mid d}\mu(d/e)$, which is one if $h=d$ and zero otherwise.
When $\operatorname{ord}_d(2)=n$, all irreducible factors of $P_d$
have degree $n$ and encode the classes in
Proposition~\ref{prop:primitive-scalar-classification}.
For example the index-three counts in dimensions $8,10,12,14$ give
$3,9,33,123$ map classes, respectively.

The graph interpretation clarifies the distinction between subgroup
membership and joint generation. Write
\[
 \Gamma(q,d)=\operatorname{Cay}(\mathbb F_q^+,C_d),
 \qquad x\sim y\ \Longleftrightarrow\ y-x\in C_d.
\]
This is the usual generalized Paley graph (including the complete
graph when $d=q-1$). Its connectivity and ordinary triangle counts
are classical \cite{LP09,DaMc21}. For a triangle with ordered vertices
$(x,y,z)$, define its ratio subgroup to be
\[
 K(\{x,y,z\})=\left\langle a,1+a\right\rangle,
 \qquad a=\frac{z-x}{y-x}.
\]
This subgroup is independent of the vertex ordering. We call the
triangle \emph{generating} when its ratio subgroup is $C_d$.
The definition uses normalized ratios, not the subgroup generated
by the unnormalized edge differences.

\begin{proposition}[Generating triangles and companion realization]
\label{prop:Paley-generating-triangles}
For $d>1$ dividing $q-1$, $\Gamma(q,d)$ is connected exactly when
$\operatorname{ord}_d(2)=n$. Each edge belongs to $T_d$ ordinary
triangles and $|\mathcal R_d|$ generating triangles. The total number
of generating triangles is $qd|\mathcal R_d|/6$.
Thus $H_d$ is realizable exactly when the graph is connected and
has a generating triangle.
If $d>3$, the action of $H_d$ on generating triangles is free,
with $|\mathcal R_d|/6$ orbits; in particular $6\mid|\mathcal R_d|$.
\end{proposition}
\begin{proof}
The additive span of $C_d$ is closed under products, contains one,
and is a finite subfield. It is therefore $\mathbb F_2[C_d]
=\mathbb F_{2^{\operatorname{ord}_d(2)}}$.
The components are its additive cosets, proving the connectivity
assertion (compare \cite[Theorem~2.2]{LP09}).

For the edge $\{0,1\}$, its third vertices are precisely the $a$
with $a,1+a\in C_d$, and the generating ones are precisely
$a\in\mathcal R_d$. Reordering vertices replaces $a$ by one of
\[
 a,\quad 1+a,\quad a^{-1},\quad(1+a)^{-1},\quad
 \frac{a}{1+a},\quad\frac{1+a}{a};
\]
these all generate the same subgroup together with one plus themselves.
Translations and multiplication by $C_d$ preserve the ratios and act
transitively on edges. Counting the $qd/2$ edges and three edges per
triangle gives the total. The realization equivalence now follows
from Corollary~\ref{cor:scalar-realization}.

For $d>3$ the six transforms of $a\in\mathcal R_d$ are distinct.
A transposition can fix only an excluded value $0$ or $1$, whereas
a three-cycle can fix only $a^2+a+1=0$, whose pair generates $C_3$.
Any affine stabilizer of a generating triangle induces a vertex
permutation and would make the corresponding transform fix $a$.
It must therefore fix the vertices individually, forcing its
multiplier to be one and its translation to be zero.
The action is free. Each orbit has size $qd$, giving the orbit count
and divisibility. No claim about the full graph automorphism group
is required here.
\end{proof}

In the family of Corollary~\ref{cor:quarter-power-barrier}, the graph
is connected and every edge has exactly $t-2$ third vertices. For an
edge $\{x,y\}$ they are $x+(y-x)a$ with
$a\in\mathbb F_t\setminus\{0,1\}$. Thus all its triangles lie in the
clique $x+(y-x)\mathbb F_t$, and none is generating.
These graphs are the already studied $(t+1)$-power graphs
$\Gamma_{t,4}(1)$ of Podest\'a and Videla \cite{PV24}; their ordinary
intersection parameter $t-2$ follows also from
\cite[Theorem~5.2]{PV24}. The descent proof supplies the exact subfield
location and its companion-group obstruction. Connectedness and
ordinary triangles alone do not guarantee realization.
The graph family, character-sum method, and ordinary triangle counts
are not claimed as new, and priority of the companion-specific
applications remains to be established.


\subsection{Realization at the inner quarter-power boundary}
\label{subsec:inner-critical}

The obstruction at $k=t+1$ has a positive counterpart at $k=t-1$.
Throughout this subsection put
\[
 t=2^s\ (s\ge2),\quad K=\mathbb F_t,\quad E=\mathbb F_{t^2},
 \quad L=\mathbb F_{t^4},\quad A=t+1,\quad B=t^2+1,\quad d=AB.
\]
Then $\gcd(A,B)=1$ and $C_d=\ker N_{L/K}$.
Write $U=\ker N_{E/K}=C_A$ and $V=\ker N_{L/E}=C_B$.
The next results prove the stronger assertion that some $a\in C_d$
has order $d$ and satisfies $1+a\in C_d$.

\begin{proposition}[Quadratic parametrization]\label{prop:inner-quadratic}
The elements $a\in L\setminus E$ with $a,1+a\in C_d$ correspond,
two at a time, to the ordered pairs $(b,w)\in U^2$ satisfying
\[
 c=1+b+w\ne0,\qquad
 \operatorname{Tr}_{E/\mathbb F_2}(b/c^2)=1.
\]
The two elements are the roots of $X^2+cX+b$, and
$N_{L/E}(a)=b$, $N_{L/E}(1+a)=w$.
The remaining elements, in $E$, are precisely the two roots of
$X^2+X+1$ when $s$ is odd, and there are none when $s$ is even.
These remaining elements generate only $C_3$.
\end{proposition}
\begin{proof}
If $a\notin E$, write $b=a^{t^2+1}$ and $c=a+a^{t^2}$.
Then $c\ne0$, the minimal polynomial is $X^2+cX+b$, and
$(1+a)^{t^2+1}=1+b+c=w$. Norm transitivity puts $b,w$ in $U$.
After division by $c^2$, the quadratic is an Artin--Schreier equation;
it is irreducible exactly when the displayed absolute trace is one.
Conversely, that trace condition makes the quadratic irreducible over
$E$, with both roots in $L\setminus E$. Their two relative norms are
$b,w$, so norm transitivity gives $a,1+a\in C_d$.
The relative norms determine the ordered pair uniquely.

For $a\in E$, the same norm conditions reduce to $a,1+a\in U$, since
$N_{L/K}(a)=N_{E/K}(a)^2$. Thus $a^t=a^{-1}$ and
$(1+a)^t=(1+a)^{-1}$. Substitution gives
$(1+a)(1+a^{-1})=1$, or $a^2+a+1=0$.
Such roots lie in $U$ exactly when $3\mid t+1$, equivalently when
$s$ is odd. Their two norms satisfy $1+b+w=0$.
In particular, the omitted $c=0$ branch contributes no realizing slope.
\end{proof}

\begin{lemma}[A genus-one norm fiber]\label{lem:inner-fiber}
Fix $b\in U$ with $b^3\ne1$ and put
\[
 X_b(K)=\{a\in L:N_{L/E}(a)=b,\ N_{L/K}(1+a)=1\}.
\]
This set is the full set of $K$-points of a smooth projective
geometrically connected curve $C_b/K$ of genus one. Over $\overline K$
the affine part has equation
\begin{equation}\label{eq:inner-fiber}
 (x+1)(x+b)(y+1)(y+b^{-1})=xy
 \qquad (x,y\ne0).
\end{equation}
Its completion in $\mathbb P^1\times\mathbb P^1$ has eight boundary
points, none rational over $K$. Consequently
\[
 \bigl|\#X_b(K)-(t+1)\bigr|\le2\sqrt t.
\]
\end{lemma}
\begin{proof}
The relative-norm equation defines a variety over $K$ by restriction
of scalars. In the four split coordinates of $\operatorname{Res}_{L/K}
\mathbb G_m$, write
$(a_0,a_1,a_2,a_3)=(x,y,b/x,b^{-1}/y)$.
The remaining norm equation is exactly \eqref{eq:inner-fiber}.
Let $P=(x+1)(x+b)$ and $Q=(y+1)(y+b^{-1})$.
At a singular point in the torus, $PQ=xy\ne0$ and
\[
 (b+1)Q+y=0,\qquad (b^{-1}+1)P+x=0.
\]
Multiplication would force $b+b^{-1}=1$, hence $b^3=1$.
This contradicts the hypothesis $b^3\ne1$.
The boundary points are
\[
\begin{split}
 D_x&=\{(0,1),(0,b^{-1}),(\infty,1),(\infty,b^{-1})\},\\
 D_y&=\{(1,0),(b,0),(1,\infty),(b,\infty)\}.
\end{split}
\]
They are distinct and smooth: in the appropriate affine charts the
derivative along the finite coordinate is a nonzero multiple of
$1+b$ or $1+b^{-1}$. No corner of the square is on the curve.
The smooth bidegree-$(2,2)$ curve is geometrically connected:
distinct components with zero intersection would all have to be
fibers of the same projection, contrary to the total bidegree.
It is therefore geometrically irreducible and has genus
$(2-1)(2-1)=1$. Its smooth completion descends to $K$.

The $t$-power Frobenius on these coordinates is
$(x,y)\mapsto(b/y^t,x^t)$. A rational point satisfies $y=x^t$,
which excludes all eight boundary points. Hence
$C_b(K)=X_b(K)$, with the stated interpretation in $L$.
The final bound is the Hasse--Weil bound.
\end{proof}

\begin{lemma}[Uniform character bound on a norm fiber]
\label{lem:inner-character}
Let $b$ be as in Lemma~\ref{lem:inner-fiber}, and let $u\in U$ be
the unique square root of $b$. For every nontrivial multiplicative
character $\chi$ of $V$, one has
\begin{equation}\label{eq:inner-character}
 \left|\sum_{a\in X_b(K)}\chi(a/u)\right|\le8\sqrt t.
\end{equation}
\end{lemma}
\begin{proof}
We use the classical Weil bound for a geometrically nontrivial
finite-order rank-one character on a curve: on a genus-$g$ curve,
tame ramification at $r$ geometric points gives the bound
$(2g-2+r)\sqrt t$, when the character is nontrivial at every one
of those points. This is the curve Artin-character estimate
discussed in \cite[\S4]{FFSTV13}; the Lang-torsor interpretation
of finite-field multiplicative characters is also used in \cite{Katz89}.
For completeness, the dimension underlying this estimate is
$2g-2+r$: compactly supported cohomology in degrees zero and two
vanishes, and the tame Euler-characteristic formula gives this
dimension in degree one. The trace formula and the Weil theorem
for curves bound each eigenvalue by $\sqrt t$.
For finite-order characters this is a character component of a
finite cover of the curve. We now check all these hypotheses here.

Let $\mathcal T=\operatorname{Res}_{E/K}
(\ker N_{L/E})$ denote the two-dimensional norm torus.
It has $\mathcal T(K)=V$, and $a\mapsto a/u$ maps $X_b$ to
$\mathcal T$, since $N_{L/E}(u)=u^2=b$.
Over $\overline K$ use coordinates $(z_0,z_1)$ on the torus;
its Frobenius is
\[
 F_{\mathcal T}(z_0,z_1)=(z_1^{-t},z_0^t).
\]
The Lang isogeny $z\mapsto F_{\mathcal T}(z)/z$ has constant
kernel $V$. Pulling it back along $a\mapsto a/u$, and writing
$r_0=x/u$, $r_1=y/u^t$, gives
\[
 \frac{z_1^{-t}}{z_0}=r_0,\qquad
 \frac{z_0^t}{z_1}=r_1,
 \qquad\text{hence}\qquad
 z_0^B=\frac{r_1^t}{r_0}.
\]
Thus this pullback is geometrically a Kummer cover with defining
function $h=(y/u^t)^t/(x/u)=y^t/x$, since $u^{t^2}=u$.
For a character of exact order $m\mid B$, its quotient cover has
equation $z^m=h$ over $\overline K$.
This geometric Kummer presentation must be distinguished from the
$K$-structure of the quotient cover: the latter is inherited from
the pullback of the $K$-defined Lang isogeny, with its corresponding
Frobenius descent datum. In particular, $h$ need not belong to the
function field $K(C_b)$ in these split coordinates. All statements
about geometric ramification and geometric connectedness are checked
after base change to $\overline K$, whereas the character sum is
taken using the descended $K$-cover.

The Frobenius descent can be made explicit, so that the character
trace is not merely inferred from a convention.  On geometric points,
the $K$-Frobenius on the base and on the pullback Lang cover is
\begin{equation}\label{eq:inner-lang-frobenius}
 \begin{aligned}
 F_{C_b}(x,y)&=\left(\frac b{y^t},x^t\right),\\
 F_Y(x,y,z_0)&=\left(\frac b{y^t},x^t,\frac xu z_0\right).
 \end{aligned}
\end{equation}
Indeed, the two Lang equations give $z_1^{-t}=r_0z_0$, while
$F_{\mathcal T}$ acts on the first split coordinate as
$z_0\mapsto z_1^{-t}$.  The displayed formula is also compatible
with the Kummer equation, since $u^B=b$ and
\[
 h\left(\frac b{y^t},x^t\right)
  =\frac{x^{t^2}y^t}{b}
  =\left(\frac xu\right)^B h(x,y).
\]
For a $K$-rational point, $(x,y)=(a,a^t)$ with $a^B=b$;
thus Frobenius fixes the base point and acts on its geometric
Lang fiber by $z_0\mapsto(a/u)z_0$.  This is precisely the deck
transformation indexed by $a/u\in V$.  In the character eigensheaf
associated with $\chi$, its Frobenius trace is therefore
$\chi(a/u)$, or $\chi^{-1}(a/u)$ if the inverse Frobenius convention
is used.  The latter choice is immaterial for the asserted bound.
In particular, this computation verifies the Frobenius trace
identification for the actual $K$-descent of the cover, rather
than for the bare geometric Kummer equation.

The zeros and poles of $x$ are simple and supported on $D_x$;
those of $y$ are simple and supported on the disjoint set $D_y$.
The valuations can be read point by point:
\[
 \begin{array}{c|c}
 \text{points}&\operatorname{ord}_P(h)\\\hline
 (0,1),\ (0,b^{-1})&-1\\
 (\infty,1),\ (\infty,b^{-1})&1\\
 (1,0),\ (b,0)&t\\
 (1,\infty),\ (b,\infty)&-t
 \end{array}
\]
Thus the eight nonzero valuations are $1,-1,t,-t$, each twice.
Since $m$ is odd and $\gcd(m,t)=1$, the character is tamely and
nontrivially ramified at all eight points. A valuation of absolute
value one also makes this cover geometrically connected, excluding
a constant-field character or a geometrically trivial pullback.
The conductor degree is consequently exactly eight, independently
of $m,t,b$. Lemma~\ref{lem:inner-fiber} and the curve bound now give
$(2\cdot1-2+8)\sqrt t=8\sqrt t$.
\end{proof}


\begin{proposition}[Four-torsion symmetry of a norm fiber]
\label{prop:inner-four-torsion}
In the setting of Lemma~\ref{lem:inner-fiber}, let $u\in U$ satisfy
$u^2=b$. On the geometric $(x,y)$-model of $C_b$ set
\[
 \rho(x,y)=\left(\frac1y,\frac xb\right),\qquad
 h=\frac{(y/u^t)^t}{x/u}=\frac{y^t}{x}.
\]
Then $\rho$ descends to a $K$-automorphism of order four. After
choosing a $K$-rational origin, it is translation by a $K$-rational
point of exact order four. The eight boundary points form exactly
two $\rho$-orbits of length four:
\[
 \begin{split}
 (0,1)&\longmapsto(1,0)\longmapsto(\infty,b^{-1})
     \longmapsto(b,\infty)\longmapsto(0,1),\\
 (0,b^{-1})&\longmapsto(b,0)\longmapsto(\infty,1)
     \longmapsto(1,\infty)\longmapsto(0,b^{-1}).
 \end{split}
\]
Moreover, as identities in the geometric function field,
\begin{equation}\label{eq:inner-rho-kummer}
 \rho^*h=(uy)^B h^{-t},\qquad (\rho^2)^*h=h^{-1}.
\end{equation}
For every character $\chi$ of $V$, the sums in
Lemma~\ref{lem:inner-character} satisfy
\[
 S_b(\chi):=\sum_{a\in X_b(K)}\chi(a/u)
      =S_b(\chi^{-t})=S_b(\chi^{-1})\in\mathbb R.
\]
Every nontrivial character belongs to an orbit of length four
under $\chi\mapsto\chi^{-t}$. If
$N_b^{\mathrm{prim}}$ denotes the number of $a\in X_b(K)$
with $\operatorname{ord}(a)=d$, then
$4\mid\#X_b(K)$ and $4\mid N_b^{\mathrm{prim}}$.
\end{proposition}
\begin{proof}
The field automorphism $\sigma:a\mapsto a^t$ of $L/K$
extends as $\sigma\otimes1$ to an $R$-linear automorphism of
$L\otimes_K R$ for every $K$-algebra $R$;
composing this extension with inversion gives a $K$-automorphism
of the ambient restriction-of-scalars torus. On $K$-points its
action is $a\mapsto(a^t)^{-1}$. Here this notation refers to the
Galois automorphism $\sigma$, not the relative Frobenius morphism
on the torus. The functorial norm identities
\[
 N_{L/E}((\sigma a)^{-1})=b^{-t}=b,\qquad
 N_{L/K}(1+(\sigma a)^{-1})
   =\frac{N_{L/K}(1+a)}{N_{L/K}(a)}=1
\]
show that it preserves both defining conditions, since
$N_{L/K}(a)=N_{E/K}(b)=1$. It gives the displayed $\rho$
in split coordinates. Equivalently, the split-coordinate
$K$-Frobenius descent is $(x,y)\mapsto(b/y^t,x^t)$, with which
$\rho$ commutes. Direct
iteration gives
\[
 \rho^2(x,y)=(b/x,b^{-1}/y),\qquad \rho^4=1.
\]
The eight boundary points are paired without fixed points by
$\rho^2$. Any remaining geometric fixed point would have
$x^2=b$, $y^2=b^{-1}$, hence $(x,y)=(u,u^{-1})$.
Substitution in the equation of $C_b$ forces
$b+b^{-1}=1$, equivalently $u^2+u+1=0$, whence $b^3=1$,
contrary to $b^3\ne1$. Thus $\rho^2$,
and hence $\rho$, has no geometric fixed point. The bound
$\#C_b(K)\ge t+1-2\sqrt t>0$ supplies a rational origin.
After choosing a rational origin, any automorphism of a genus-one
curve is a translation followed by an origin-preserving
automorphism. If the latter is nontrivial, its difference from the
identity is a nonzero, hence surjective, elliptic-curve
endomorphism, so the composite has a geometric fixed point.
Consequently $\rho$ is a translation, necessarily by a
$K$-rational point of exact order four.

The two boundary cycles follow by substitution in the rational
formula, interpreted on $\mathbb P^1\times\mathbb P^1$. Since
$u^{t+1}=1$, $u^B=u^2=b$ and $b^t=b^{-1}$,
\[
 \rho^*h=\frac{(x/b)^t}{1/y}
   =b x^t y=(uy)^B\left(\frac{y^t}{x}\right)^{-t},
\]
while the formula for $\rho^2$ yields $(\rho^2)^*h=h^{-1}$.
For $a=uv$ with $v\in V$, $\rho$ sends $a$ to
$uv^{-t}$. It therefore permutes $X_b(K)$ and preserves
$\operatorname{ord}(a)$; changing variables in the character sum
proves the stated equalities. Finally $B=t^2+1$ is odd,
$(-t)^2\equiv-1\pmod B$. For a character of order $m>1$,
an orbit of length one or two would imply $m\mid2$,
which is impossible since $m\mid B$ is odd. Thus every
nontrivial character has an orbit of length four. Translation
by a point of exact order four acts freely, proving both
divisibilities.
\end{proof}

These symmetries do not by themselves reduce the bound in
Lemma~\ref{lem:inner-character}: grouping the characters into
four-element orbits leaves the total absolute sieve coefficients
unchanged. The estimate used below remains $8\sqrt t$.

\begin{theorem}[Primitive slopes in the inner critical family]
\label{thm:inner-primitive-count}
Put $\theta(B)=\varphi(B)/B$ and $W(B)=2^{\omega(B)}$.
For each $b\in U$ of order $A$, let
\[
 N_b^{\mathrm{prim}}=\#\{a\in X_b(K):\operatorname{ord}(a)=d\}.
\]
Then
\begin{equation}\label{eq:inner-fiber-count}
 \bigl|N_b^{\mathrm{prim}}-\theta(B)(t+1)\bigr|
 \le\theta(B)(8W(B)-6)\sqrt t.
\end{equation}
In particular, with
$N_d^{\mathrm{prim}}=\#\{a\in C_d:\operatorname{ord}(a)=d,\ 1+a\in C_d\}$,
\begin{equation}\label{eq:inner-total-count}
 \bigl|N_d^{\mathrm{prim}}-\varphi(A)\theta(B)(t+1)\bigr|
 \le\varphi(A)\theta(B)(8W(B)-6)\sqrt t.
\end{equation}
Thus $N_d^{\mathrm{prim}}\sim\varphi(A)\theta(B)t$ as $s\to\infty$.
\end{theorem}
\begin{proof}
Every $a\in X_b(K)$ has a unique expression $a=uv$ with
$v\in V$, where $u$ is the square root of $b$ in $U$.
Since $\operatorname{ord}(u)=A$ and $\gcd(A,B)=1$, it has order
$d$ exactly when $v$ has order $B$.
The usual character indicator in a cyclic group of order $B$ is
\[
 \mathbf1_{\operatorname{ord}(v)=B}
 =\theta(B)\sum_{m\mid B}\frac{\mu(m)}{\varphi(m)}
   \sum_{\operatorname{ord}(\chi)=m}\chi(v).
\]
Indeed the sum factors over primes $p\mid B$: the corresponding
factor is zero on $V^p$ and $p/(p-1)$ off $V^p$, which cancels
the factor $(p-1)/p$ in $\theta(B)$.
The trivial-character term is $\theta(B)\#X_b(K)$.
For each nontrivial squarefree divisor $m$ the sum of the absolute
values of its coefficients is one. There are $W(B)-1$ such
divisors. Lemmas~\ref{lem:inner-fiber} and~\ref{lem:inner-character}
therefore give an error of at most
$\theta(B)(2+8(W(B)-1))\sqrt t$, proving
\eqref{eq:inner-fiber-count}.

An element of order $d$ has relative norm of order $A$, because
$N_{L/E}(a)=a^B$. Hence these $\varphi(A)$ norm fibers partition
the elements counted by $N_d^{\mathrm{prim}}$. Summing proves
\eqref{eq:inner-total-count}. Finally $2^{\omega(N)}=N^{o(1)}$:
for any $\epsilon>0$, all but finitely many primes satisfy
$2\le p^\epsilon$, and the remaining primes contribute only a
constant. Since $B\asymp t^2$, the relative error tends to zero.
\end{proof}

\begin{theorem}[All-dimensional realization at index $t-1$]
\label{thm:inner-all}
For every $s\ge2$, there exists $a\in\mathbb F_{2^{4s}}$ with
\[
 \operatorname{ord}(a)=(2^s+1)(2^{2s}+1),\qquad
 (1+a)^{(2^s+1)(2^{2s}+1)}=1.
\]
Consequently $\mathcal R_d\ne\varnothing$ at $k=t-1$ for every
such $s$. The proper primitive group $L^+\rtimes C_d$ is generated
by $f(x)=ax+c$ and $F(f)$ for every $c\in L^\times$.
\end{theorem}
\begin{proof}
Every prime divisor $p$ of $B=t^2+1$ is $1$ modulo $4$, since
$t^2\equiv-1\pmod p$. If $p\ge256$, then $2\le p^{1/8}$.
For the smaller primes a direct integer product gives
\[
 \prod_{\substack{p<256\ \mathrm{prime}\\p\equiv1\pmod4}}
 \frac{256}{p}<32^8.
\]
The 24 primes in this product and the exact integer comparison
are supplied with the finite verifier in Appendix~\ref{app:inner-cert}.
It follows for every $B$ under consideration that
\[
 W(B)<32B^{1/8}.
\]
For $s\ge33$, $t^2\ge2^{66}$ and $t^2+1<2t^2$ imply
$t^4>2^{64}(t^2+1)$. Equivalently,
$\sqrt t>256B^{1/8}>8W(B)$.
The lower bound in \eqref{eq:inner-fiber-count} is then positive
for every primitive $b\in U$.

For $2\le s\le32$, Appendix~\ref{app:inner-cert} gives explicit
field towers $E=\mathbb F_2[X]/(M_s)$ and
$L=E[Z]/(Z^2+cZ+b)$, in which $a=Z$ is a certificate.
Irreducibility, both norm conditions and the exact order $d$
are verified by deterministic finite arithmetic, with complete
factorizations and recursively certified primes. These 31 cases
exhaust the remaining range; they are not an extrapolation.

Finally $d>t^2-1$ implies $\mathbb F_2[a]=L$: every proper
subfield of $L$ has at most $t^2$ elements. Thus $C_d$ acts
irreducibly, while $k=t-1>1$ makes it a proper multiplier group.
The affine conclusion follows from
Corollary~\ref{cor:scalar-realization}.
\end{proof}

The two adjacent indices $t-1$ and $t+1$ therefore have opposite
realization behavior for every $s\ge2$.
This does not settle the universal question for all admissible
$k\le q^{1/4}$. Also, $N_d^{\mathrm{prim}}$ counts primitive individual slopes;
it is a subset count within $|\mathcal R_d|$. The joint and
double full-order counts below require the mixed-character estimate.
Ordinary intersections for index $t-1$ belong to cyclotomic
number theory, including the general norm-based treatment of
Kudo and Nogata \cite{KN26}. The additional exact-order condition
is the focus here. The classical Weil bound, Lang construction,
and primitive-character sieve are attributed above; priority
for this particular realization theorem is not asserted.


\paragraph{Mixed characters and joint densities.}
We now retain all three varying character components on a norm fiber.
For $a\in X_b(K)$ put
\begin{equation}\label{eq:inner-mixed-variables}
 w=(1+a)^B\in U,\qquad v=a^A\in V,\qquad v_+=(1+a)^A\in V.
\end{equation}
\begin{auxlemma}[Character decomposition]
The map $C_d\longrightarrow U\times V$, $z\mapsto(z^B,z^A)$,
is an isomorphism. Every pair $\xi,\zeta\in\widehat{C_d}$ has a
unique decomposition
\[
 \xi(z)=\alpha(z^B)\chi(z^A),\qquad
 \zeta(z)=\eta(z^B)\psi(z^A),
 \quad \alpha,\eta\in\widehat U,\quad\chi,\psi\in\widehat V.
\]
Define the two-character sum and the norm-character exception set by
\begin{equation}\label{eq:inner-pair-sum}
 \begin{aligned}
 S_b(\xi,\zeta)&=\sum_{a\in X_b(K)}\xi(a)\zeta(1+a),\\
 \mathcal E_{\mathrm{norm}}
 &=\{(\alpha\circ(z\mapsto z^B),1):\alpha\in\widehat U\}.
 \end{aligned}
\end{equation}
For every $b\in U$ the sum is
\[
 S_b(\xi,\zeta)=\alpha(b)
       \sum_{a\in X_b(K)}\eta(w)\chi(v)\psi(v_+).
\]
Moreover $(\xi,\zeta)\in\mathcal E_{\mathrm{norm}}$ if and only if
$\eta=\chi=\psi=1$.
\end{auxlemma}
\begin{proof}
The two powers have orders dividing $A$ and $B$, respectively.
An element in the kernel has both its $A$-th and $B$-th powers equal
to one, hence is one because $\gcd(A,B)=1$.
The source and target have the same order $AB$, so the map is an
isomorphism. Taking characters gives the unique decomposition.
On the fiber $a^B=b$; substitution gives the identity for the sum.
Uniqueness also proves the stated equivalence for the exception set.
\end{proof}

\begin{proposition}[Mixed characters on a norm fiber]
\label{prop:inner-mixed-character}
Let $b\in U$ satisfy $b^3\ne1$. For characters
$\eta\in\widehat U$ and $\chi,\psi\in\widehat V$, not all trivial,
\begin{equation}\label{eq:inner-mixed-bound}
 \left|\sum_{a\in X_b(K)}\eta(w)\chi(v)\psi(v_+)\right|
 \le8\sqrt t.
\end{equation}
The product Lang cover used for this sum is geometrically connected,
and hence its smooth projective normalization is geometrically irreducible.
Every character quotient is tame, with conductor degree at most eight.
\end{proposition}
\begin{proof}
Let $C_b^\circ=C_b\setminus(D_x\cup D_y)$, with the eight points of
Lemma~\ref{lem:inner-fiber}. This open curve has
$C_b^\circ(K)=X_b(K)$. Write $\sigma\in\operatorname{Gal}(L/K)$
for $z\mapsto z^t$, extended as $\sigma\otimes1$ on every
$L\otimes_K R$. The algebraic map
\[
 \nu(z)=z\,\sigma(z)
\]
sends the norm-one torus $\ker N_{L/K}$ to
$\mathcal T=\operatorname{Res}_{E/K}(\ker N_{L/E})$:
\[
 N_{L/E}(\nu(z))=N_{L/E}(z)\,\sigma(N_{L/E}(z))=N_{L/K}(z)=1.
\]
Its value on $K$-points is $z^A$. This construction is needed for
the varying norm of $1+a$: exponentiation by $A$ on geometric
points would give relative norm $w^A$, which is not the constant
function one.

Set $w_0=(x+1)(x+b)/x$. The split coordinates of the three
$K$-defined maps to the norm tori are
\[
\begin{array}{c|cc}
 &\text{first coordinate}&\text{second coordinate}\\\hline
 N_{L/E}(1+a)&w_0&w_0^{-1}\\
 \nu(a)&xy&by/x\\
 \nu(1+a)&(x+1)(y+1)&(y+1)(x+b)/x
\end{array}
\]
The first row is a map to $\mathcal U=\ker N_{E/K}$, and
$\mathcal U(K)=U$, $\mathcal T(K)=V$.
On $\mathcal U$ the Frobenius is $z\mapsto z^{-t}$; its Lang
equation is $z^{-A}=w_0$. On $\mathcal T$ the Frobenius is
$(z_0,z_1)\mapsto(z_1^{-t},z_0^t)$; the Lang equations
$z_1^{-t}/z_0=r_0$, $z_0^t/z_1=r_1$ give
$z_0^B=r_1^t/r_0$.
Consequently the product pullback is, over $\overline K$,
\begin{equation}\label{eq:inner-mixed-cover}
 z_U^A=w_0^{-1},\qquad z_1^B=g_1,\qquad z_2^B=g_2,
\end{equation}
where
\begin{equation}\label{eq:inner-mixed-functions}
 g_1=b^{-1}\frac{y^{t-1}}{x^A},\qquad
 g_2=w_0^t\frac{(y+1)^{t-1}}{(x+1)^A}.
\end{equation}
For the second formula, substitute
$x+b=xw_0/(x+1)$ in $r_1^t/r_0$.

These geometric equations inherit their $K$-structure from the
Lang pullbacks. More explicitly, above
$F_{C_b}(x,y)=(b/y^t,x^t)$ their Frobenius sends the three
displayed covering coordinates to
\begin{equation}\label{eq:inner-mixed-frobenius}
 \bigl(w_0z_U,\;xy z_1,\;(x+1)(y+1)z_2\bigr).
\end{equation}
This follows from $F(z)=r z$ in the Lang equations.
It is compatible with \eqref{eq:inner-mixed-cover}, since
\[
 w_0(F_{C_b})=w_0^{-t},\quad
 g_1(F_{C_b})=(xy)^B g_1,\quad
 g_2(F_{C_b})=((x+1)(y+1))^B g_2.
\]
At a rational point $(x,y)=(a,a^t)$, the three deck multipliers
in \eqref{eq:inner-mixed-frobenius} are precisely $w,v,v_+$.
Thus the character component has trace $\eta(w)\chi(v)\psi(v_+)$;
the inverse Frobenius convention simultaneously inverts the
characters and gives the same bound.

We check geometric connectedness and the conductor directly.
The divisors of $w_0,g_1,g_2$ are supported on the eight boundary
points. Their valuations are
\begin{equation}\label{eq:inner-mixed-valuations}
\begin{array}{c|rrr}
 P&\operatorname{ord}_P(w_0)&\operatorname{ord}_P(g_1)&
       \operatorname{ord}_P(g_2)\\\hline
 (0,1)&-1&-A&-1\\
 (0,b^{-1})&-1&-A&-t\\
 (\infty,1)&-1&A&t\\
 (\infty,b^{-1})&-1&A&1\\
 (1,0)&1&t-1&-1\\
 (b,0)&1&t-1&t\\
 (1,\infty)&1&-(t-1)&-t\\
 (b,\infty)&1&-(t-1)&1
\end{array}
\end{equation}
Indeed $x,y$ have simple zeros and poles on $D_x,D_y$;
$x+1$ has simple zeros at $(1,0),(1,\infty)$ and simple poles
at the two points with $x=\infty$, and the corresponding assertion
holds for $y+1$. Substitution in \eqref{eq:inner-mixed-functions}
gives each entry of the table.

For every $m\mid B$, a relation $g_1^r g_2^s=f^m$ in
$\overline K(C_b)^\times$ would imply, at $(0,1)$ and
$(0,b^{-1})$,
\[
 m\mid Ar+s,\qquad m\mid Ar+ts.
\]
Now $\gcd(B,t-1)=\gcd(t^2+1,t-1)=1$ and $\gcd(A,B)=1$.
Subtracting forces $m\mid s$, and then $m\mid r$.
The valuation $-1$ of $w_0$ likewise excludes every nontrivial
Kummer power relation for an order dividing $A$.
One can verify simultaneous independence in a common exponent:
\[
 w_0^{Be_0}g_1^{Ae_1}g_2^{Ae_2}
       \in\overline K(C_b)^{\times AB}
 \quad\Longrightarrow\quad
 A\mid e_0,\ B\mid e_1,\ B\mid e_2.
\]
Reduction of the first two valuations modulo $A$ gives $A\mid e_0$;
their difference modulo $B$ gives $B\mid e_2$, and either valuation
then gives $B\mid e_1$.
This is independence of the relevant functions modulo powers in
$\overline K(C_b)^\times$, verified by their local valuations.

At $P_1=(0,1)$ and $P_2=(0,b^{-1})$, the geometric inertia
elements, written in $G=\mathbb Z/A\mathbb Z\times
\mathbb Z/B\mathbb Z\times\mathbb Z/B\mathbb Z$, are
\[
 v_1=(-1,-A,-1),\qquad v_2=(-1,-A,-t).
\]
Their difference is $(0,0,t-1)$, which generates the third
factor because $\gcd(t-1,B)=1$. Combining a suitable multiple of
this element with $v_1$ gives $(-1,-A,0)$. Its order is $AB$:
the first coordinate has order $A$, the second has order $B$
since $\gcd(A,B)=1$. Thus $v_1,v_2$ generate all of $G$.
Equivalently, no nontrivial character annihilates both inertia
elements. The subgroup generated by geometric inertia is therefore
the full deck group, proving geometric connectedness of the product
cover and of every character quotient.

All covering degrees are odd, so wild inertia is trivial.
For a rank-one character its local Artin conductor is one if
the inertia character is nontrivial, and zero otherwise.
Thus its conductor degree is the number $r$ of active rows of
\eqref{eq:inner-mixed-valuations}, and $r\le8$.
Cancellation can occur at individual rows; we do not presume
ramification at every boundary point. For a nontrivial character
triple the inertia test proves geometric nontriviality and $r>0$.
The genus-one Artin-character Weil bound now gives
$(2\cdot1-2+r)\sqrt t\le8\sqrt t$; see
\cite[\S4, Lemma~1]{FFSTV13} and the Lang construction in
\cite[pp.~198--199]{Katz89}.
In cohomological terms, extend over the inactive boundary points.
On the complement of the $r$ active points, $H_c^0,H_c^2$ vanish
and the tame Euler-characteristic formula gives
$\dim H_c^1=2g-2+r=r$.
The eigenvalues have absolute value at most $\sqrt t$, by the
Weil theorem on the character component of the finite cover.
No active boundary point is $K$-rational, so this trace sum is
exactly the sum over $X_b(K)$, with no rational boundary correction.
\end{proof}

\begin{corollary}[The precise two-character exception]
\label{cor:inner-mixed-exception}
Let $b\in U$ satisfy $b^3\ne1$ and let $\xi,\zeta\in\widehat{C_d}$.
Use the character decomposition lemma and \eqref{eq:inner-pair-sum}.
If $(\xi,\zeta)\in\mathcal E_{\mathrm{norm}}$, then
\[
 S_b(\xi,\zeta)=\alpha(b)\#X_b(K).
\]
In every other case $|S_b(\xi,\zeta)|\le8\sqrt t$.
In particular, the assertion for every nontrivial two-character
pair is false: a nontrivial $\alpha\circ N_{L/E}$ with $\zeta=1$
is constant on the norm fiber.
\end{corollary}
\begin{proof}
The decomposition lemma identifies the exception set with the
trivial varying triple. That triple gives the exact constant sum.
Every pair outside the set has a nontrivial varying triple, so
Proposition~\ref{prop:inner-mixed-character} gives the bound.
\end{proof}

\begin{auxlemma}[Local averaging in a cyclic group]
Put
\[
 \begin{aligned}
 \theta_2(n)&=\prod_{p\mid n}(1-p^{-2}),&
 J_b&=\{p\mid A:b\in U^p\},\\
 \lambda_A(b)&=\prod_{p\in J_b}(1-p^{-1}),&
 H(A)&=\prod_{p\mid A}\frac{p+2}{p+1}.
 \end{aligned}
\]
The following are exact averages over all $A$ elements of $U$:
\begin{equation}\label{eq:inner-local-averages}
 \begin{aligned}
  \frac1A\sum_{b\in U}\lambda_A(b)&=\theta_2(A),\\
  \frac1A\sum_{b\in U}\lambda_A(b)2^{|J_b|}
       &=\theta_2(A)H(A).
 \end{aligned}
\end{equation}
In particular the average of $\lambda_A(b)\theta_2(B)$ is
$\theta_2(A)\theta_2(B)=\theta_2(d)$.
\end{auxlemma}
\begin{proof}
Fix a generator $u$ and write $b=u^j$, $j\in\mathbb Z/A\mathbb Z$.
For $p\mid A$ the condition $b\in U^p$ is equivalent to $p\mid j$.
Reduction of uniform $j$ modulo $\operatorname{rad}(A)$ is uniform;
the Chinese remainder theorem makes the residues modulo the distinct
primes independent. The probability of $p\in J_b$ is therefore $1/p$.
The local expectation for $\lambda_A$ is
\[
 (1-1/p)+\frac1p(1-1/p)=1-p^{-2}.
\]
For $\lambda_A(b)2^{|J_b|}$ it is
\[
 (1-1/p)+\frac2p(1-1/p)
       =(1-p^{-2})\frac{p+2}{p+1}.
\]
Multiplication of these independent local expectations proves both
identities. Since $\gcd(A,B)=1$, their prime supports are disjoint,
which proves the last equality.
\end{proof}

\begin{theorem}[Joint generation density at the inner boundary]
\label{thm:inner-joint-density}
Use $\theta_2,\lambda_A,J_b,H(A)$ from the local averaging lemma,
and put $W(n)=2^{\omega(n)}$, $r_0=\gcd(3,A)\le3$.
Let $\mathcal R_{d,b}=\mathcal R_d\cap X_b(K)$.
For every $b\in U$ with $b^3\ne1$,
\begin{equation}\label{eq:inner-joint-fiber}
 \bigl||\mathcal R_{d,b}|-\lambda_A(b)\theta_2(B)A\bigr|
 \le\lambda_A(b)\theta_2(B)
       \bigl(8\,2^{|J_b|}W(B)-6\bigr)\sqrt t.
\end{equation}
For the full family,
\begin{equation}\label{eq:inner-joint-total}
 \bigl||\mathcal R_d|-\theta_2(d)A^2\bigr|
 \le\theta_2(d)A\bigl(8H(A)W(B)-6\bigr)\sqrt t+2r_0A.
\end{equation}
Since $H(A)\le W(A)$, this implies
\[
 |\mathcal R_d|=\theta_2(d)t^2+O(t^{3/2}W(d)),\qquad
 |\mathcal R_d|\sim\theta_2(d)t^2.
\]
The fixed-fiber population is $O(t)$, and its coefficient depends
on $b$. If $b$ is primitive in $U$, then $J_b$ is empty and
the coefficient is $\theta_2(B)$.
Moreover $|T_d-A^2|\le2A\sqrt t+r_0A$, so the actual proportion
of jointly generating slopes among the ordinary intersection is
$|\mathcal R_d|/T_d=\theta_2(d)+O(W(d)/\sqrt t)$.
\end{theorem}
\begin{proof}
The group isomorphism gives
\[
 \langle a,1+a\rangle=C_d
 \quad\Longleftrightarrow\quad
 \langle b,w\rangle=U\ \text{and}\ \langle v,v_+\rangle=V.
\]
For $p\mid A$ with $b\notin U^p$ the first condition is automatic.
For $p\in J_b$ it requires $w\notin U^p$. By orthogonality its
indicator is the product of the factors
\[
 1-\frac1p\sum_{\eta^p=1}\eta(w),\qquad p\in J_b.
\]
Each factor has constant coefficient $1-1/p$ and total absolute
coefficient sum $2(1-1/p)$.
The second condition has indicator
\[
 \prod_{p\mid B}\left(1-\frac1{p^2}
        \sum_{\chi^p=\psi^p=1}\chi(v)\psi(v_+)\right).
\]
Here each factor has constant coefficient $1-1/p^2$ and total
absolute coefficient sum $2(1-1/p^2)$.
Characters belonging to different prime-primary components have
unique products, so multiplying these expansions gives constant
coefficient $c_b=\lambda_A(b)\theta_2(B)$ and total absolute
coefficient sum $c_b2^{|J_b|}W(B)$.
The Hasse error for the trivial triple is $2c_b\sqrt t$;
Proposition~\ref{prop:inner-mixed-character} bounds all other
terms by $8\sqrt t$. This proves \eqref{eq:inner-joint-fiber}.

The local averaging lemma supplies \eqref{eq:inner-local-averages}.
Multiplication by $\theta_2(B)$ turns $\theta_2(A)$ into
$\theta_2(d)$; it also turns the weighted average into
$\theta_2(d)H(A)$. These identities average over all $b\in U$,
before separating the exceptional fibers.

There are precisely $r_0$ exceptional $b$ with $b^3=1$.
For every $b$, each $w\in U$ allows at most two roots of
$X^2+(1+b+w)X+b$. Thus $\#X_b(K)\le2A$ even in an exceptional
fiber. Its contribution differs from the formal main term
$c_bA\in[0,A]$ by at most $2A$.
Sum \eqref{eq:inner-joint-fiber} over the good fibers, bound
the positive error coefficients by their sum over all $b$, and
use the two exact averages. This gives
\eqref{eq:inner-joint-total}. It includes all exceptional and
subfield points, rather than applying the smooth-fiber bound to them.
Finally $W(d)=d^{o(1)}=t^{o(1)}$, while $\theta_2(d)\ge1/2$:
the latter follows by comparison with the telescoping product
$\prod_{n=2}^\infty(1-n^{-2})=1/2$.
Replacing $A^2=(t+1)^2$ by $t^2$ gives the stated asymptotic.
For $T_d$ sum the Hasse bound over the good fibers and use
$|\#X_b(K)-A|\le A$ on each exceptional fiber.
This proves its bound and, by division, the density assertion.
\end{proof}

\begin{theorem}[Double full-order statistics and all-dimensional existence]
\label{thm:inner-double-count}
Define $\theta(n)=\varphi(n)/n$ and
\[
 D_{d,b}=\#\{a\in X_b(K):\operatorname{ord}(a)
                          =\operatorname{ord}(1+a)=d\},\qquad
 D_d=\sum_{b\in U}D_{d,b}.
\]
If $b$ is not primitive in $U$, then $D_{d,b}=0$.
For primitive $b$, put $c=\theta(A)\theta(B)^2$. Then
\begin{equation}\label{eq:inner-double-fiber}
 |D_{d,b}-cA|\le c\bigl(8W(A)W(B)^2-6\bigr)\sqrt t.
\end{equation}
In particular,
\begin{equation}\label{eq:inner-double-total}
 |D_d-\theta(d)^2A^2|
 \le\theta(d)^2A\bigl(8W(A)W(B)^2-6\bigr)\sqrt t,
\end{equation}
and $D_d\sim\theta(d)^2t^2$ as $s\to\infty$.
Define the explicit analytic range
\begin{equation}\label{eq:inner-double-range}
 \mathcal S_{\mathrm{an}}
 =\{s\in\mathbb Z_{\ge2}:A>(8W(A)W(B)^2-6)\sqrt t\},
\end{equation}
with $t=2^s$, $A=t+1$, and $B=t^2+1$.
For every $s\in\mathcal S_{\mathrm{an}}$ and every primitive
$b\in U$, $D_{d,b}>0$. This set contains all sufficiently large
integers $s$. Moreover, $D_d>0$ for every integer $s\ge2$.
\end{theorem}
\begin{proof}
Order $d$ forces $a^B=b$ to have order $A$, which proves the
vanishing assertion. For primitive $b$, the two full-order conditions
are exactly that $w$ is primitive in $U$ and $v,v_+$ are primitive
in $V$. Apply the primitive-character indicator proved in
Theorem~\ref{thm:inner-primitive-count} separately to these three
variables. For a cyclic group of order $n$ it has constant coefficient
$\theta(n)$ and total absolute coefficient sum $\theta(n)W(n)$.
The product therefore has constant coefficient $c$ and total
absolute coefficient sum $cW(A)W(B)^2$.
The trivial triple contributes $c\#X_b(K)$, with Hasse error
$2c\sqrt t$; every other triple has the mixed bound
$8\sqrt t$. This proves \eqref{eq:inner-double-fiber}.
Sum over the $\varphi(A)=A\theta(A)$ primitive norm fibers and
use $\theta(d)=\theta(A)\theta(B)$ to obtain
\eqref{eq:inner-double-total}.
The relative error tends to zero because
$W(A)W(B)^2=t^{o(1)}$. The stated inequality makes the lower
bound positive in every primitive norm fiber.
For all-dimensional existence, Proposition~\ref{prop:inner-diagonal-double}
below gives $D_d\ge\Delta_t>0$ for $s\ge33$, independently of
the preceding three-character sieve. The exact certificates in
Appendix~\ref{app:inner-cert} prove $D_d>0$ for each $2\le s\le32$.
\end{proof}

\begin{proposition}[A trace-one double full-order subfamily]
\label{prop:inner-diagonal-double}
For $t=2^s$, $s\ge2$, use $A=t+1$, $B=t^2+1$ and $d=AB$ as above.
Put
\[
 \Delta_t=\#\{\gamma\in V:\tau=\gamma+\gamma^{-1}\in U,
                 \ \operatorname{ord}(\gamma)=B,
                 \ \operatorname{ord}(\tau)=A\}.
\]
Every element counted by $\Delta_t$ gives, injectively,
\begin{equation}\label{eq:inner-diagonal-construction}
 a=(1+\gamma)^{-1}\in C_d,\qquad
 1+a=a^{t^2},\qquad
 \operatorname{ord}(a)=\operatorname{ord}(1+a)=d.
\end{equation}
In particular $D_d\ge\Delta_t$. For any characters $\eta$ of $U$
and $\chi$ of $V$ with $(\eta,\chi)\ne(1,1)$,
\begin{equation}\label{eq:inner-diagonal-character}
 \left|\sum_{\substack{\gamma\in V\\
                       \gamma+\gamma^{-1}\in U}}
       \eta(\gamma+\gamma^{-1})\chi(\gamma)\right|\le4\sqrt t.
\end{equation}
Consequently
\begin{equation}\label{eq:inner-diagonal-count}
 |\Delta_t-\theta(d)(t+1)|
       \le\theta(d)\bigl(4W(d)-2\bigr)\sqrt t.
\end{equation}
The sufficient condition $\sqrt t\ge4W(d)$ implies $\Delta_t>0$;
it holds for every $s\ge33$ by the integer bounds in
Appendix~\ref{app:inner-cert}.
\end{proposition}
\begin{proof}
\textit{The full-order construction.}\par
Here $\gamma^{t^2}=\gamma^{-1}$ and $\gamma\ne1$.
Thus $a^{t^2}=\gamma/(1+\gamma)=1+a$.
Also
\[
 N_{L/E}(1+\gamma)=\gamma+\gamma^{-1}=\tau,
 \qquad a^B=\tau^{-1},\qquad a^{t^2-1}=\gamma.
\]
Since $N_{E/K}(\tau)=1$, the norm of $a$ to $K$ is one, so
$\operatorname{ord}(a)$ divides $d$. The two displayed powers
force both $A$ and $B$ to divide this order. Their gcd is one,
and hence the order is $d$. The conjugate $1+a$ has the same
order. The formula for $a$ is injective in $\gamma$.

\smallskip\noindent\textit{The curve and its four boundary points.}\par
We analyze the norm fiber $X_1$ directly, rather than applying
Proposition~\ref{prop:inner-mixed-character}, whose good-fiber
hypothesis excludes this parameter. Write its variable as $\gamma$.
Over $\overline K$, the relative norm-one torus has coordinates
$(x,y,x^{-1},y^{-1})$. Its $K$-points have
$x=\gamma$, $y=\gamma^t$, $\gamma\in V$, and its Frobenius is
\[
 F_Y(x,y)=(y^{-t},x^t).
\]
The condition $N_{E/K}(\gamma+\gamma^{-1})=1$ gives the
projective closure $Y_{\overline K}\subset\mathbb P^1\times\mathbb P^1$ of
\begin{equation}\label{eq:inner-diagonal-curve}
 (x+1)^2(y+1)^2=xy.
\end{equation}
Here $Y$ denotes the $K$-curve with the inherited torus descent
$F_Y$, not the split curve equipped with coordinatewise
$t$-power Frobenius. The descent extends to the projective
closure and preserves its boundary. In particular the relation
$y=x^t$ is imposed on $K$-points, not in its geometric function field.
On the affine chart the two partial derivatives are $y$ and $x$;
their only common zero $(0,0)$ is not on the curve. The boundary is
\[
 P_0=(0,1),\quad P_\infty=(\infty,1),\quad
 Q_0=(1,0),\quad Q_\infty=(1,\infty).
\]
At $P_0$, with $u=y+1$, the equation is
$(1+x)^2u^2+x(1+u)=0$, whose derivative with respect to $x$
is one at the point. Thus it is smooth and $v_{P_0}(x)=2$.
Reciprocating $x$ or $y$, and interchanging them, gives the
same smoothness test at the other boundary points.
This proves that the bidegree-$(2,2)$ curve is smooth.
It is geometrically irreducible. Smoothness implies that
the curve is reduced and that distinct irreducible components
would be disjoint. Every component has a nonzero effective
bidegree $(a_i,b_i)$, with $a_i,b_i\ge0$ and
$\sum_i(a_i,b_i)=(2,2)$. If one component has both entries
positive, its intersection $a_ib_j+a_jb_i$ with every other
component is positive. If no component has both entries
positive, the sum $(2,2)$ forces a pure vertical component
and a pure horizontal component, whose intersection is again
positive. Both possibilities contradict disjointness.
Adjunction on $\mathbb P^1\times\mathbb P^1$ now gives genus
$(2-1)(2-1)=1$.
Frobenius permutes the four boundary points in the cycle
\[
 P_0\longmapsto Q_0\longmapsto P_\infty
      \longmapsto Q_\infty\longmapsto P_0.
\]
There is no $K$-rational boundary point. Therefore the open curve
$Y^\circ$ has exactly the points in the sum
\eqref{eq:inner-diagonal-character}, and
$|\#Y^\circ(K)-(t+1)|\le2\sqrt t$ by Hasse--Weil.

\smallskip\noindent\textit{Actual Lang covers and their character traces.}\par
Let $\mathcal U=\ker N_{E/K}$ and
$\mathcal V=\operatorname{Res}_{E/K}(\ker N_{L/E})$;
their $K$-point groups are $U$ and $V$. On $Y^\circ$ use the
maps $\tau=\gamma+\gamma^{-1}$ to $\mathcal U$ and $\gamma$
to $\mathcal V$ itself. Write
\[
 w_0=(x+1)^2/x,\qquad h=y^t/x.
\]
The split coordinates of $\tau$ are $(w_0,w_0^{-1})$.
The two Lang pullbacks, with their inherited $K$-structures, are
\begin{equation}\label{eq:inner-diagonal-lang}
 z_U^A=w_0^{-1},\qquad z_V^B=h.
\end{equation}
For the first, the Lang equation is $z_U^{-t}/z_U=w_0$.
For the second, the split equations are
$z_1^{-t}/z_0=x$, $z_0^t/z_1=y$; elimination gives
$z_0^B=y^t/x$, with $z_V=z_0$.
Thus the actual functions are $w_0^{-1}$ and $h$, not merely
a geometrically convenient Kummer model.
Above $F_Y$, Frobenius on the two covering coordinates is
\[
 (z_U,z_V)\longmapsto(w_0z_U,xz_V).
\]
This follows from the split Lang equations. It also respects
\eqref{eq:inner-diagonal-lang}, since
$w_0(F_Y)=w_0^{-t}$ and $h(F_Y)=x^Bh$.
At a $K$-point the deck multipliers are exactly
$\tau$ and $\gamma$. Choosing the character convention for
these multipliers gives trace $\eta(\tau)\chi(\gamma)$;
the inverse convention just inverts both characters.

\smallskip\noindent\textit{Geometric connectedness and conductor.}\par
The complete boundary valuation table is
\begin{equation}\label{eq:inner-diagonal-valuations}
\begin{array}{c|rr}
 &w_0&h\\\hline
 P_0&-2&-2\\
 P_\infty&-2&2\\
 Q_0&2&2t\\
 Q_\infty&2&-2t
\end{array}
\end{equation}
For example $x+1$ has simple zeros at $Q_0,Q_\infty$ and
a double pole at $P_\infty$; $y$ has a double zero at $Q_0$
and a double pole at $Q_\infty$. These observations give every
entry of the table.
Suppose for $0\le e_0<A$, $0\le e_1<B$ that
\[
 (w_0^{-1})^{Be_0}h^{Ae_1}
       \in\overline K(Y)^{\times AB}.
\]
Its valuation at $P_0$ is $2Be_0-2Ae_1$.
Reduction modulo $A$ and modulo $B$, using that $A,B$ are
odd and coprime, gives $A\mid e_0$ and $B\mid e_1$.
This single valuation separates the two cyclic components
because reduction modulo each order kills the other component.
Thus $e_0=e_1=0$. The Kummer criterion proves that the
product cover in \eqref{eq:inner-diagonal-lang} is geometrically
connected with full group $U\times V$; in particular every
nontrivial character quotient is geometrically nontrivial.

All cover degrees are odd, so the ramification is tame.
At every boundary point the valuation of $w_0^{-1}$ is one of
$\pm2$, and that of $h$ is one of $\pm2$ or $\pm2t$.
Since $A$ and $B$ are odd, $\gcd(2,A)=\gcd(2t,B)=1$; each
coordinate therefore generates its cyclic deck factor. A nontrivial
character on either factor is nontrivial on that coordinate, and the
two character values cannot cancel because their orders divide the
coprime integers $A$ and $B$. Thus every nontrivial pair is ramified
at each of the four boundary points, with Swan conductor zero and
local Artin conductor one.
Its total conductor is exactly four, and it is unramified on
$Y^\circ$.
The associated rank-one finite-monodromy sheaf is pure of weight
zero and geometrically nontrivial. Thus $H_c^0=H_c^2=0$, and
Grothendieck--Ogg--Shafarevich gives
$\dim H_c^1=2g(Y)-2+4=4$.
The classical curve character bound
\cite[Lemma~1]{FFSTV13}, with the Lang construction as in
\cite{Katz89}, bounds each Frobenius eigenvalue by $\sqrt t$.
This proves \eqref{eq:inner-diagonal-character}.

\smallskip\noindent\textit{The two primitive indicators and positivity.}\par
Apply the primitive-character indicator to $\tau\in U$ and
$\gamma\in V$. Their product has constant coefficient
$\theta(A)\theta(B)=\theta(d)$ and total absolute coefficient
sum $\theta(d)W(A)W(B)=\theta(d)W(d)$.
The trivial pair has Hasse error $2\theta(d)\sqrt t$.
The other pairs have total coefficient mass
$\theta(d)(W(d)-1)$ and character bound $4\sqrt t$.
Their sum is exactly the error in \eqref{eq:inner-diagonal-count}.
If $\sqrt t\ge4W(d)$, then
$(t+1)-(4W(d)-2)\sqrt t>0$, proving positivity.
For $33\le s\le113$, the complete integer certificates in
Appendix~\ref{app:inner-cert} give
$\omega(d)\le r_s$ with $s\ge4+2r_s$, hence this condition.
For $s\ge114$, the same appendix proves for every odd $n$ that
$W(n)<4096n^{1/8}$. Since $d<2t^3$,
\[
 (4W(d))^8<2^{113}t^3<t^4=(\sqrt t)^8.
\]
These two ranges prove $\Delta_t>0$ for every $s\ge33$.
\end{proof}

The construction \eqref{eq:inner-diagonal-construction} is the
classical trace-one reciprocal-trace change of variable
\cite[\S4]{Coh}, applied here with the additional condition
$N_{E/K}(\tau)=1$ and the smaller exact order $A=t+1$.
The estimate established above concerns this specified fiber;
the Frobenius identity itself is not a priority claim.
It also implies
$\operatorname{Tr}_{L/K}(a)=(a+a^{t^2})+(a^t+a^{t^3})=1+1=0$,
so these particular slopes are not normal over $K$.

The three populations $|\mathcal R_d|$, $N_d^{\mathrm{prim}}$
and $D_d$ have different main terms. Theorem~\ref{thm:inner-double-count}
now gives $D_d>0$ for every $s\ge2$, using the trace-one subfamily
of Proposition~\ref{prop:inner-diagonal-double} for $s\ge33$
and the exact double-order certificates for $2\le s\le32$.
This subfamily has a linear main term $\theta(d)(t+1)$; it does
not replace the quadratic full-family main term
$\theta(d)^2(t+1)^2$. Nor does all-dimensional existence assert
positivity in every primitive norm fiber at small $s$.
Indeed $\Delta_t=0$ at $s=3,5$; the finite certificates at these
parameters use slopes outside this trace-one subfamily.


\begin{theorem}[The six-dimensional Singer application]\label{thm:agld}
$\langle T,D\rangle=\mathrm{AGL}(1,64)$; its order is $4032=64\cdot63$.
\end{theorem}
\begin{proof}
By Proposition~\ref{prop:singer}, in field coordinates the linear part of
$T$ is multiplication by a primitive $\alpha\in\mathbb F_{64}$.
The translation is $c=T(0)=63\ne0$. Thus $U=V$ and the least common multiple
in Corollary~\ref{cor:affine-actions} is already $63$ because
$\operatorname{ord}(\alpha)=63$. The corollary proves the claim.
\end{proof}

The general structure allows reducible and nonsemisimple matrices.
For example, the four-dimensional finite check gives the following distribution:
\[
\begin{array}{c|rrrrrr}
\text{group order}&3&12&15&48&192&240\\
\hline
\text{number of maps}&112&3360&4032&5040&20160&60480
\end{array}
\]
Of the $93184$ affine complete mappings, $80640$
give transitive groups; $60480$ give primitive groups, all doubly transitive
in this dimension. These counts are computational evaluations of the
structure theorem, reproduced by the additional verifier in Appendix~\ref{app:code};
the action criteria themselves are proved above for every dimension.

\section[Uniform statement for 2k-bit words]{Uniform statement for $2k$-bit words}\label{sec:2k}
Identify the $2k$-dimensional $\mathbb F_2$-space $V_{2k}$ with
$\mathbb{F}_{2^{2k}}$; write $M=2^{2k}-1=(2^k-1)(2^k+1)$. In field
coordinates an affine Singer layer is a map $T(x)=\alpha x\oplus b$ whose linear part
is multiplication by a primitive element $\alpha$ (order $M$).
\begin{lemma}\label{lem:ord}
For every $k\ge1$, $\operatorname{ord}_{2^k+1}(2)=2k$.
\end{lemma}
\begin{proof}
Let $m=\operatorname{ord}_{2^k+1}(2)$. Since $2^k\equiv-1\pmod{2^k+1}$,
we have $m\mid2k$, while $m\nmid k$. Suppose, for contradiction, that
$m<2k$. If $m$ were odd, $m\mid2k$ would imply $m\mid k$.
Thus $m$ is even; write $m=2d$
with $d\mid k$, and $m<2k$ gives $d<k$. The subgroup
$\langle2\rangle\subset(\mathbb Z/(2^k+1)\mathbb Z)^\times$ is cyclic of order
$2d$, so its unique element of order $2$ is $2^d$. But $2^k\equiv-1$ also has
order $2$ in this subgroup; hence $2^d\equiv-1\pmod{2^k+1}$. Thus
$2^k+1\mid 2^d+1$, contradicting
$0<2^d+1<2^k+1$. Therefore $m=2k$.
\end{proof}
For \(q=2^k\), define
\[
\operatorname{PTE}(k):\quad
\exists\beta\in\mathbb{F}_{q^2}^*,\quad
\operatorname{ord}(\beta)=q^2-1,\quad
\operatorname{Tr}_{\mathbb{F}_{q^2}/\mathbb{F}_q}(\beta)=1.
\]
The following theorem supplies the required existence input; its proof is given in Appendix~\ref{app:pte}.
\begin{samepage}
\begin{theorem}[Classical PTE existence]\label{thm:pte}
For every \(k\ge1\), \(\operatorname{PTE}(k)\) holds.
\end{theorem}
\begin{proof}
See Appendix~\ref{app:pte}.
\end{proof}
\end{samepage}
\begin{theorem}\label{thm:2k}
For every $k\ge1$ there exists a primitive $\alpha\in\mathbb F_{2^{2k}}$
such that, for every $b\in\mathbb F_{2^{2k}}^*$, the affine Singer layer
$T_b(x)=\alpha x\oplus b$ has
$D_b=(\mathrm{id}\oplus T_b)^{-1}$ and $K_b=(\mathrm{id}\oplus D_b)^{-1}$ defined.
The maps $T_b,D_b$ have cycle type $1+M$, while $K_b$ has cycle type
$1+(2^k-1)\times(2^k+1)$, and $\langle T_b,D_b\rangle=\mathrm{AGL}(1,2^{2k})$.
The triality operator satisfies $F^3=\mathrm{id}$ with
$T_b\mapsto D_b\mapsto K_b\mapsto T_b$.
\end{theorem}
\begin{proof}
Let $b\ne0$ be arbitrary and write $T=T_b$, $D=D_b$, $K=K_b$.
The linear parts of $T$, $D$, $K$ are
$\alpha$, $(1\oplus\alpha)^{-1}$, and
$\kappa=1\oplus\alpha^{-1}=(1\oplus\alpha)/\alpha$
respectively (independent of the translation $b$); the three slopes are cycled by
the map $z\mapsto 1/(z\oplus1)$, so $F^3=\mathrm{id}$. Since
$\mathbb{F}_{2^{2k}}^*$ is cyclic, an affine map with slope $\mu\ne1$ has cycle type
$1+\frac{M}{\mathrm{ord}(\mu)}\times\mathrm{ord}(\mu)$ (the translation only moves the
fixed point; a shorter orbit would force $T^d$, $d<\mathrm{ord}(\mu)$, to have a second
fixed point, impossible since $\mu^d\ne1$); the stated types therefore follow once
$\mathrm{ord}(\alpha)=\mathrm{ord}(1\oplus\alpha)=M$ and
$\mathrm{ord}(\kappa)=2^k+1$. Existence of such an $\alpha$ follows directly from Theorem~\ref{thm:pte}.
Put $q=2^k$ and choose a primitive $\beta\in\mathbb F_{q^2}$ with
$\beta^q\oplus\beta=1$. Take $\alpha=\beta$. Then $1\oplus\alpha=\alpha^q$,
so
\[
\mathrm{ord}(\alpha)=\mathrm{ord}(1\oplus\alpha)=q^2-1=M,
\qquad
\kappa=\frac{1\oplus\alpha}{\alpha}=\alpha^{q-1},
\qquad \mathrm{ord}(\kappa)=q+1.
\]
Here the last equality follows from
$\mathrm{ord}(\alpha^{q-1})=M/\gcd(M,q-1)=q+1$.
These are exactly the required slope orders. To make the dependence on $b$
explicit, set $S_b(x)=bx$. Then $T_b=S_bT_1S_b^{-1}$.
Indeed, for any invertible additive map $S$ and $g=SfS^{-1}$,
\[
(\mathrm{id}\oplus g)(x)=S\bigl((\mathrm{id}\oplus f)(S^{-1}x)\bigr).
\]
Taking inverses gives $F(g)=SF(f)S^{-1}$ whenever $F(f)$ is defined.
Applying this with $S=S_b$ gives $D_b=S_bD_1S_b^{-1}$ and
$K_b=S_bK_1S_b^{-1}$. Their fixed points are
\[
\mathrm{Fix}(T_b)=\{b/(1\oplus\alpha)\},\qquad
\mathrm{Fix}(D_b)=\{b/\alpha\},\qquad
\mathrm{Fix}(K_b)=\{b\}.
\]
Thus all nonzero $b$ give the stated cycle types; \S\ref{sec:group} gives
the group-generation claim.
\end{proof}
Since the $K$-slope has order $2^k+1$, the degree over $\mathbb{F}_2$ of its
minimal polynomial is $\operatorname{ord}_{2^k+1}(2)=2k$ by
Lemma~\ref{lem:ord}; this polynomial is therefore an irreducible degree-$2k$
factor of $\Phi_{2^k+1}$.
\textbf{Remark (the $\tau$-collapse).}
Put $q=2^k$. For $\gamma\in\mathbb F_{q^2}^*$ of order $q+1$,
one has $\gamma^q=\gamma^{-1}$; hence
$\tau=\gamma\oplus\gamma^{-1}$ satisfies $\tau^q=\tau$ and belongs to
$\mathbb F_q$. The Frobenius identities
$(1\oplus\gamma)^{2^k-1}=\gamma^{-1}$ and $(1\oplus\gamma)^{2^k+1}=\tau$,
together with $\gcd(2^k-1,2^k+1)=1$, give the following order calculation.
For an element $z$ of finite order $N$ and coprime positive integers $r,s$,
$\operatorname{ord}(z^r)=N/\gcd(N,r)$ implies
$N=\operatorname{lcm}(\operatorname{ord}(z^r),\operatorname{ord}(z^s))$.
Consequently
$\mathrm{ord}(1\oplus\gamma)=\mathrm{lcm}(\mathrm{ord}(\gamma^{-1}),\mathrm{ord}(\tau))
=\mathrm{lcm}(2^k+1,2^k-1)=M$ whenever $\tau$ is primitive in $\mathbb{F}_{2^k}$;
the pair $(u^{2^k-1},u^{2^k+1})=(\gamma^{-1},\tau^{-1})$ for
$u=\gamma/(1\oplus\gamma)$ gives the same order for $u$. Since
$\alpha=(1\oplus\gamma)^{-1}$ and
$\kappa=1\oplus\alpha^{-1}=\gamma$, once $\gamma$
has order $2^k+1$ the three order conditions on $\alpha$ reduce to the single
condition that $\tau$ be primitive. The order assumption on $\gamma$ is
essential; it is not implied by $\tau$ being primitive alone.
Existence of such a $\gamma$ is the classical square-order statement
(see \cite{Mor} and \cite[Theorem~1, p.~1]{Coh90}; for an equivalent square-order formulation, see
\cite[Theorem~1.5, p.~170]{Coh}); it is also supplied directly by
Theorem~\ref{thm:pte}.
For $q=2^k$ the appendix proof is self-contained, so the paper does not use these
classical results as external existence inputs.

The theorem is an existence statement in admissible coordinates (equivalently, for an
admissible ordered generator basis). The argument of \S\ref{sec:group} applies verbatim
with $Q=2^{2k}$ (taking $b\ne0$ in the proof), so the generation
$\langle T,D\rangle=\mathrm{AGL}(1,2^{2k})$ holds for every $k$. For the concrete
coordinate presentation of \S\ref{sec:T} the orders and cycle data were verified
for $k=1,2,3,4$ (the bases are listed in the appendix):
$k=1$ (2 bits) order $12=4\cdot3$, $K$ type $1+1\cdot3$, factor
$\Phi_3$ (basis $(1,3)$: $\mathrm{Fix}(T)=\{1\}$, $\mathrm{Fix}(D)=\{2\}$,
$\mathrm{Fix}(K)=\{3\}$); $k=2$ (4 bits) order $240=16\cdot15$, $K$ type
$1+3\cdot5$, factor $\Phi_5$; $k=3$ (6 bits) order $4032=64\cdot63$, $K$ type
$1+7\cdot9$, factor $\Phi_9$ (this paper); $k=4$ (8 bits) order
$65280=256\cdot255$, $K$ type $1+15\cdot17$, linear part an irreducible degree-$8$
factor of $\Phi_{17}$.

The basis is not canonical even in dimension $6$: among the $720$ orderings of these
six generators, $572$ do not even make $\mathrm{id}\oplus T$ a permutation; of the $148$
that do, only $54$ give $T$ of order $63$, and only $9$ give the exact cycle-type triple
$T=D=1+63$, $K=1+7\cdot9$. For instance $(4\,6)$ keeps $T=1+63$ but changes
$\mathrm{Fix}(T)$ from $38$ to $11$ and the cycle types of $D$ and $K$ to $D=1+3\cdot21$,
$K=1+63$. Hence the particular entries of $G$ and the anchors $\{38,41,63\}$ are coordinates
of one representative, not invariants, and no uniqueness statement is made. The structural results of \S\ref{sec:group}--\ref{sec:2k},
however, are representative-independent: $\langle T,D\rangle=\mathrm{AGL}(1,64)$, the
third-iterate identity $F^3=\mathrm{id}$, and the $2k$-bit existence statement hold for any
admissible basis; only the numerical anchors change with the basis.

\begingroup
\emergencystretch=1em
\noindent\textbf{Remark (constructiveness).}\\
Theorem~\ref{thm:2k} is an existence statement; its existence input is proved
in Appendix~\ref{app:pte}. The $\tau$-collapse turns it into a finite search with
certification. For $\tau=\gamma\oplus\gamma^{-1}\in\mathbb{F}_{2^k}$, the
Artin--Schreier identity
(\emph{$x^2+\tau x+1$ is irreducible over $\mathbb{F}_{2^k}$ iff
$\mathrm{Tr}_{\mathbb{F}_{2^k}/\mathbb{F}_2}(1/\tau)=1$}) gives $2^{k-1}$
candidate values of $\tau$. For a candidate, fix an explicit representation of
$\mathbb{F}_{2^{2k}}$ and take a root $\gamma$ of $x^2+\tau x+1$
(found, for instance, by exhaustive search); the collapse gives the required orders for $1\oplus\gamma$ and
$\gamma/(1\oplus\gamma)$ precisely when $\tau$ is primitive \emph{and}
$\mathrm{ord}(\gamma)=2^k+1$. The second condition is not automatic from the
first: for example, when $k=6$ and
$\mathbb{F}_{64}=\mathbb{F}_2[t]/(t^6+t+1)$, the element
$\tau=t^5+t^3+1$ is primitive and $\mathrm{Tr}(1/\tau)=1$, but the roots of
$x^2+\tau x+1$ have order $13$, not $65$. The theorem of
Appendix~\ref{app:pte} guarantees that a jointly successful candidate exists, so a
deterministic enumeration that tests both conditions terminates. Certifying them requires factoring $2^k-1$ and
$2^k+1$ (or an equivalent exact-order test). This procedure supplies a
deterministic construction with certification; no factorization-free closed
formula is supplied here.
A necessary restatement is that $u=1/\tau$ is a primitive element of
$\mathbb{F}_{2^k}$ with $\mathrm{Tr}_{\mathbb{F}_{2^k}/\mathbb{F}_2}(u)=1$;
equivalently, its degree-$k$ minimal polynomial has $x^{k-1}$-coefficient
$1$. This condition alone is still not sufficient for admissibility, because
the root-order condition above must also be checked.
\endgroup

\section[Fixed mappings and orbits under F]{Fixed mappings and orbits under $F$}\label{sec:rem}
Lemma~\ref{lem:F3} implies that every orbit under $F$ has length one or three.
We first give the scalar classification, then describe all fixed complete
mappings geometrically and count the orbits in the entire affine class.
An \emph{$F$-fixed map} means $F(f)=f$; this is distinct from the unique
fixed point of the permutation $f$ in Lemma~\ref{lem:fp}.

\subsection*{Scalar affine maps}
\begin{proposition}\label{prop:F-orbits}
Let $k\ge1$, $Q=2^{2k}$, and
$\mathcal A_Q=\{f_{a,c}:x\mapsto ax\oplus c\mid a\in\mathbb F_Q\setminus\{0,1\},\ c\in\mathbb F_Q\}$.
The operator $F$ acts on $\mathcal A_Q$ with two singleton orbits and
$\bigl(Q(Q-2)-2\bigr)/3$ orbits of length three. Its fixed maps are exactly
$f_{a,0}$ with $a^2+a+1=0$. On the invariant subset $c\ne0$, every orbit
has length three.
\end{proposition}
\begin{proof}
In coefficient pairs,
\[
F(a,c)=\left((1+a)^{-1},\frac{c}{1+a}\right).
\]
The new slope again lies outside $\{0,1\}$, and Lemma~\ref{lem:F3} gives
$F^3=\mathrm{id}$, so each orbit has length one or three. A fixed pair
must satisfy $a=(1+a)^{-1}$, equivalently $a^2+a+1=0$; the constant-term
equation is $c(1+a)=c$, so $ac=0$ and $c=0$.
Conversely these conditions give a fixed pair. The polynomial $a^2+a+1$
has two distinct roots in $\mathbb F_Q$, since $\mathbb F_4\subseteq\mathbb F_Q$.
There are $Q(Q-2)$ pairs in total, yielding the stated counts.
Multiplication of $c$ by $(1+a)^{-1}\ne0$ preserves $c\ne0$, and no fixed
pair lies in that subset.
\end{proof}
\subsection*{Fixed maps and line spreads}
A \emph{line spread} of $V=\mathbb F_2^n$ is a collection of two-dimensional
$\mathbb F_2$-subspaces whose nonzero vectors partition $V\setminus\{0\}$.
Thus a line here is two-dimensional as a vector space; it is a projective
line in $\mathrm{PG}(n-1,2)$.
Each such subspace has exactly three nonzero vectors $x,y,x\oplus y$.
A cyclic orientation on a line specifies one of the two $3$-cycles on them.

\begin{proposition}\label{prop:F-spreads}
For $n\ge1$, the $F$-fixed complete mappings of $\mathbb F_2^n$ are in
bijection with line spreads equipped with a cyclic orientation on each line.
Every such map fixes $0$ and has only $3$-cycles on the nonzero vectors.
They exist exactly when $n$ is even. If $s_n$ denotes the number of line spreads
in a fixed coordinate space, their number, for even $n$, is
\[
s_n\,2^{(2^n-1)/3}.
\]
Nonlinear $F$-fixed complete mappings exist for every even $n\ge4$.
\end{proposition}
\begin{proof}
If $F(f)=f$, then $f^{-1}=\mathrm{id}\oplus f$. Evaluating at $f(x)$ gives
\[
f^2(x)=x\oplus f(x)=f^{-1}(x),\qquad f^3=\mathrm{id}.
\]
Let $p$ be the unique fixed point from Lemma~\ref{lem:fp}.
The displayed equality at $p$ gives $p=p\oplus p=0$.
Thus each nonzero cycle is $(x,f(x),x\oplus f(x))$, the nonzero part of
a two-dimensional subspace. Disjoint cycles supply a line spread and
their orientations. Conversely, orient each line of a spread and fix $0$.
The resulting permutation satisfies $f^2=\mathrm{id}\oplus f=f^{-1}$,
so it is complete and $F(f)=f$. This proves the bijection and the count.

A spread requires $3\mid2^n-1$, hence even $n$. If $n=2m$, identify
$V$ with $\mathbb F_4^m$; its one-dimensional $\mathbb F_4$-subspaces
give a line spread over $\mathbb F_2$, proving existence.
For $m\ge2$ (equivalently, even $n\ge4$), start with the linear map
$A:x\mapsto\omega x$, where
$\omega^2+\omega+1=0$ in $\mathbb F_4$, and reverse the orientation
on exactly one spread line $U$. The resulting map $f$ equals $A^2$ on
$U\setminus\{0\}$ and equals $A$ elsewhere, and remains $F$-fixed.
For $0\ne u\in U$ and $v\notin U$, one has $u\oplus v\notin U$, and
\[
f(u\oplus v)=Au\oplus Av\ne A^2u\oplus Av=f(u)\oplus f(v),
\]
since $(A^2+A)u=u\ne0$. Hence $f$ is nonlinear.
\end{proof}

If $V$ is identified with $\mathbb F_{2^n}$ for even $n$, the spread obtained
from $\mathbb F_4$ consists of the multiplicative cosets of $\mathbb F_4^*$.
An orientation on each coset is multiplication by either $\omega$ or $\omega^2$.
Thus these examples are a particular family of the first-order cyclotomic
complete mappings considered in \cite{BW21}. Proposition~\ref{prop:F-spreads}
describes fixed maps from arbitrary line spreads, without choosing a field
structure or restricting to this coset partition.

For $n=4$, exact-cover enumeration of the $35$ lines on $15$ nonzero vectors
gives $s_4=56$, recovering the classical spread count also formalized in
Coq by Magaud \cite{Mag22}. The proposition then gives $56\cdot2^5=1792$ fixed
complete mappings. Enumeration of their orientations finds $112$ linear
and $1680$ nonlinear maps, in agreement with the affine count below.
These finite counts, the complete list of spreads, and a nonlinear witness
are reproduced by \path{verify_complete_mapping_orbits.py}.

\subsection*{All affine complete mappings}
Write $h_j=|\mathrm{GL}(j,2)|=\prod_{i=0}^{j-1}(2^j-2^i)$,
where $h_0=1$ is the empty product,
and let
\[
\mathcal C_n^{\mathrm{aff}}
=\{x\mapsto Lx\oplus c:\ L,I+L\in\mathrm{GL}(n,2),\ c\in\mathbb F_2^n\}.
\]
This includes matrices that are not scalar multiplication in a single
field structure on $\mathbb F_2^n$.

\begin{proposition}\label{prop:F-affine}
For $n\ge1$, put
\[
d_n=h_n\sum_{j=0}^n\frac{(-1)^j2^{j(j-1)/2}}{h_j},\qquad
r_n=\begin{cases}
0,&n\text{ odd},\\
h_n/|\mathrm{GL}(n/2,4)|,&n\text{ even}.
\end{cases}
\]
There are $2^n d_n$ affine complete mappings.
Under $F$ they form $r_n$ singleton orbits and
$(2^n d_n-r_n)/3$ orbits of length three.
The fixed maps are exactly $x\mapsto Lx$ with $L^2+L+I=0$.
\end{proposition}
\begin{proof}
The number $d_n$ counts invertible matrices with no nonzero fixed vector.
Its displayed expression is the zero-dimensional fixed-space case of the
classical Rudvalis--Shinoda count, recorded in \cite[Theorem~2.1]{FS15}.
Here is a direct proof. For a $j$-dimensional subspace $U$ put
$\mu(U)=(-1)^j2^{j(j-1)/2}$. If $W$ has dimension $r$, the
Gaussian-binomial identity gives
\[
\sum_{U\le W}\mu(U)
=\sum_{j=0}^r\binom rj_2(-1)^j2^{j(j-1)/2}
=\prod_{i=0}^{r-1}(1-2^i).
\]
This is $1$ when $W=0$ and $0$ otherwise. Apply it to
$W=\ker(I+L)$ and sum over $L\in\mathrm{GL}(n,2)$.
For a fixed $j$-subspace $U$, the matrices fixing $U$ pointwise have block
form $\left(\begin{smallmatrix}I&B\\0&C\end{smallmatrix}\right)$,
with $2^{j(n-j)}$ choices of $B$ and $h_{n-j}$ choices of $C$.
Since there are $\binom nj_2$ choices of $U$ and
\[
\binom nj_2\,2^{j(n-j)}h_{n-j}=h_n/h_j,
\]
interchanging the two finite sums yields the formula for $d_n$.
Each such $L$ permits all $2^n$ translations $c$.

For $f(x)=Lx\oplus c$, the condition $F(f)=f$ is equivalent to
$(\mathrm{id}\oplus f)\circ f=\mathrm{id}$. Comparing linear and
constant parts gives $L^2+L=I$ and $Lc=0$, so $c=0$.
The polynomial $z^2+z+1$ is irreducible over $\mathbb F_2$.
Thus such an $L$ equips $V$ with an $\mathbb F_4$-vector-space structure,
which requires even $n$. Conversely, for even $n$ multiplication by a
root of this polynomial supplies such an $L$.
All these matrices form one conjugacy class in $\mathrm{GL}(n,2)$:
choose an $\mathbb F_4$-basis in each structure and compare the resulting
$\mathbb F_2$-bases. The invertible maps commuting with $L$ are exactly
the $\mathbb F_4$-linear ones, so the centralizer has order
$|\mathrm{GL}(n/2,4)|$. This proves the formula for $r_n$.
Finally Lemma~\ref{lem:F3} restricts the other orbits to length three.
\end{proof}

The following small values are exact evaluations of the formulas, independently
checked by enumerating all binary matrices and translations in these dimensions:
\begin{center}
\begin{tabular}{r|r|r|r|r}
$n$ & $d_n$ & $|\mathcal C_n^{\mathrm{aff}}|$ & $F$-fixed maps & $3$-element orbits\\
\hline
2 & 2 & 8 & 2 & 2\\
3 & 48 & 384 & 0 & 128\\
4 & 5824 & 93184 & 112 & 31024\\
\end{tabular}
\end{center}
Every $F$-fixed affine map is linear. In particular the $1680$ nonlinear fixed
maps in dimension four lie outside the affine class, showing the additional
scope of Proposition~\ref{prop:F-spreads}.

\subsection*{Linear conjugacy of oriented spreads}
We now distinguish the $F$-orbits of maps from equivalence under a linear
change of coordinates. Write $f\sim g$ if $g=SfS^{-1}$ for some
$S\in\mathrm{GL}(n,2)$. Such conjugacy preserves $F(f)=f$.
For a spread $\mathcal S$, let $G_{\mathcal S}$ be its setwise stabilizer
in $\mathrm{GL}(n,2)$, and choose a reference cyclic orientation on each line.
An orientation is then a vector $(\epsilon_L)_{L\in\mathcal S}$ over $\mathbb F_2$,
with $1$ meaning reversal of the reference orientation.
For $g\in G_{\mathcal S}$ define $\delta_g(L)\in\mathbb F_2$ by comparing the
transported reference orientation on $gL$ with the chosen reference on $gL$.
Thus the transported bits satisfy
\[
\epsilon'_{gL}=\epsilon_L\oplus\delta_g(L).
\]

\begin{proposition}[Orientation-sensitive orbit formula]\label{prop:oriented-Burnside}
The correspondence of Proposition~\ref{prop:F-spreads} identifies linear
conjugacy with equivalence of oriented spreads. Let $c(g)$ be the number of
cycles in the permutation of spread lines induced by $g$, and put
\[
e(g)=\begin{cases}
1,&\displaystyle\bigoplus_{L\in C}\delta_g(L)=0
       \text{ for every line cycle }C,\\
0,&\text{otherwise}.
\end{cases}
\]
The number of linear-conjugacy classes supported on the orbit of $\mathcal S$ is
\[
a(\mathcal S)=\frac{1}{|G_{\mathcal S}|}
      \sum_{g\in G_{\mathcal S}}e(g)2^{c(g)}.
\]
Summing this over representatives of the linear equivalence classes of
underlying spreads gives all $F$-fixed linear-conjugacy classes.
The stabilizer of a fixed map is the stabilizer of its oriented spread.
\end{proposition}
\begin{proof}
Linear conjugacy sends each cycle $(x,f(x),x\oplus f(x))$ to its transported
cycle, so it sends the spread and its orientations to exactly those of
$SfS^{-1}$. To fix an orientation, the displayed bit relation must hold
with $\epsilon'=\epsilon$. Around a line cycle the successive equations
are consistent if and only if the sum of the $\delta_g(L)$ on that cycle
is zero. When consistent, there are two choices for its first bit and all
other bits are forced. Thus $g$ fixes $e(g)2^{c(g)}$ orientations.
Burnside's orbit-counting lemma gives the formula. Restricting to the
stabilizer of one underlying spread suffices for its full linear orbit;
distinct underlying spread orbits cannot give conjugate maps.
The same transport correspondence gives the stabilizer assertion.
\end{proof}

\begin{proposition}[Local criterion for an invariant orientation]
\label{prop:orientation-local}
Let $\mathcal S$ be any line spread and
$G=G_{\mathcal S}$ its setwise stabilizer in $\mathrm{GL}(V)$.
For $L\in\mathcal S$, write $G_L=\{g\in G:gL=L\}$.
The following conditions are equivalent:
\begin{enumerate}
\item There is an orientation of every line in $\mathcal S$
that is fixed by all of $G$.
\item Every $g\in G_L$, for every $L\in\mathcal S$,
induces an even permutation of the three nonzero vectors of $L$.
\item The orientation-translation cocycle associated with any
choice of reference orientations is a coboundary.
\end{enumerate}
More precisely, put $E=\mathbb F_2^{\mathcal S}$ with its natural
permutation $G$-module structure $(g\cdot v)_{gL}=v_L$.
If $D_g\in E$ is defined by $(D_g)_{gL}=\delta_g(L)$,
then $D_{gh}=D_g+g\cdot D_h$, and condition (3) means that
there exists $v\in E$ with $D_g=v+g\cdot v$ for every $g\in G$.
\end{proposition}
\begin{proof}
An orientation of a line is a cyclic ordering of its three
nonzero vectors, up to cyclic rotation. Its stabilizer within
the permutation group on those vectors is the alternating
group of degree three. Hence (1) implies (2).

Assume (2). Choose one representative $L$ from each $G$-orbit
on $\mathcal S$ and choose an arbitrary cyclic orientation on it.
Transport the chosen orientation to $gL$ using $g\in G$.
If $gL=hL$, then $h^{-1}g\in G_L$ and (2) says that this
element preserves the chosen orientation. Thus the transport
is independent of the choice of $g$, and the resulting orientation
on all lines is fixed by $G$. This proves (2)$\Rightarrow$(1).

For the last equivalence, the transported bit vector is
$g\cdot v+D_g$. Composing two transports gives
$D_{gh}=D_g+g\cdot D_h$. The bit vector $v$ is fixed by
every $g$ precisely when $v=g\cdot v+D_g$, or equivalently
$D_g=v+g\cdot v$. This is exactly the coboundary condition.
\end{proof}

The proposition distinguishes the existence of one orientation
fixed by the entire spread automorphism group from the individual
fixed-orientation tests $e(g)$ in
Proposition~\ref{prop:oriented-Burnside}; the latter still supply
the general orbit count even when no globally invariant orientation exists.

\begin{theorem}[Four-dimensional classification]\label{thm:fixed-n4-conjugacy}
There are exactly three conjugacy classes under $\mathrm{GL}(4,2)$ of $F$-fixed
complete mappings of $\mathbb F_2^4$. Identify $V$ with $\mathbb F_4^2$, choose
$\omega^2+\omega+1=0$, and start with $f_0(x)=\omega x$ on its five spread lines.
Representatives are obtained by reversing any $r$ lines, with $r=0,1,2$.
Their orbit sizes, stabilizer orders and stabilizer groups are
\begin{center}
\begin{tabular}{c|r|r|c|c}
$r$ & orbit size & stabilizer order & stabilizer group & linear\\
\hline
0 & 112 & 180 & $C_3\times A_5$ & yes\\
1 & 560 & 36 & $C_3\times A_4$ & no\\
2 & 1120 & 18 & $C_3\times S_3$ & no
\end{tabular}
\end{center}
Here $C_3$ is cyclic of order three, and $A_j,S_j$ denote alternating and
symmetric groups. In particular the $1680$ nonlinear maps form two classes.
\end{theorem}
\begin{proof}
First every four-dimensional spread is equivalent to the displayed one.
Choose two spread lines as complementary spaces $E\oplus E$, where
$E=\mathbb F_2^2$. The other three lines are graphs of invertible maps
$P_i:E\to E$, with all $P_i+P_j$ invertible. A change of coordinates makes
one of these maps $I$. Among the six matrices of $\mathrm{GL}(2,2)$, the only
ones for which $I+P$ is invertible are
\[
P=\begin{pmatrix}0&1\\1&1\end{pmatrix},\qquad
P^2=\begin{pmatrix}1&1\\1&0\end{pmatrix}.
\]
They satisfy $P^2+P+I=0$. The three graphs are consequently those of
$I,P,P^2$, giving the five one-dimensional $\mathbb F_4$-subspaces of
$\mathbb F_4^2$.

For this standard spread, the only linear maps inducing cyclic orientations
on its lines are $f_0$ and $f_0^2$. Indeed on each line such a map is
multiplication by $\omega$ or $\omega^2$; additivity at $e_1+e_2$ forces
the choices on the two coordinate lines to agree, which determines the
map on their direct sum. Therefore any linear stabilizer of the underlying
spread conjugates $f_0$ to either $f_0$ or $f_0^2$.
The commuting maps are $\mathrm{GL}(2,4)$; composing with the coordinatewise
Frobenius map $\sigma(x_1,x_2)=(x_1^2,x_2^2)$ accounts for the other coset.
Thus the full spread stabilizer is $\Gamma\mathrm L(2,4)$ of order
$2|\mathrm{GL}(2,4)|=360$. Since $|\mathrm{GL}(4,2)|=20160$, there are
$20160/360=56$ spreads, recovering the classical count \cite{Mag22}.

On the five lines, $\mathrm{GL}(2,4)$ has scalar kernel $C_3$ and image
$\mathrm{PGL}(2,4)$ of order $60$. This index-two subgroup of $S_5$ is $A_5$.
Its elements preserve every reference orientation. Frobenius reverses them
all and induces the transposition exchanging $\omega$ and $\omega^2$ on
$\mathbb P^1(\mathbb F_4)$, fixing $0,1,\infty$.
Consequently even line permutations permute the reversed-line subset,
whereas odd line permutations also complement it.
$A_5$ is transitive on subsets of each fixed size: an arbitrary permutation
carrying one subset to another can have its parity adjusted by a transposition
inside the subset or its complement. Hence the orientation orbits are
exactly the weight pairs $\{0,5\}$, $\{1,4\}$, $\{2,3\}$, of sizes
$2,10,20$. Multiplying by the $56$ underlying spreads gives the three
displayed orbit sizes and orbit--stabilizer gives their stabilizer orders.

For the representatives of weights $r=0,1,2$, an odd line permutation
cannot stabilize the orientation because it changes the weight to $5-r$.
Within $\mathrm{GL}(2,4)$ the projective stabilizers are $A_5$, $A_4$,
and the two-subset stabilizer $S_3$, respectively.
The determinant splits $\mathrm{GL}(2,4)=C_3\times\mathrm{SL}(2,4)$;
the projective action identifies $\mathrm{SL}(2,4)$ with $A_5$, giving
the displayed direct products. The earlier two-linear-map observation
shows that only the first class is linear.
\end{proof}

The proof gives representatives and stabilizers, rather than only the total
number of maps. Independently, \path{verify_affine_groups_and_conjugacy.py}
constructs all $1792$ maps and conjugates the three representatives by all
$20160$ matrices; their disjoint orbits exhaust the list. It also compares
direct fixed-orientation counts with Proposition~\ref{prop:oriented-Burnside}
for all $360$ standard-spread automorphisms.
For dimension six, underlying spreads have already been classified by
Mateva and Topalova \cite{MT09} into $131044$ types. We next treat the
standard Desarguesian type, one of these types. The other underlying types
are outside the six-dimensional computation reported here.


\subsection*{Standard spreads in arbitrary even dimension}
Let $V=\mathbb F_4^m$, let $\omega^2+\omega+1=0$, and let $\mathcal S_m$
consist of its one-dimensional $\mathbb F_4$-subspaces, viewed as
two-dimensional $\mathbb F_2$-subspaces. Put
\[
N_m=\frac{4^m-1}{3},\qquad
P_m=\mathrm{PGL}(m,4),\qquad
p_m=|P_m|=\frac{1}{3}\prod_{i=0}^{m-1}(4^m-4^i).
\]
Orient each line by $x\mapsto\omega x$ before choosing the reversed-line
subset. Write $a_m$ for the number of full linear-conjugacy classes supported
on spreads equivalent to $\mathcal S_m$.

\begin{theorem}[Standard-spread action and asymptotics]\label{thm:standard-asymptotic}
For $m\ge2$, the setwise stabilizer of $\mathcal S_m$ in
$\mathrm{GL}(2m,2)$ is $\Gamma\mathrm L(m,4)$.
Its linear coset permutes orientation bits; its semilinear coset also
complements every bit and fixes no orientation. Every oriented stabilizer
contains the scalar group $C_3=\mathbb F_4^*$ and is contained in
$\mathrm{GL}(m,4)$. If $c(g)$ counts the cycles of $g$ on the $N_m$
projective points, then
\begin{equation}\label{eq:standard-class}
a_m=\frac{1}{2p_m}\sum_{g\in P_m}2^{c(g)}
    =\frac{2^{N_m}}{2p_m}(1+\varepsilon_m),
\qquad
0\le\varepsilon_m\le(p_m-1)2^{-2^{2m-3}}.
\end{equation}
For a uniformly chosen orientation,
\begin{equation}\label{eq:standard-exception}
\Pr\bigl(\operatorname{Stab}_{\mathrm{GL}(2m,2)}(f)\ne C_3\bigr)
\le(p_m-1)2^{-2^{2m-3}}.
\end{equation}
Consequently almost every orientation has stabilizer exactly $C_3$ as
$m\to\infty$, and
\[
a_m\sim\frac{3\,2^{(4^m-1)/3}}{2|\mathrm{GL}(m,4)|}.
\]
Such orientations exist for every $m\ge4$ and are nonlinear.
Furthermore, exactly $a_m-1$ of the $a_m$ supported conjugacy classes
are nonlinear for every $m\ge2$.
\end{theorem}
\begin{proof}
A linear map giving cyclic orientations on the lines of $\mathcal S_m$
must be either $f_0(x)=\omega x$ or $f_0^2(x)=\omega^2x$.
On each line it has one of these two restrictions. Applying additivity
to independent coordinate vectors and their sum forces their choices to
agree, and the two $\mathbb F_2$-basis vectors on each coordinate line
then determine the map on $V$. A spread stabilizer conjugates $f_0$ to
one of these two maps. The centralizer of $f_0$ is $\mathrm{GL}(m,4)$;
coordinatewise Frobenius $\sigma(x)=x^2$ conjugates $f_0$ to $f_0^2$.
This proves the spread-stabilizer assertion.

Linear elements preserve the reference orientations; elements of
$\mathrm{GL}(m,4)\sigma$ reverse all of them. Such an element sends an
orientation of weight $r$ to one of weight $N_m-r$. As $N_m$ is odd, it
fixes none. Scalars act trivially on the orientation set, giving
the first equality in \eqref{eq:standard-class} by
Proposition~\ref{prop:oriented-Burnside}.
The orientation stabilizer is therefore the inverse image in
$\mathrm{GL}(m,4)$ of its subset stabilizer in $P_m$.

For a nonscalar matrix $A$, its fixed projective points are the lines in
its eigenspaces for the three nonzero elements of $\mathbb F_4$.
If their dimensions are $d_\lambda$, the number is
$\sum_\lambda(4^{d_\lambda}-1)/3$.
Their dimensions sum to at most $m$ and each is at most $m-1$.
Convexity gives the upper bound $(4^{m-1}-1)/3+1$.
Thus every nonidentity projective element moves at least $4^{m-1}-1$
points. Each nontrivial permutation cycle has at least two points, so
\[
N_m-c(g)\ge
\left\lceil\frac{4^{m-1}-1}{2}\right\rceil=2^{2m-3}.
\]
It fixes a proportion $2^{c(g)-N_m}$ of all orientations.
Summing these proportions bounds both the nonidentity contribution
in \eqref{eq:standard-class} and, by the union bound, the probability
in \eqref{eq:standard-exception}. This is the classical motion argument
\cite{RS98} in the present projective action.

Finally $p_m<2^{2m^2}/3$ and $2^{2m-3}\ge2m^2$ for $m\ge4$;
the latter holds at $m=4$ and increases thereafter.
The exceptional bound is then less than one and tends to zero.
The only two linear orientations have stabilizer $\mathrm{GL}(m,4)$,
which is larger than $C_3$, so orientations with exactly $C_3$ are nonlinear.
The two uniform orientations are conjugate by the semilinear Frobenius map,
and no other orientation is linear; hence they constitute the single
linear conjugacy class and the other $a_m-1$ classes are nonlinear.
\end{proof}

For example $p_4=987033600$, and the exceptional proportion is at most
$987033599/2^{32}<0.230$. The estimate therefore already guarantees that
more than $77\%$ of the orientations on the standard eight-dimensional
spread have exactly the scalar stabilizer. This conclusion concerns the
standard spread family, without asserting a distribution over all spreads.

\begin{corollary}[Six-dimensional standard-spread classes]\label{cor:standard-n6}
The $2^{21}$ orientations of $\mathcal S_3$ give exactly $80$ full
$\mathrm{GL}(6,2)$ conjugacy classes supported on its underlying spread orbit.
Classifying by the smaller of the two complementary orientation weights,
the numbers are
\[
\begin{array}{c|rrrrrrrrrrr}
r&0&1&2&3&4&5&6&7&8&9&10\\ \hline
\text{classes}&1&1&1&2&3&5&8&10&13&17&19.
\end{array}
\]
Only the weight-zero class is linear.
\end{corollary}
\begin{proof}
Here $|P_3|=60480$. Its projective cycle types and multiplicities are
listed below; a superscript is the number of cycles of that length.
\begin{center}
\small
\begin{tabular}{l|l|r}
matrix description & projective cycle type & multiplicity\\ \hline
scalar & $1^{21}$ & 1\\
$J_2(a)\oplus a$ & $1^5\,2^8$ & 315\\
$J_3(a)$ & $1\,2^2\,4^4$ & 3780\\
$aI_2\oplus b$, $a\ne b$ & $1^6\,3^5$ & 672\\
$J_2(a)\oplus b$, $a\ne b$ & $1^2\,2^2\,3\,6^2$ & 10080\\
three distinct eigenvalues & $1^3\,3^6$ & 2240\\
quadratic block, normalized order 5 & $1\,5^4$ & 8064\\
quadratic block, normalized order 15 & $1\,5\,15$ & 16128\\
cubic block, projective order 3 & $3^7$ & 1920\\
cubic block, projective order 7 & $7^3$ & 5760\\
cubic block, projective order 21 & $21$ & 11520
\end{tabular}
\end{center}
For completeness these multiplicities follow directly from rational
canonical form. In $\mathrm{GL}(3,4)$, the first six types have respectively
$3,3,3,6,6,1$ conjugacy classes. Their corresponding centralizer orders are
\[
181440,\quad576,\quad48,\quad540,\quad36,\quad27.
\]
Divide their total element counts by three
to pass to projective matrices.
For a quadratic block plus eigenvalue $a$, divide the quadratic roots
by $a$. The $12$ elements of $\mathbb F_{16}\setminus\mathbb F_4$
have orders $5$ or $15$, with $4$ or $8$ such roots.
They give $6$ or $12$ matrix classes with centralizer order $45$.
For a cubic block, roots outside $\mathbb F_4$ have projective orders
$3,7,21$, represented by $6,18,36$ field elements.
They give $2,6,12$ matrix classes with centralizer order $63$.
For the Jordan types, $J_2(a)\oplus a$ has projective order two and five
fixed points. The powers of $J_3(a)$ have respectively one, five, and
twenty-one fixed points at exponents $1,2,4$, giving $1\,2^2\,4^4$.
For $J_2(a)\oplus b$, the corresponding counts at exponents $1,2,3,6$
are $2,6,5,21$, giving $1^2\,2^2\,3\,6^2$.
The remaining cycle types follow from the orders on the invariant line,
quadratic component, or cubic field. All multiplicities sum to $60480$,
so the types are exhaustive.

Summing multiplicity times $2^{\text{number of cycles}}$ gives $9676800$;
\eqref{eq:standard-class} gives $9676800/120960=80$.
The number of $r$-subset classes under $P_3$ is the coefficient of $z^r$ in
\[
\frac{1}{60480}\sum_{g\in P_3}
       \prod_{C\text{ a point cycle of }g}(1+z^{|C|}).
\]
Using the table gives the displayed row. Semilinear conjugation pairs
weights $r$ and $21-r$, so each entry also counts the oriented classes
of the corresponding weight pair. Only the two uniform orientations are
linear, and they belong to the same class.
\end{proof}

An independent complete traversal of the $2^{21}$ orientation vectors
supplies all $80$ representatives, their full linear stabilizer orders,
and their orbit sizes. Among these, four classes have stabilizer $C_3$,
accounting for $483840$ orientations on the fixed standard spread.
These representative and stabilizer data are finite computations, distinct
from the general proof of Theorem~\ref{thm:standard-asymptotic}.
The supplied verifier also independently enumerates the $60480$ projective
matrices and recovers the preceding cycle table and weight polynomial.

\subsection*{Coordinate remarks and scope}
For the six-bit example, the fixed point $38$ can serve as the origin after
relabeling, while $0$ remains the additive origin.
In the value set $\{0,\dots,63\}$ the $9$-multiples are exactly the $8$
self-doubled words ($x_i=x_{i+3}$), a $3$-dimensional subspace.
To parametrize the long orbit, define
\[
\theta(j)=T^j(63),\qquad j\in\mathbb Z/63.
\]
This is a bijection onto $V\setminus\{38\}$. Apply the CRT isomorphism to the
\emph{iteration exponent} $j$, sending it to $(j\bmod7,j\bmod9)$.
Thus $\theta(0)=63$ corresponds to $(0,0)$, whereas
$0=\theta(62)$ corresponds to $(6,8)$, since $T(0)=63$.
Reduction of the integer word labels modulo $63$ does not give this parametrization.
The cyclotomic factor $\Phi_{2^k+1}$ reflects
$(2^k-1)(2^k+1)=2^{2k}-1$; the fixed point of $K$ accounts for the remaining point.

The general fixed-map description reduces counting to the line-spread
enumeration problem; no general formula for $s_n$ is supplied here.
The affine orbit count is complete for every dimension, while the scalar
Singer triples provide a distinguished family with controlled permutation
cycle types and two-generator affine groups. Theorems~\ref{thm:affine-group}, \ref{thm:fixed-n4-conjugacy} and
\ref{thm:standard-asymptotic} give the general affine group structure,
the full four-dimensional classification, and the standard-spread
asymptotics. Corollary~\ref{cor:standard-n6} counts all six-dimensional
standard-spread oriented classes; it does not classify orientations on
the other $131043$ underlying spread types (out of $131044$ in total).
The next section treats nonlinear companion groups.


\section{Nonlinear companion groups}\label{sec:nonlinear}
An $F$-fixed map generates only a cyclic group of order three with its
companion. Large companion groups therefore require a three-element
$F$-orbit. If two companions both have cycle type $1+(2^n-1)$, their
cycles have odd length, so their group is contained in $A_{2^n}$.
The following example attains this upper bound in dimension four.

\begin{proposition}[Two long cycles with distinct fixed points]
\label{prop:long-cycle-double-transitive}
Let $\Omega$ be a finite set of size $N\ge4$, and suppose
$f,g\in\mathrm{Sym}(\Omega)$ both have cycle type $1+(N-1)$.
If their unique fixed points are distinct, then
$H=\langle f,g\rangle$ acts doubly transitively on $\Omega$.
If in addition $N=2^n$ with $n\ge3$ and $H$ contains a
$3$-cycle, then $H=A_N$.
\end{proposition}
\begin{proof}
Let $a$ be the unique fixed point of $f$.
Since $g$ fixes a different point, $g(a)\ne a$.
The $(N-1)$-cycle $f$ acts transitively on
$\Omega\setminus\{a\}$, so $H$ is transitive on $\Omega$:
its orbit containing $a$ also contains $g(a)$ and every other point.
Moreover $f\in H_a$, and $f$ itself is transitive on
$\Omega\setminus\{a\}$. Thus the point stabilizer is
transitive on the remaining points and $H$ is doubly transitive.

Now let $N=2^n\ge8$. An $(N-1)$-cycle is even, so
$H\le A_N$. The action is primitive since it is doubly
transitive. The classical Jordan theorem for a primitive
degree-$N$ group containing a prime cycle of length
$p\le N-3$, applied with $p=3$, implies $A_N\le H$
(see, e.g., \cite{DM}). Therefore $H=A_N$.
\end{proof}

This sufficient condition separates the group-theoretic deduction
from the difficult existence problem for complete mappings with
prescribed companion cycle types: the distinct-fixed-point hypothesis
gives double transitivity automatically, while the existence of a
suitable short-cycle element still requires additional arguments.


\begin{proposition}[A nonlinear alternating-group example]\label{prop:nonlinear-A16}
On $V=\mathbb F_2^4$, with integer labels and XOR addition, let $f$ be
the permutation in the first row below and let $g=F(f)$.
\begin{center}
\small
\begin{tabular}{c|rrrrrrrrrrrrrrrr}
$x$&0&1&2&3&4&5&6&7&8&9&10&11&12&13&14&15\\ \hline
$f(x)$&1&5&12&3&11&7&15&10&2&14&9&13&0&8&6&4\\
$g(x)$&3&0&5&10&1&13&11&9&14&6&8&15&12&7&2&4
\end{tabular}
\end{center}
Then $f$ is complete and nonlinear. Each of $f,F(f),F^2(f)$ has cycle
type $1+15$, with fixed points $3,12,1$ respectively, and
\[
\langle f,F(f)\rangle=A_{16}.
\]
\end{proposition}
\begin{proof}
The first row and $x\oplus f(x)$ each permute the labels $0,\ldots,15$;
inverting the latter gives the second row.
The cycles other than the indicated fixed points are
\begin{align*}
f:\;&(0\ 1\ 5\ 7\ 10\ 9\ 14\ 6\ 15\ 4\ 11\ 13\ 8\ 2\ 12),\\
g:\;&(0\ 3\ 10\ 8\ 14\ 2\ 5\ 13\ 7\ 9\ 6\ 11\ 15\ 4\ 1),\\
F^2(f):\;&(0\ 12\ 14\ 7\ 2\ 10\ 13\ 6\ 8\ 5\ 4\ 11\ 15\ 9\ 3).
\end{align*}
The third cycle follows from $F^2(f)=\mathrm{id}\oplus f^{-1}$.
An affine map would satisfy
$f(1\oplus2)=f(1)\oplus f(2)\oplus f(0)$.
Here these values are $3$ and $8$, so $f$ is nonlinear.
Since $F$ preserves the affine class, all three companions are nonlinear.

Use right-to-left composition, write $[u,v]=uvu^{-1}v^{-1}$, and put
$w=f^2g$. Directly from the table,
\[
w=(0\ 3\ 14)(1\ 5\ 2\ 10\ 12)
   (4\ 7\ 6\ 8\ 15\ 13\ 9),
\qquad w(11)=11.
\]
Thus $s=w^{35}=(3\ 0\ 14)$.
All generators are even, so $H=\langle f,g\rangle\le A_{16}$.
The cycles $f^jsf^{-j}=(3\ f^j(0)\ f^j(14))$ belong to $H$,
and their supports contain every point because $f$ cycles all points
other than $3$.

To see that these cycles generate $A_{16}$, make a graph on the other
$15$ points, joining $a,b$ whenever $(3\ a\ b)$ is available.
There are no isolated vertices. If two graph components have edges
$ab$ and $cd$, the identity
$[(3\ a\ b),(3\ c\ d)]=(3\ a\ c)$ joins them.
After the graph is connected, conjugating $(3\ b\ c)$ by $(3\ a\ b)$
gives $(a\ 3\ c)$, so paths supply $(3\ a\ b)$ for every pair.
Finally $(3\ a\ b)(3\ b\ c)=(a\ b\ c)$.
These are all three-cycles, which generate $A_{16}$.
Hence $H\ge A_{16}$, proving equality without a group-order search.
\end{proof}

Alternatively, the distinct fixed points $3$ and $12$ give double
transitivity by Proposition~\ref{prop:long-cycle-double-transitive};
the explicit word $w^{35}$ is a $3$-cycle. The Jordan-theorem
conclusion in that proposition therefore yields $A_{16}$ directly.
The elementary proof above is retained independently of Jordan's theorem.

Complete mapping polynomials over $\mathbb F_{16}$ and their reduced-degree
distribution were computationally studied by Yuan--Tong--Zhang \cite{YTZ07};
Bors--Wang studied cycle types and coset-wise affine constructions of
complete mappings \cite{BW21,BWcos22}. These works provide important
precedents for the finite-field maps and cycle-type methods. The result here
concerns the additional group generated by one specified map and its
prescribed companion, which is not the object counted by those cited results.
A construction in infinitely many dimensions with simultaneous cycle-type
and companion-group control is not supplied here. The groups generated by
the full collections of complete mappings in \cite{BW22} concern a different
set of generators.

\begin{proposition}[Translations of fixed maps]\label{prop:translated-fixed}
Let $f_0$ be an $F$-fixed complete mapping on $\mathbb F_2^n$ and
let $b\ne0$. Write $\tau_b(x)=x\oplus b$ and $f=\tau_b f_0$.
Then $f$ is complete and
\[
F(f)=f_0\tau_b,\qquad
F^2(f)=\tau_b f_0\tau_b,\qquad
\langle f,F(f)\rangle=\langle f_0,\tau_b\rangle.
\]
The orbit of $0$ under this group is exactly the four-point line
\[
U_b=\{0,b,f_0(b),f_0^2(b)\}.
\]
The induced action on $U_b$ is $A_4$.
Thus this construction cannot give a transitive companion group
when $n>2$.
\end{proposition}
\begin{proof}
The fixed-map identities give $f_0^3=\mathrm{id}$ and
$\mathrm{id}\oplus f_0=f_0^2$.
Therefore $\mathrm{id}\oplus f=\tau_b f_0^2$ is a permutation,
and its inverse is $g=f_0\tau_b$.
Also
\[
(\mathrm{id}\oplus g)(x)
=x\oplus f_0(x\oplus b)
=b\oplus f_0^2(x\oplus b),
\]
so taking its inverse gives the formula for $F^2(f)$.
Since $gf=f_0^2$, the group generated by $f,g$ contains $f_0$
and then $\tau_b=ff_0^{-1}$; the reverse inclusion is immediate.

Proposition~\ref{prop:F-spreads} makes $U_b$ a two-dimensional
subspace, invariant under both $f_0$ and $\tau_b$.
Starting at $0$, translation reaches $b$ and $f_0$ reaches the other
two nonzero points, so the orbit is all of $U_b$.
On these points $f_0=(b\ f_0(b)\ f_0^2(b))$ and
$\tau_b=(0\ b)(f_0(b)\ f_0^2(b))$.
These generate the three double transpositions and a three-cycle,
hence $A_4$.
\end{proof}

\begin{thebibliography}{99}
\bibitem{HP55} M.~Hall, L.~J.~Paige, \emph{Complete mappings of finite groups},
Pacific J.~Math.~5 (1955) 541--549,
\href{https://doi.org/10.2140/pjm.1955.5.541}{doi:10.2140/pjm.1955.5.541}.
\bibitem{MW22} B.~D.~McKay, I.~M.~Wanless,
\emph{Enumeration of Latin squares with conjugate symmetry},
J.~Combin.~Des.~30 (2022) 105--130,
\href{https://doi.org/10.1002/jcd.21814}{doi:10.1002/jcd.21814}.
\bibitem{Mag22} N.~Magaud,
\emph{Proof Pearl: Formalizing Spreads and Packings of the Smallest Projective
Space PG(3,2) Using the Coq Proof Assistant}, in ITP 2022,
LIPIcs~237 (2022) 25:1--25:17,
\href{https://doi.org/10.4230/LIPIcs.ITP.2022.25}{doi:10.4230/LIPIcs.ITP.2022.25}.
\bibitem{LN} R.~Lidl, H.~Niederreiter, \emph{Finite Fields}, 2nd ed., Cambridge, 1997.
\bibitem{vW02} P.~van Wamelen, \emph{Jacobi sums over finite fields},
Acta Arith.~102(1) (2002) 1--20,
\href{https://doi.org/10.4064/aa102-1-1}{doi:10.4064/aa102-1-1}.
\bibitem{LP09} T.~K.~Lim, C.~E.~Praeger,
\emph{On generalized Paley graphs and their automorphism groups},
Michigan Math.~J.~58(1) (2009) 293--308,
\href{https://doi.org/10.1307/mmj/1242071694}{doi:10.1307/mmj/1242071694}.
\bibitem{DaMc21} M.~L.~Dawsey, D.~McCarthy,
\emph{Generalized Paley graphs and their complete subgraphs of orders
three and four}, Res.~Math.~Sci.~8 (2021), Article~18,
\href{https://arxiv.org/abs/2006.14716}{arXiv:2006.14716}.
\bibitem{PV24} R.~A.~Podest\'a, D.~E.~Videla,
\emph{Spectral properties of generalized Paley graphs of
$(q^\ell+1)$-th powers and applications},
\href{https://arxiv.org/abs/1812.03332v7}{arXiv:1812.03332v7} (2024).
\bibitem{Katz89} N.~M.~Katz, \emph{An estimate for character sums},
J.~Amer.~Math.~Soc.~2 (1989) 197--200,
\href{https://web.math.princeton.edu/~nmk/old/estcharsums.pdf}{author's copy}.
\bibitem{FFSTV13} R.~R.~Farashahi, P.-A.~Fouque, I.~E.~Shparlinski,
M.~Tibouchi, J.~F.~Voloch, \emph{Indifferentiable deterministic hashing
to elliptic and hyperelliptic curves}, Math.~Comp.~82 (2013) 491--512,
\href{https://doi.org/10.1090/S0025-5718-2012-02606-8}{doi:10.1090/S0025-5718-2012-02606-8}.
\bibitem{CK21} S.~D.~Cohen, G.~Kapetanakis,
\emph{Finite field extensions with the line or translate property for
$r$-primitive elements}, J.~Aust.~Math.~Soc.~111 (2021) 313--319,
\href{https://doi.org/10.1017/S1446788720000099}{doi:10.1017/S1446788720000099}.
\bibitem{CKR22} S.~D.~Cohen, G.~Kapetanakis, L.~Reis,
\emph{The existence of $\mathbb F_q$-primitive points on curves using
freeness}, C.~R.~Math.~360 (2022) 641--652,
\href{https://doi.org/10.5802/crmath.328}{doi:10.5802/crmath.328}.
\bibitem{AN22} J.~J.~R.~Aguirre, V.~G.~L.~Neumann,
\emph{Pairs of $r$-primitive and $k$-normal elements in finite fields},
arXiv:2210.11504v1 (2022),
\href{https://arxiv.org/abs/2210.11504}{arXiv:2210.11504}.
\bibitem{Ha21} D.~Hachenberger,
\emph{Ovoids and primitive normal bases for quartic extensions of Galois fields},
J.~Algebr.~Comb.~53 (2021) 85--114,
doi:10.1007/s10801-019-00920-8.
\bibitem{SSL24} J.~Sharma, R.~Sarma, S.~Laishram,
\emph{Arithmetic progressions of $r$-primitive elements in a field},
Bull.~Braz.~Math.~Soc.~(N.S.)~55 (2024), Article~40,
\href{https://doi.org/10.1007/s00574-024-00412-9}{doi:10.1007/s00574-024-00412-9}.
\bibitem{CSST25} K.~Chatterjee, H.~Sharma, A.~Shukla, S.~K.~Tiwari,
\emph{Arithmetic progression in a finite field with prescribed norms},
Forum Math.~37(1) (2025) 345--357,
\href{https://doi.org/10.1515/forum-2024-0026}{doi:10.1515/forum-2024-0026}.
\bibitem{CSh25} A.~Choudhary, J.~Sharma,
\emph{Arithmetic progressions of some special elements in finite fields},
Commun.~Algebra (published online 2025),
\href{https://doi.org/10.1080/00927872.2025.2561951}{doi:10.1080/00927872.2025.2561951}.
\bibitem{KN26} H.~Kudo, Y.~Nogata, \emph{Cyclotomic numbers of order
$q-1$ over $\mathbb F_{q^r}$}, arXiv:2604.25461 (2026),
\url{https://arxiv.org/abs/2604.25461}.
\bibitem{DM} J.~D.~Dixon, B.~Mortimer, \emph{Permutation Groups}, Springer, 1996.
\bibitem{Coh} S.~D.~Cohen, \emph{The orders of related elements of a finite field}, Ramanujan J.~7 (2003) 169--183, doi:10.1023/A:1026295128600.
\bibitem{Mor} O.~Moreno, \emph{On the existence of a primitive quadratic of trace $1$ over $GF(p^m)$}, J.~Combin.~Theory Ser.~A 51 (1989) 104--110.
\bibitem{Coh90} S.~D.~Cohen, \emph{Primitive elements and polynomials with arbitrary trace}, Discrete Math.~83 (1990) 1--7.
\bibitem{CP05} S.~D.~Cohen, M.~Pre\v{s}ern,
\emph{Primitive finite field elements with prescribed trace},
Southeast Asian Bull.~Math.~29(2) (2005) 283--300,
\href{https://researchportal.hw.ac.uk/en/publications/primitive-finite-field-elements-with-prescribed-trace/}{publication record}.
\bibitem{GPT} C.~Gravel, D.~Panario, D.~Thomson, \emph{Unicyclic strong permutations}, arXiv:1809.03551 (2018); Cryptogr.~Commun.~11 (2019) 1211--1231.
\bibitem{BW21} A.~Bors, Q.~Wang, \emph{Cycle types of complete mappings of finite fields},
J.~Algebra~591 (2022) 577--610,
\href{https://doi.org/10.1016/j.jalgebra.2021.09.017}{doi:10.1016/j.jalgebra.2021.09.017}.
\bibitem{BW22} A.~Bors, Q.~Wang, \emph{Compositions and parities of complete mappings and of orthomorphisms},
J.~Combin.~Theory Ser.~A~196 (2023) 105723,
\href{https://doi.org/10.1016/j.jcta.2022.105723}{doi:10.1016/j.jcta.2022.105723}.
\bibitem{FS15} J.~Fulman, D.~Stanton, \emph{On the distribution of the number of fixed vectors for the finite classical groups},
Ann.~Comb.~20 (2016) 755--773,
\href{https://doi.org/10.1007/s00026-016-0336-7}{doi:10.1007/s00026-016-0336-7}.
\bibitem{MT09} Z.~T.~Mateva, S.~T.~Topalova, \emph{Line spreads of PG(5,2)},
J.~Combin.~Des.~17(1) (2009) 90--102,
\href{https://doi.org/10.1002/jcd.20198}{doi:10.1002/jcd.20198}.
\bibitem{Muy26} A.~M\"uyesser,
\emph{Cycle type in Hall--Paige: a proof of the Friedlander--Gordon--Tannenbaum conjecture},
Forum Math.~Sigma~14 (2026) e50,
\href{https://doi.org/10.1017/fms.2026.10197}{doi:10.1017/fms.2026.10197}.

\bibitem{RS98} A.~Russell, R.~Sundaram,
\emph{A note on the asymptotics and computational complexity of graph
distinguishability}, Electron.~J.~Combin.~5(1) (1998), Research Paper~23,
\href{https://doi.org/10.37236/1361}{doi:10.37236/1361}.
\bibitem{CNS84} P.~J.~Cameron, P.~M.~Neumann, J.~Saxl,
\emph{On groups with no regular orbits on the set of subsets},
Arch.~Math.~43 (1984) 295--296,
\href{https://doi.org/10.1007/BF01196649}{doi:10.1007/BF01196649}.
\bibitem{Se97} \'{A}.~Seress,
\emph{Primitive groups with no regular orbits on the set of subsets},
Bull.~London Math.~Soc.~29 (1997) 697--704,
\href{https://doi.org/10.1112/S0024609397003536}{doi:10.1112/S0024609397003536}.
\bibitem{BWcos22} A.~Bors, Q.~Wang,
\emph{Coset-wise affine functions and cycle types of complete mappings},
Finite Fields Appl.~83 (2022) 102088,
\href{https://doi.org/10.1016/j.ffa.2022.102088}{doi:10.1016/j.ffa.2022.102088}.
\bibitem{YTZ07} Y.~Yuan, Y.~Tong, H.~Zhang,
\emph{Complete mapping polynomials over finite field $\mathbb F_{16}$},
in \emph{Arithmetic of Finite Fields}, WAIFI 2007, Lecture Notes in Computer
Science 4547, Springer (2007), 147--158,
\href{https://doi.org/10.1007/978-3-540-73074-3_12}{doi:10.1007/978-3-540-73074-3\_12}.
\end{thebibliography}

\appendix
\section{Computational appendix}\label{app:code}
All numerical claims (fixed points, cycle types, group order $4032$, and the
characteristic polynomials of $L_T$, $L_D$, and $L_K$) were verified
computationally. The script \path{verify_paperA_structure.py} reproduces
the coordinate, cycle, group-order, basis-ordering, and $\tau$-collapse data,
as well as the $\mathbb F_{64}$ instance of Proposition~\ref{prop:F-orbits}
and the CRT exponent convention. It computes the characteristic polynomials
as $\det(tI+L)$ by the permutation expansion in characteristic two.
The separate \path{verify_k1_29.py} checks the finite PTE certificate.
The script \path{verify_complete_mapping_orbits.py} enumerates all binary
matrices and translations for $n\le4$, checks Lemmas~\ref{lem:fp}--\ref{lem:F3}
and the coordinate-permutation identities of Proposition~\ref{prop:F-parastrophy},
and compares the resulting counts with Proposition~\ref{prop:F-affine}.
For $n=4$ it enumerates line spreads by exact cover, orients every line,
checks $F(f)=f$, and tests additivity, reproducing the counts following
Proposition~\ref{prop:F-spreads}.
The additional \path{verify_affine_groups_and_conjugacy.py} partitions all
admissible matrices for $n\le4$ into linear-conjugacy classes and directly
closes the two-generator permutation group for every translation at each
matrix representative. The $4,16,80$ closure cases for $n=2,3,4$, respectively,
cover all $8,384,93184$ affine pairs by linear conjugacy. These closures are
compared with Theorem~\ref{thm:affine-group}; transitivity, block systems,
and double transitivity are also tested directly. The same script checks
Theorem~\ref{thm:fixed-n4-conjugacy} using all $1792$ maps and all $20160$
linear changes of coordinates, and verifies the signed Burnside count.
It shares basic permutation and spread routines with
\path{verify_complete_mapping_orbits.py}; the direct closure computation
does not use the proposed group formula. Its data file
\path{affine_groups_and_conjugacy_results.json} includes every matrix
representative and translation case, map representatives, stabilizer matrices,
and the fixed-orientation counts.
These deterministic scripts require only the Python standard library and were
run with Python~3.12.14. The ancillary data include the $64$ rows of $T,D,K$,
all $720$ basis orderings, the $29$ PTE certificates with their factorizations,
the $56$ four-dimensional spreads, and a nonlinear fixed-map witness.
The script \path{export_paperA_data.py} regenerates the coordinate and
certificate tables; it reuses the two original verifiers rather than constituting
an independent implementation. The scripts, data, proof sources, reproduction
instructions, and the current SHA-256 manifest accompany this manuscript as
ancillary files. The single command \path{python reproduce_paperA.py} checks
the manifest, runs all finite verifiers, regenerates the data, and checks their
hashes again.
The project repository is
\begin{center}
\small\url{https://github.com/63331849121/yijing-mathematical-structures}.
\end{center}
Relative to the standard basis $e_1,\dots,e_6$, the linear part of the main map is
\[
L_T=\begin{pmatrix}
0&0&1&0&0&0\\
1&1&0&1&1&0\\
0&0&1&0&1&0\\
0&1&1&1&1&0\\
1&0&1&1&0&1\\
0&0&1&1&0&0
\end{pmatrix}.
\]
\emph{Formalization.} \path{CompleteMappingTriality.lean} checks
the function-level rotation, $F^2(f)=\mathrm{id}\oplus f^{-1}$,
$F^3(f)=f$, and the unique-fixed-point assertion of
Lemmas~\ref{lem:fp}--\ref{lem:F3} for additive groups of exponent two.

The affine group-generation argument (for a field $\mathbb K$ with a primitive
element $\alpha$ and $c\ne0$, the maps $T$ and $D$ generate $\mathrm{AGL}(1,\mathbb K)$, the core
of Theorem~\ref{thm:agld} and \S\ref{sec:2k}) and the algebraic lemmas used in it (the commutator
is a translation; conjugation of translations by $T$) have been formalized and
machine-checked in \path{AGL.lean}.

The conditional file
\path{CohenEvenConditional.lean}
formalizes the field-level Frobenius identities and the resulting
order-theoretic $\tau$-collapse under an explicit order interface; it does not
formalize the paper proof of the existence theorem. The concrete cycle-type data of
\S\ref{sec:T}--\ref{sec:tdk} were verified by independent computation. The existence
part of Theorem~\ref{thm:2k} is proved in Appendix~\ref{app:pte}; its analytic part
is not formalized in Lean, and its finite range is discharged by the deterministic
certificate there. Thus Cohen's theorem is not used as a proof dependency; the
classical references are cited only for historical attribution. The general matrix-group structure, its action criteria, the signed Burnside
formula, four-dimensional conjugacy classification, coordinate-permutation
description, line-spread classification, and affine orbit-count proof are
paper proofs, not Lean proofs. The inverse-realization and arithmetic-sieve criteria, the
primitive-group nonrealizability theorem, the positive Jacobi-sum/sieve
bound, the fixed-index realization theorem and index-three threshold,
the quantitative density and growing-index results, norm descent and
the quarter-power obstruction, primitive inverse classification and
the generating-triangle correspondence,
the local invariant-orientation
criterion, and the long-cycle double-transitivity criterion are also
paper proofs, not part of the Lean formalization.

All three Lean files were checked using Lean~4.31.0 and mathlib commit
\hashpart{fabf563a}%
\hashpart{7c95a166}%
\hashpart{b8d7b6ef}%
\hashpart{ca11c8b4}%
\hashpart{dc9d911f}. The ancillary \path{lean/} directory contains these files,
\path{lean-toolchain}, \path{lakefile.lean}, and
\path{lake-manifest.json}, pinning the proof environment.
These file checks used an existing dependency cache; a fresh dependency download
and cloud CI run are not claimed. Code is supplied under the MIT license, and
the manuscript and data under CC~BY~4.0.


The verifier \path{verify_standard_spread_orbits.py} independently enumerates
the $60480$ projective matrices for the standard six-dimensional spread
and traverses all $2^{21}$ orientations using elementary matrix generators
and the signed Frobenius action. It checks the projective cycle table,
the weight polynomial, and the $80$ disjoint orientation orbits.
The data \path{standard_spread_results.json} contain every representative,
its map, full linear stabilizer order, and both its orientation-orbit
size and its full $\mathrm{GL}(6,2)$ conjugacy-class size.
Only this underlying spread type is covered.
The verifier \path{verify_nonlinear_companions.py} checks the explicit
sixteen-point tables, all three cycle types, the three-cycle word
certificate, and the translated-fixed identities.
It does not enumerate $A_{16}$: the group equality has the elementary
proof in Proposition~\ref{prop:nonlinear-A16}.
The arbitrary-dimensional asymptotic and translated-fixed results
are paper proofs, outside the supplied Lean formalization.

The verifier \path{verify_scalar_realization.py} compares two independent
finite counting methods: binary-polynomial gcds and explicit multiplicative
subgroups in finite-field coordinates. It covers every divisor of $2^n-1$
for $n=2,4,6,8,10,12,14$, checks the Jacobi bounds and slope-generation
sieve, and verifies the degree-eight witness in
Corollary~\ref{cor:index-three-realization}. The data
\path{scalar_realization_results.json} include the index-three counts
$|\mathcal R_d|=0,0,24,90,396,1722$ for
$n=4,6,8,10,12,14$, respectively. These computations verify finite
arithmetic certificates; the infinite-family and threshold assertions
use the paper proofs in Section~\ref{sec:group}.

The verifier \path{verify_scalar_depth.py} audits the quantitative
main term and error by exact rational arithmetic in even dimensions
$2$ through $14$. It checks norm descent by polynomial gcds for
$2\le r\le8$, all proper divisors $s\mid r$, and all $u\mid2^s-1$.
The quarter-power counts are checked for $s=1,2,3,4,5$ by polynomial
gcds and the divisor sieve, with independent field enumeration and
subfield-location checks for $s=1,2,3,4$.
For primitive subgroup orders in even dimensions $2$ through $14$,
it checks Frobenius and six-ratio orbits, the degree of the polynomial
encoding the realizing slopes, and the conjugacy-class counts.
It also directly enumerates generating triangles and their affine
stabilizers on $4$ and $16$ vertices. The corresponding data are
\path{scalar_depth_results.json}. These are finite arithmetic audits;
the uniform exponent obstruction and classification use the paper proofs.


\subsection*{Finite certificates for the inner critical theorem}
\label{app:inner-cert}
Theorem~\ref{thm:inner-all} uses the analytic argument for $s\ge33$
and the following complete finite certificate for $2\le s\le32$.
A hexadecimal integer encodes a binary polynomial with its least
significant bit the constant coefficient. Each row means
\[
 E=\mathbb F_2[X]/(M_s),\qquad
 L=E[Z]/(Z^2+(1+b+w)Z+b),\qquad a=Z.
\]
Thus only $M_s,b,w$ need be specified; all field operations are exact.
\begin{center}\small
\begin{tabular}{r|lll}
$s$ & $M_s$ (hex) & $b$ (hex) & $w$ (hex)\\\hline
2 & \texttt{13} & \texttt{8} & \texttt{1}\\
3 & \texttt{43} & \texttt{6} & \texttt{3b}\\
4 & \texttt{11b} & \texttt{2f} & \texttt{1}\\
5 & \texttt{409} & \texttt{9b} & \texttt{2c9}\\
6 & \texttt{1009} & \texttt{c78} & \texttt{8e5}\\
7 & \texttt{4021} & \texttt{3e8a} & \texttt{1}\\
8 & \texttt{1002b} & \texttt{2878} & \texttt{689}\\
9 & \texttt{40009} & \texttt{303be} & \texttt{12d39}\\
10 & \texttt{100009} & \texttt{a3085} & \texttt{1}\\
11 & \texttt{400003} & \texttt{662cc} & \texttt{3205f9}\\
12 & \texttt{100001b} & \texttt{d0842d} & \texttt{1}\\
13 & \texttt{400001b} & \texttt{32f9317} & \texttt{1}\\
14 & \texttt{10000003} & \texttt{15a3baa} & \texttt{6c6a1f9}\\
15 & \texttt{40000003} & \texttt{317d6a85} & \texttt{36559252}\\
16 & \texttt{10000008d} & \texttt{4893e13d} & \texttt{2e37b719}\\
17 & \texttt{40000001b} & \texttt{dcbd299} & \texttt{91607c0b}\\
\end{tabular}
\end{center}

\begin{center}\small
\begin{tabular}{r|lll}
$s$ & $M_s$ (hex) & $b$ (hex) & $w$ (hex)\\\hline
18 & \texttt{1000000035} & \texttt{78c0283c4} & \texttt{1}\\
19 & \texttt{4000000063} & \texttt{1153a3d1ef} & \texttt{16004f10d3}\\
20 & \texttt{10000000039} & \texttt{1f9215dca7} & \texttt{1}\\
21 & \texttt{40000000027} & \texttt{2657c0cd035} & \texttt{374caba3695}\\
22 & \texttt{100000000021} & \texttt{fe622270c1c} & \texttt{1}\\
23 & \texttt{400000000003} & \texttt{17c45ffbd76} & \texttt{29bd79b1ae4}\\
24 & \texttt{100000000002d} & \texttt{fed891dbe6c9} & \texttt{ec7a6a1920af}\\
25 & \texttt{400000000001d} & \texttt{2907b6c341c09} & \texttt{1}\\
26 & \texttt{10000000000009} & \texttt{c6d27ccd2b0df} & \texttt{432bf4829f702}\\
27 & \texttt{4000000000007d} & \texttt{28047cef28c34e} & \texttt{1}\\
28 & \texttt{100000000000095} & \texttt{a586528c5b1c9} & \texttt{1}\\
29 & \texttt{400000000000063} & \texttt{243f16ec4e21119} & \texttt{1}\\
30 & \texttt{1000000000000003} & \texttt{8c388b10c413f59} & \texttt{1}\\
31 & \texttt{4000000000000069} & \texttt{12dfcd22c3b1728b} & \texttt{1}\\
32 & \texttt{1000000000000001b} & \texttt{ae3111b3bb81428f} & \texttt{c171e75cc47c2b2e}\\
\end{tabular}
\end{center}

The ancillary file \path{inner_critical_certificates.json} also records
the full factorizations of $A,B$, every nonidentity residue
$a^{d/p}$ for a prime $p\mid d$, and recursive Lucas primality
certificates for those primes. Verification consists of the following
finite identities, carried out by \path{verify_inner_critical.py}.
The additional \path{verify_inner_symmetry.py} checks the
four-torsion identities and rational Lang-fiber multiplier at
all 31 certificate points, and enumerates every primitive-norm
fiber for $s=2,\ldots,6$. It checks four-point orbits, order
preservation, divisibility by four and the exact integer
multiplicity distributions underlying the character symmetries.
The new verification logs and runtime metadata are supplied with
the ancillary files; these finite checks supplement the paper
proof of Proposition~\ref{prop:inner-four-torsion}.
First, $M_s$ has degree $n=2s$, and
\[
 X^{2^n}\equiv X\pmod {M_s},\qquad
 \gcd(M_s,X^{2^{n/r}}-X)=1\quad(r\mid n\text{ prime}),
\]
which certify irreducibility. Second, $b^A=w^A=1$,
$b^{A/p}\ne1$ for every prime $p\mid A$, $c=1+b+w\ne0$,
and $\operatorname{Tr}_{E/\mathbb F_2}(b/c^2)=1$.
Third, in the quadratic field one checks
\[
 a^d=(1+a)^d=1,\quad a^B=b,\quad(1+a)^B=w,\quad
 a^{d/p}\ne1\quad(p\mid d\text{ prime}).
\]
Primality is certified without probable-prime assumptions: the base
case is $2$; for an integer $p>2$ one supplies a complete recursively
certified factorization of $p-1$ and an integer $v$ with
$v^{p-1}\equiv1\pmod p$ and
$\gcd(v^{(p-1)/r}-1,p)=1$ for every prime $r\mid p-1$.
For any prime divisor $\ell$ of $p$, these conditions force the order
of $v$ modulo $\ell$ to be $p-1$, hence $\ell\ge p$ and $p$ is prime.
The independent verifier uses different field multiplication and
powering implementations from the certificate finder; it requires
only the Python standard library, run without the \texttt{-O} flag.

The small-prime list used in the analytic constant is
\[
\begin{gathered}
5,13,17,29,37,41,53,61,73,89,97,101,\\
109,113,137,149,157,173,181,193,197,229,233,241.
\end{gathered}
\]
Their product is
$6045408114863853991373345088496297320256672385$;
the integer comparison is $256^{24}<32^8$ times this product.
The verifier checks this comparison as well.

As a separate small-field audit, multiplicative-subgroup enumeration
and polynomial gcds with the M\"obius sieve agree on the following
counts. Here $N_d^{\mathrm{prim}}$ has the meaning in
Theorem~\ref{thm:inner-primitive-count}.
\[
\begin{array}{c|r|r|r|r|r}
s&d&T_d&|\mathcal R_d|&N_d^{\mathrm{prim}}&D_d\\\hline
2&85&24&24&16&8\\
3&585&86&72&48&24\\
4&4369&240&240&224&208\\
5&33825&1214&960&520&240\\
6&266305&4104&3960&2928&1992
\end{array}
\]
The data also verify the genus-one Hasse bounds on every primitive
norm fiber in this small range. They confirm that the diagonal
$b=w$ gives no realizing slope for $s=3,5$.
These supplementary counts do not replace the general proof.
The present all-dimensional theorem uses the classical Weil bound
and a finite exact certificate, and is outside the Lean formalization.

\paragraph{Additional double-order certificates.}
The file \path{inner_double_order_certificates.json} uses the same
$M_s,b$ as the preceding table and replaces $w$ by $b^{e_s}$, with
the following exponents. The root $a=Z$ of
$Z^2+(1+b+b^{e_s})Z+b$ and its translate $1+a$ both have exact
order $d$, for every $2\le s\le32$.
\begin{center}\small
\begin{tabular}{rr|rr}
$s$&$e_s$&$s$&$e_s$\\\hline
2&1&18&1\\
3&4&19&5\\
4&1&20&1\\
5&4&21&5\\
6&3&22&1\\
7&1&23&8\\
8&4&24&4\\
9&26&25&1\\
10&1&26&3\\
11&7&27&1\\
12&1&28&1\\
13&1&29&1\\
14&8&30&1\\
15&4&31&1\\
16&4&32&3\\
17&7&&\\
\end{tabular}
\end{center}
The independent \path{verify_inner_mixed.py} checks irreducibility,
the trace, both norms and both nonidentity residues
$a^{d/p},(1+a)^{d/p}$ for every prime $p\mid d$. It uses the
complete factorizations and recursive Lucas certificates of
\path{inner_critical_certificates.json}; all new residues and all
norm-fiber counts for $s=2,\ldots,6$ are recorded in
\path{inner_mixed_verification.json}. It also checks the corrected
mixed Lang functions and trace multipliers at the 31 certificate points.
These computations establish precisely the stated finite range.
Together with Proposition~\ref{prop:inner-diagonal-double} for $s\ge33$,
they complete the all-dimensional double full-order theorem.
The new verifier \path{verify_inner_double_all.py} independently
rechecks all 31 double-order certificates, including their recursive
Lucas primality proofs, without calling a certificate finder.

\subsection*{Integer bounds for double full-order existence}
The bound in Proposition~\ref{prop:inner-diagonal-double} is positive
when $\sqrt t\ge4W(d)$. We cover $33\le s\le113$ by exact divisor
upper bounds, and $s\ge114$ by a uniform inequality for odd integers.
No large residual factor is assumed to be prime.

For the uniform bound let $P$ be the product of all 53 odd primes
less than $256$, namely
\[
\begin{gathered}
3,5,7,11,13,17,19,23,29,31,37,41,43,47,53,59,61,67,71,73,79,83,\\
89,97,101,103,107,109,113,127,131,137,139,149,151,157,163,167,173,\\
179,181,191,193,197,199,211,223,227,229,233,239,241,251.
\end{gathered}
\]
The exact integer comparison $256^{53}<4096^8P$ gives, for every
odd $n$,
\[
 \frac{W(n)^8}{n}
 =\prod_{p^e\parallel n}\frac{256}{p^e}
 \le\prod_{3\le p<256}\frac{256}{p}
 =\frac{256^{53}}{P}<4096^8.
\]
Here the product on the right is over primes: each factor for a
prime at least $256$ is at most one, and adding missing smaller
prime factors only increases the bound. Thus $W(n)<4096n^{1/8}$.
For $s\ge114$, $d<2t^3$ and $t>2^{113}$ prove the sufficient
condition, as computed in Proposition~\ref{prop:inner-diagonal-double}.

For each $33\le s\le113$, divide $A$ and $B$ by every prime at
most $4096$, with its full multiplicity. Let $j_A,j_B$ be the
numbers of distinct small factors and $R_A,R_B$ the residuals.
All prime factors of each residual exceed $4096$. If
$4097^{k_A}\le R_A<4097^{k_A+1}$ and similarly for $B$, then
\[
 \omega(d)=\omega(A)+\omega(B)
    \le r_s:=j_A+k_A+j_B+k_B.
\]
The following table lists all 81 bounds; each satisfies
$s\ge4+2r_s$, so $\sqrt t\ge4\,2^{r_s}\ge4W(d)$.
\begin{center}\small
\begin{tabular}{rr|rr|rr}
$s$&$r_s$&$s$&$r_s$&$s$&$r_s$\\\hline
33&11&60&15&87&24\\
34&9&61&16&88&21\\
35&13&62&15&89&23\\
36&11&63&21&90&27\\
37&11&64&15&91&26\\
38&11&65&20&92&23\\
39&13&66&19&93&25\\
40&9&67&18&94&24\\
41&11&68&16&95&26\\
42&14&69&19&96&24\\
43&12&70&20&97&26\\
44&11&71&18&98&26\\
45&17&72&19&99&29\\
46&13&73&20&100&26\\
47&13&74&19&101&26\\
48&12&75&23&102&27\\
49&15&76&19&103&26\\
50&14&77&23&104&25\\
51&15&78&22&105&32\\
52&12&79&21&106&27\\
53&14&80&19&107&28\\
54&16&81&23&108&29\\
55&16&82&21&109&27\\
56&13&83&22&110&29\\
57&16&84&23&111&30\\
58&15&85&24&112&27\\
59&15&86&22&113&29\\
\end{tabular}
\end{center}
The full small factorizations, residual integers and their integer
power bounds are supplied in
\path{inner_double_omega_certificates.json}. The independent verifier
\path{verify_inner_double_all.py} reconstructs every small prime by
deterministic trial division, checks the exact products and all 81
inequalities, and checks $256^{53}<4096^8P$.
Its results, log and runtime information are
\path{inner_double_all_verification.json},
\path{inner_double_all_run.log} and
\path{double_all_verification_environment.json}.
It also exhaustively checks the trace-one subfamily for $2\le s\le6$:
its counts are respectively $8,0,16,0,24$. In particular the
all-dimensional claim is for $D_d$, not for $\Delta_t$.
The disjoint ranges $2\le s\le32$, $33\le s\le113$, and
$s\ge114$ cover every integer $s\ge2$.

\subsection*{Bases for the coordinate presentation}
For the coordinate presentation $T(x)=(2^{2k}-1)\oplus\mathrm{coords}_G(x)$ on
$2k$ bits the following bases give the stated group order $\langle T,D\rangle$ and
$K$-cycle type:
\begin{center}
\begin{tabular}{c|c|c|c|c}
$k$ & basis $G$ & $\langle T,D\rangle$ & $K$ type & factor \\
\hline
1 & $(1,3)$ & $12=4\cdot3$ & $1+1\cdot3$ & $\Phi_3$ \\
2 & $(1,2,5,13)$ & $240=16\cdot15$ & $1+3\cdot5$ & $\Phi_5$ \\
3 & $(63,33,18,49,1,21)$ & $4032=64\cdot63$ & $1+7\cdot9$ & $\Phi_9$ \\
4 & $(134,136,241,71,34,65,3,143)$ & $65280=256\cdot255$ & $1+15\cdot17$ & a degree-$8$ factor of $\Phi_{17}$ \\
\end{tabular}
\end{center}
All four rows were verified by independent computation. The $720$-case count
was reproduced by enumerating all orderings of the six displayed vectors, solving
the corresponding $\mathbb{F}_2$ linear system for $\mathrm{coords}_G$, and
recording the cycle types of $T,D,K$ whenever the two inverses exist.

\begin{samepage}
\subsection*{Explicit $\tau$-collapse data}
With $\mathbb{F}_{2^{2k}}=\mathbb{F}_2[t]/(f)$ and elements written as polynomials in $t$,
the following rows realize the order conditions of Theorem~\ref{thm:2k}:
\begin{center}
\begin{tabular}{c|c|c|c|c}
$k$ & $f$ & $\gamma$ & $\tau=\gamma\oplus\gamma^{-1}$ & $\alpha=(1\oplus\gamma)^{-1}$ \\
\hline
1 & $t^2+t+1$ & $t$ & $1$ & $t$ \\
2 & $t^4+t+1$ & $t^3$ & $t^2+t+1$ & $t$ \\
3 & $t^6+t+1$ & $t^2+t$ & $t^4+t^3+1$ & $t^5+t^3+t^2$ \\
4 & $t^8+t^4+t^3+t^2+1$ & $t^3+t^2+t+1$ & $t^7+t^4+t^3+1$ & $t^6+t^4+t^3+t^2+1$ \\
\end{tabular}
\end{center}
\end{samepage}
In each row $\mathrm{ord}(\gamma)=\mathrm{ord}(1\oplus\alpha^{-1})=2^k+1$, $\tau$ is
primitive in $\mathbb{F}_{2^k}$, and
$\mathrm{ord}(\alpha)=\mathrm{ord}(1\oplus\alpha)=2^{2k}-1$ (all verified). The defining polynomial $f$ is a field-representation choice and need not coincide with the characteristic polynomial $p_T$ of \eqref{eq:pT}; for instance, the $k=3$ row uses $f=t^6+t+1$, while $p_T(t)=t^6+t^5+t^3+t^2+1$.



\section{Self-contained proof for the even-characteristic case}\label{app:pte}
\begin{proof}[Proof of Theorem~\ref{thm:pte}]
The statement is classical (see \cite{Mor} and
\cite[Theorem~1, p.~1]{Coh90}; it also follows from the equivalent square-order formulation in
\cite[Theorem~1.5, p.~170]{Coh}); we give a self-contained proof in the
even-characteristic case used above. Cohen and Presern \cite{CP05} also give a
self-contained treatment of prescribed traces. This appendix reconstructs the
classical trace-one argument and makes its finite verification explicit; no
priority is claimed for the existence result or the character-sum method.
Put $q=2^k$, $L=\mathbb{F}_{q^2}$, $K_0=\mathbb{F}_q$, and
$M=q^2-1=(q-1)(q+1)$. Let
\[
E_1=\{x\in L:x^q+x=1\}.
\]
The trace map $L\to K_0$ is surjective and $K_0$-linear, so $|E_1|=q$. If
$x\in E_1$, then $x\notin K_0$ and its minimal polynomial over $K_0$ is
$X^2+X+xx^q$. Hence it is enough to prove that $E_1$ contains an element of
order $M$.

Let $\psi$ be the nontrivial additive character of $\mathbb F_2$, and put
$\psi_L(x)=\psi(\operatorname{Tr}_{L/\mathbb F_2}(x))$ and
$\psi_{K_0}(a)=\psi(\operatorname{Tr}_{K_0/\mathbb F_2}(a))$. For a multiplicative
character $\chi$ of $L^*$, extend $\chi$ to $0$ by $\chi(0)=0$ and define
\[
g_L(\chi)=\sum_{x\in L^*}\chi(x)\psi_L(x),\qquad
g_{K_0}(\chi)=\sum_{a\in K_0^*}\chi(a)\psi_{K_0}(a),\qquad
S(\chi)=\sum_{x\in E_1}\chi(x).
\]
If $\chi$ is nontrivial on $L^*$, then $|g_L(\chi)|=q$; if
$\chi|_{K_0^*}$ is nontrivial, then $|g_{K_0}(\chi)|=\sqrt q$. Indeed, the usual
squared-modulus calculation gives
\[
|g_L(\chi)|^2
=\sum_{t\in L^*}\chi(t)^{-1}\sum_{a\in L^*}\psi_L(a(1-t))=q^2,
\]
and the same argument in $K_0$ gives $q$.

We first compute $S(\chi)$. For $\alpha\in K_0$ set
$E_\alpha=\{x\in L:x^q+x=\alpha\}$; these sets partition $L$, and
$E_\alpha=\alpha E_1$ for $\alpha\ne0$. If $\chi^{q+1}\ne1$, then the
restriction of $\chi$ to $K_0^*$ is nontrivial. Since $E_0=K_0$,
its contribution is $\sum_{x\in E_0}\chi(x)=0$; hence
\[
g_L(\chi)
=\sum_{\alpha\in K_0}\psi_{K_0}(\alpha)\sum_{x\in E_\alpha}\chi(x)
=S(\chi)g_{K_0}(\chi),
\]
so $|S(\chi)|=q/\sqrt q=\sqrt q$.

Suppose now that $\chi^{q+1}=1$ and $\chi\ne1$. Then $\chi|_{K_0^*}=1$.
Let $H=\{h\in L^*:h^{q+1}=1\}$; since $\gcd(q-1,q+1)=1$, the group $L^*$
is the direct product $K_0^*\times H$. Write $x=ch$ with $c\in K_0^*$ and
$h\in H$. The condition $x^q+x=1$ becomes
$c(h+h^{-1})=1$; thus $h\ne1$ and $c=(h+h^{-1})^{-1}$. This gives a
bijection between $H\setminus\{1\}$ and $E_1$, and therefore
\[
S(\chi)=\sum_{h\in H,\ h\ne1}\chi(h)
=\sum_{h\in H}\chi(h)-1=-1,
\]
because $\chi|_H$ is nontrivial.

We shall use the primitive-element indicator
\begin{equation}\label{eq:pte-indicator}
\mathbf 1_{\operatorname{ord}(x)=M}
=\frac{\varphi(M)}M
\sum_{\delta\mid M}\frac{\mu(\delta)}{\varphi(\delta)}
\sum_{\operatorname{ord}(\chi)=\delta}\chi(x).
\end{equation}
To see this, fix a generator $g$ of $L^*$ and write $x=g^a$. The left side
is the indicator of $\gcd(a,M)=1$. For a character with
$\chi(g)=e^{2\pi i j/M}$ and order $\delta=M/\gcd(M,j)$, its unnormalized
Fourier coefficient is the Ramanujan sum
$c_M(j)=\mu(\delta)\varphi(M)/\varphi(\delta)$; Fourier inversion gives
\eqref{eq:pte-indicator}. (The displayed identity for $c_M(j)$ follows by
M\"obius inversion.)

Assume for contradiction that $E_1$ has no element of order $M$. Summing
\eqref{eq:pte-indicator} over $E_1$ gives
\[
0=\frac{\varphi(M)}M
\sum_{\delta\mid M}\frac{\mu(\delta)}{\varphi(\delta)}S_\delta,
\qquad
S_\delta=\sum_{\operatorname{ord}(\chi)=\delta}S(\chi).
\]
Let $s=\omega(M)$ and $s'=\omega(q+1)$ be the numbers of distinct prime
divisors. The $\delta=1$ term is $q$. If $\delta\mid q+1$ and $\delta>1$,
then each character of order $\delta$ satisfies $\chi^{q+1}=1$, so the
corresponding contribution is $-\mu(\delta)$; summing over
$\delta\mid q+1$, $\delta>1$, gives $1$. For the remaining squarefree
divisors $\delta\nmid q+1$, the character sums are bounded by $\sqrt q$,
and their number is at most $2^s-2^{s'}$. Consequently,
\[
q+1\le(2^s-2^{s'})\sqrt q. \tag{B.1}\label{eq:pte-bound}
\]
Since $q+1>\sqrt q$, this implies
\[
2^s>2^{k/2}, \qquad\text{so}\qquad s>k/2.
\]
In particular $k<2s$. Let $p_1<p_2<\cdots$ be the primes and
$P_s=p_1\cdots p_s$. Since $M=2^{2k}-1$ has exactly $s$ distinct prime
divisors, $M\ge P_s$. We have
\[
P_{16}=32589158477190044730>2^{64}.
\]
For $n\ge16$, the ratio
\[
\frac{P_{n+1}/2^{4(n+1)}}{P_n/2^{4n}}
=\frac{p_{n+1}}{16}
\]
is greater than $1$, because $p_{n+1}\ge59$. Hence $P_s>2^{4s}$ for
every $s\ge16$. But $k<2s$ gives
\[
M=2^{2k}-1<2^{2k}<2^{4s},
\]
so $s\ge16$ is impossible. Thus $s\le15$ and $k<2s\le30$, i.e.
$k\le29$.

It remains to cover this finite range. Each row below specifies an irreducible
degree-$k$ polynomial $m_k$ over $\mathbb F_2$ and an element $a_k\in K_0$
with $\operatorname{Tr}_{K_0/\mathbb F_2}(a_k)=1$, where
$K_0=\mathbb F_2[x]/(m_k)$. The displayed hexadecimal integers encode
their binary coefficient vectors.
For each row let $t$ be a root of $X^2+X+a_k$ in
$K_0[t]/(X^2+X+a_k)$. The deterministic verifier checks that $m_k$ is
irreducible and that $\operatorname{Tr}_{K_0/\mathbb F_2}(a_k)=1$.
The trace condition implies that $X^2+X+a_k$ has no root in $K_0$
and is irreducible there; the quotient is therefore a quadratic
field extension. The verifier then checks that
$t^{2^k}=t+1$, and that $t^M=1$ while $t^{M/p}\ne1$ for every prime
$p\mid M$. The last two checks show that $\operatorname{ord}(t)=M$; the
third shows that $\operatorname{Tr}_{L/K_0}(t)=1$. A verifier for the table, \path{verify_k1_29.py},
is included as an ancillary file. These checks are performed row by row
using polynomial arithmetic in characteristic two. Its SHA-256 hash is
\hashpart{042CE7F7}%
\hashpart{4EDFF13A}%
\hashpart{DE5C6285}%
\hashpart{E42CB3D1}%
\hashpart{5491FC25}%
\hashpart{C65B99BC}%
\hashpart{670D2560}%
\hashpart{47EE8D9F}.
\begin{center}
\small
\begin{tabular}{rrr}
$k$ & $m_k$ & $a_k$\\
\hline
1 & \texttt{0x3} & \texttt{0x1}\\
2 & \texttt{0x7} & \texttt{0x3}\\
3 & \texttt{0xb} & \texttt{0x5}\\
4 & \texttt{0x13} & \texttt{0xd}\\
5 & \texttt{0x25} & \texttt{0x13}\\
6 & \texttt{0x43} & \texttt{0x34}\\
7 & \texttt{0x83} & \texttt{0x47}\\
8 & \texttt{0x11b} & \texttt{0xc8}\\
9 & \texttt{0x203} & \texttt{0x2b}\\
10 & \texttt{0x409} & \texttt{0xcb}\\
11 & \texttt{0x805} & \texttt{0x525}\\
12 & \texttt{0x1009} & \texttt{0x39e}\\
13 & \texttt{0x201b} & \texttt{0x1ec4}\\
14 & \texttt{0x4021} & \texttt{0x3717}\\
15 & \texttt{0x8003} & \texttt{0x62bf}\\
16 & \texttt{0x1002b} & \texttt{0x1e6e}\\
17 & \texttt{0x20009} & \texttt{0x1155b}\\
18 & \texttt{0x40009} & \texttt{0x2ee84}\\
19 & \texttt{0x80027} & \texttt{0xb9af}\\
20 & \texttt{0x100009} & \texttt{0x6aee2}\\
21 & \texttt{0x200005} & \texttt{0x15750b}\\
22 & \texttt{0x400003} & \texttt{0x24a150}\\
23 & \texttt{0x800021} & \texttt{0x6c92f9}\\
24 & \texttt{0x100001b} & \texttt{0x389ec4}\\
25 & \texttt{0x2000009} & \texttt{0x6801ad}\\
26 & \texttt{0x400001b} & \texttt{0x2707d2e}\\
27 & \texttt{0x8000027} & \texttt{0x6e9ab1e}\\
28 & \texttt{0x10000003} & \texttt{0xfac7cd2}\\
29 & \texttt{0x20000005} & \texttt{0xc69b55e}\\
\end{tabular}
\end{center}
The table therefore supplies the required element for every $k\le29$.
Combining this with the contradiction for $k\ge30$ proves the theorem.
\end{proof}

\end{document}
